\hypertarget{nace_kartlar_16-18}{%
\chapter{Faz 3, parça 3 --- Ağaç, kâğıt ve basım: Gümrük Birliği'nden çıkışın sınıf düzeyinde analizi (NACE 16--18)}\label{nace_kartlar_16-18}}

\emph{Durum: Faz 3 / 9, parça 3 / 16; 22 Kasım 2026. Önceki parçalar: \belgebag{nace_kartlar_10-12}{}{\kod{nace\_\allowbreak{}kartlar\_\allowbreak{}10-\allowbreak{}12.\allowbreak{}md}}, \belgebag{nace_kartlar_13-15}{}{\kod{nace\_\allowbreak{}kartlar\_\allowbreak{}13-\allowbreak{}15.\allowbreak{}md}}. Önceki fazlar: \belgebag{01_teshis}{}{\kod{01\_\allowbreak{}teshis.\allowbreak{}md}}, \belgebag{02a_politika_uzayi}{}{\kod{02a\_\allowbreak{}politika\_\allowbreak{}uzayi.\allowbreak{}md}}, \belgebag{02b_secenekler_hukuk}{}{\kod{02b\_\allowbreak{}secenekler\_\allowbreak{}hukuk.\allowbreak{}md}}, \belgebag{02c_hazirlik}{}{\kod{02c\_\allowbreak{}hazirlik.\allowbreak{}md}}. Kümülatif defter: \belgebag{defter}{}{\kod{defter.\allowbreak{}md}}. Bütün 230 imalat sınıfının ana matrisi: \href{../veri/nace_matris.csv}{\kod{veri/\allowbreak{}nace\_\allowbreak{}matris.\allowbreak{}csv}}.}

\hypertarget{nace_kartlar_16-18-yuxf6netici-uxf6zeti}{%
\section{Yönetici özeti}\label{nace_kartlar_16-18-yuxf6netici-uxf6zeti}}

Ağaç ürünleri, kâğıt ve basım, Gümrük Birliği'nden çıkışın en az etkilediği sanayi kollarından biri. Üç bölümün 18 sınıfında 245 bin kişi çalışıyor; bu, imalat istihdamının \%4,9 kadarı. Yıllık ihracat 5,8 milyar dolar ve bunun yalnızca \%24 kadarı AB'ye gidiyor. 2025'te AB'ye satılan 1.295 milyon avroluk ürünün tamamı Gümrük Birliği'nin sanayi rejiminde, ama anlaşmasız çıkışta AB'nin bu ürünlere uygulayacağı tarife ortalama yalnızca \%1,3. İmalatın genelinde bu oran \%5,6. AB kâğıt ve kâğıt ürünlerinde MFN tarifesini sıfırlamış durumda \etiketV{}. Ahşap levhada ve kontrplakta tarife \%7--10, ama bu ürünlerde Türkiye'nin ana pazarları AB dışında.

Bu yüzden çıkışın bu parçadaki istihdam ve cari denge etkisi sıfıra yakın, hatta küçük ölçüde olumlu. Anlaşmasız çıkışta (S4) üç bölümün AB'ye ihracatı T+5'te yılda 134 milyon dolar azalıyor. Buna karşılık istihdam etkisi medyanda +471 kişi ve cari dengeye katkı +101 milyon dolar. Bunun nedeni ithalat tarafı. Türkiye kâğıt ve kartonda net ithalatçı ve bu ithalatın \%63 kadarı AB'den geliyor. S4'te Türkiye'nin c2 rejimi, yurt içi üretimi güçlü kâğıt kalitelerinde AB menşeli ürüne ek gümrük vergisi (İGV) uyguluyor. AB'den kâğıt ve karton ithalatı 239 milyon dolar azalıyor ve bunun bir kısmını yerli üretim karşılıyor. Önerilen bitiş durumu olan derin ve kapsamlı serbest ticaret anlaşmasında (S2) etki daha da küçük: istihdam +163 kişi, cari denge +41 milyon dolar. Dağılımın alt ucunda da (P10) kayıplar birkaç yüz kişiyi geçmiyor.

Bu sonuç, sektörün AB'ye bağlı risklerinin küçük olduğu anlamına gelmiyor. Asıl riskler Gümrük Birliği'nden bağımsız ve modelin tarife ve sınır maliyeti kanallarının dışında. Bunlardan ilki AB'nin ormansızlaşma tüzüğü (EUDR). Tüzük, beş hafta sonra, 30 Aralık 2026'da büyük ve orta ölçekli işletmeler için uygulanmaya başlıyor \etiketV{}. Kapsamında kereste, ahşap levha, kâğıt hamuru, kâğıt ve karton var; basılı ürünler 2025 sonundaki değişiklikle kapsamdan çıkarıldı \etiketV{}. Türkiye, Komisyon'un ülke sınıflandırmasında düşük riskli; bu, AB'li ithalatçı için basitleştirilmiş özen yükümlülüğü ve \%1'lik denetim oranı demek \etiketV{}. Ama ürünün hangi parselden kesilmiş odundan yapıldığını gösteren coğrafi konum verisi yine şart. Türkiye'de ormanların neredeyse tamamı devletin ve satışlar Orman Genel Müdürlüğü (OGM) üzerinden yapılıyor. Bu, bu veriyi tek elden ve ucuza üretme olanağı veriyor. Bu parçanın ilk önerisi, OGM'nin her satışta EUDR'ye uygun coğrafi konum verisini alıcıya vermesi (E-16-01).

İkinci risk ticaret savunması ve Rus odunu. AB, Rus huş kontrplağına 2021'de koyduğu anti-damping vergisini 2024'te Türkiye'den gönderilen kontrplağa da genişletti. Gerekçe, Rus kontrplağının Türkiye üzerinden AB'ye sokulmasıydı \etiketV{}. AB'de üreticilerin başvurusuyla bu vergiler için 6 Kasım 2026'da süre sonu gözden geçirmesi açıldı ve vergiler soruşturma süresince (12--15 ay) yürürlükte kalıyor \etiketV{}. Türkiye'nin kontrplak ithalatının \%45, kaplama ithalatının \%44 kadarı Rusya'dan geliyor. Bu tablo, hasmane bir AB'nin Türk levha ve kontrplak ihracatına karşı başka genişletmeler ve yaptırımları dolanma soruşturmaları açması için zemin hazırlıyor. Gümrük Birliği bu araçlara karşı bugün de koruma sağlamıyor: genişletme, Gümrük Birliği içinde yapıldı. Önerilen savunma, ihracat hatlarında Rus ve Belarus odunundan ayrışmış ve belgelenmiş bir tedarik (E-16-02) ile S2'de ticaret savunması öncesi zorunlu istişare (E-16-04).

Üçüncü bulgu Türkiye'nin kendi ticaret politikasıyla ilgili. Türkiye, oluklu mukavva kâğıdında (testliner ve fluting) 2024'ten beri üç yıllık bir korunma önlemi uyguluyor. Önlem, ithalatın büyük kısmının geldiği AB menşeli kâğıdı da kapsıyor; muafiyet yalnızca listelenen gelişmekte olan ülkelere tanınmış \etiketV{}. Yani Gümrük Birliği, Türkiye'nin AB'den gelen bir sanayi ürününe DTÖ korunma önlemi uygulamasını engellemedi. Aralık 2025'te kaplamalı kâğıtta yeni bir korunma soruşturması açıldı \etiketV{}. Bu önlemlerin ve S4'teki c2 tarifesinin bedeli aşağı akışta ödeniyor. Kâğıdı girdi olarak kullanan oluklu ambalaj üreticileri (41 bin çalışan) ve matbaalar (48 bin çalışan), kâğıt üreticilerinden (14 bin çalışan) çok daha fazla kişi çalıştırıyor. c2'nin HS6 düzeyindeki kural tabanlı uygulaması, sıvı ambalaj kartonu ve kraftliner gibi Türkiye'de üretilmeyen kalitelere de İGV koyuyor. Önerilen düzeltme, bu kalitelerde İGV'nin bütün menşeler için sıfırlanması (E-17.12-01).

Duruşlar bu tabloyu yansıtıyor. 18 sınıfın 5 tanesi için ``Geliştir'', 13 tanesi için ``Koru'' öneriliyor. Geliştirilecek sınıflar ahşap levha (Türkiye dünyanın 3. büyük lif levha ihracatçısı), hijyen kâğıtları ve bezler, kâğıt ve karton ile iki küçük sınıf. ``Koru'' duruşu bu parçada çoğunlukla ``çıkışa özgü önlem gerekmiyor, izle'' anlamına geliyor. Paketin yıllık mali maliyeti S2 yolunda 0,39 milyar TL, S4 yolunda 1,09 milyar TL. Parça, kümülatif enflasyon bütçesinin S2 yolunda 0,2, S4 yolunda önlemlerle 0,7 baz puanını kullanıyor.

İki olumsuz bulgu açıkça yazılmalı. Birincisi, modelin S4'te kâğıtta bulduğu istihdam artışı abartılı. Artışın bir kısmı, c2'nin Türkiye'de üretilmeyen kalitelerde de AB ithalatını pahalandırmasından geliyor ve bu kalitelerde yerli üretim ithalatın yerini alamaz. E-17.12-01 uygulanırsa bu kazanç da küçülür. İkincisi, bu parçanın asıl riskleri olan EUDR uyumu, ticaret savunması ve Rus odunu modelin dışında. Bunların büyüklüğü firma verisi olmadan ölçülemiyor (VB-40, VB-41). Senaryo sonuçlarının küçüklüğü, sektörün AB pazarındaki konumunun güvende olduğu biçiminde okunmamalı.

\Needspace{7\baselineskip}

\textbf{Tablo Ö.1 --- Çıkışın 16--18'e etkisi (T+5'te yıllık; bugüne göre fark)}

\nopagebreak[4]

\begingroup\gbtablo\footnotesize

\begin{longtable}[]{@{}lrlrl@{}}
\toprule\noalign{}
\textbf{Gösterge} & \textbf{S2: P50} & \textbf{S2: P10 / P90} & \textbf{S4: P50} & \textbf{S4: P10 / P90} \\
\midrule\noalign{}
\endhead
\bottomrule\noalign{}
\endlastfoot
AB'ye ihracat (\$ mn) & −75 & −168 / +28 & −134 & −219 / −58 \\
İstihdam (kişi) & +163 & −375 / +815 & +471 & −8 / +964 \\
Cari denge (\$ mn) & +41 & −67 / +178 & +101 & +3 / +197 \\
TÜFE, birinci tur (bp) & 0,2 & 0,2 / 0,3 & 0,9 & 0,8 / 1,1 \\
TÜFE, önerilen önlemlerle (bp) & 0,2 & --- & 0,7 & --- \\
Önlemlerin mali maliyeti (mlr TL/yıl) & 0,39 & --- & 1,09 & --- \\
\end{longtable}

\endgroup

\emph{Kaynak:} \etiketT{} f3e Monte Carlo (N = 2.000; yüzdelikler bölüm toplamının dağılımından). S2: derin ve kapsamlı STA; S4: anlaşmasız çıkış ve tek taraflı paket. Önlemler ve mali maliyet Bölüm 4'teki eylemlerden.

\hypertarget{nace_kartlar_16-18-1-bu-paruxe7anux131n-sorusu-ve-okuma-kux131lavuzu}{%
\section{1. Bu parçanın sorusu ve okuma kılavuzu}\label{nace_kartlar_16-18-1-bu-paruxe7anux131n-sorusu-ve-okuma-kux131lavuzu}}

Faz 3, Gümrük Birliği'nden çıkışın imalat sanayisine etkisini NACE Rev. 2'nin dört basamaklı sınıfları düzeyinde ölçüyor ve her sınıf için bir strateji öneriyor. Bu üçüncü parça ağaç ve ağaç ürünlerini (16), kâğıt ve kâğıt ürünlerini (17) ve basımı (18) kapsıyor. Sorular önceki parçalardakiyle aynı. Çıkış bu sınıfın ihracatını, üretimini, istihdamını, cari dengeye katkısını ve tüketici fiyatlarını hangi kanaldan ve ne ölçüde etkiler? Hangi yetkinlikler korunmalı, hangileri geliştirilmeli? AB ile ve üçüncü ülkelerle neyin müzakeresi yapılmalı? Hangi önlemi kim, hangi hukuki araçla ve ne zaman almalı?

Senaryolar değişmedi. \textbf{S2}, AB ile derin ve kapsamlı bir serbest ticaret anlaşması; sanayi ürünlerinde tarifeler menşe şartıyla sıfır ve çalışmanın önerdiği bitiş durumu bu. \textbf{S3}, AB'nin üye devletlerin onayı olmadan imzalayabileceği temel bir ticaret anlaşması; S2'ye giden yolda ara durum. \textbf{S4}, anlaşmasız çıkış ve Türkiye'nin tek taraflı paketi: AB, Türk mallarına MFN tarifesini uyguluyor; Türkiye de AB mallarına MFN tarifesini ve seçici ek gümrük vergisini (c2 rejimi) uyguluyor, girdilerde tarifeleri sıfırlıyor ve tüketici fiyatlarını gerektiğinde vergi takasıyla dengeliyor. AB'nin tutumu hasmane ve üyelik perspektifi yok. İki takvim var: T2028'de çıkış tarihi T = 1 Ocak 2028, T2029'da T = 1 Ocak 2029. Yöntem, veri kaynakları ve kanıt etiketleri ilk parçanın birinci bölümünde ve Ek B'sinde anlatıldı. Bu parçanın Ek B'si yalnızca yenilikleri veriyor.

Bu parça, önceki ikisinden farklı bir soru soruyor. Gıda ve giyimde çıkışın asıl maliyeti AB'nin tarife duvarıydı. Ağaç, kâğıt ve basımda bu duvar ya yok ya da alçak. Kâğıtta AB'nin MFN tarifesi sıfır, ahşap levhada \%7--10. Bu yüzden parçanın sorusu ``çıkış ne kadar ihracat kaybettirir?'' değil. Soru, ``Gümrük Birliği'nin dışında kalan ve hasmane bir AB'nin kullanabileceği araçlar bu sektörlere ne yapar, ve Türkiye'nin kendi ticaret politikası girdi tarafında kimi korur, kime maliyet yükler?'' Bu araçlar üç başlıkta toplanıyor: ormansızlaşma ve ambalaj düzenlemeleri, ticaret savunması ve Rusya yaptırımları, ve Türkiye'nin kâğıtta uyguladığı korunma önlemleri ile S4'teki c2 tarifesi.

Bu parçada modelde üç değişiklik yapıldı. Birincisi iç ikame tavanı. Model, AB'den ithalat azaldığında boşluğun bir kısmını yerli üretimin dolduracağını varsayıyor. Kâğıt hamurunda Türkiye'nin üretimi yok denecek kadar az olduğu için bu varsayım anlamsız bir istihdam artışı üretiyordu. Artık sınıf tablosunda bir tavan tanımlanabiliyor ve kâğıt hamurunda bu tavan sıfır. İkincisi CE ve onaylanmış kuruluş bağımlılığı. AB'ye satılan yapı ürünlerinin bir kısmı, AB'nin tanıdığı bir kuruluşun deneyini gerektiriyor; S4'te Türk kuruluşlarının raporları tanınmıyor. Bu maliyet şimdiye kadar bütün sınıflarda sıfırdı. Bu parçada ahşap levha, ahşap döşeme ve doğrama için tanımlandı. Üçüncüsü bir düzeltme. Model S2'de, AB'nin MFN tarifesi zaten sıfır olan ürünlerde de ihracatçıya menşe ispatı maliyeti yazıyordu. Oysa tarifesi sıfır olan üründe ihracatçı tercih istemez ve menşe ispatına da katlanmaz. Bu hata, kâğıt gibi sıfır tarifeli ürünlerin S2 kaybını abartıyordu; ilaç ve elektronik gibi sonraki parçalarda çok daha büyük bir sapma yaratacaktı. Üç değişiklik Ek B'de anlatılıyor; önceki parçaların sonuçlarına etkileri Ek D'de.

Belge önceki parçalarla aynı sırada ilerliyor. İkinci bölüm sektörün Gümrük Birliği içindeki yerini, üçüncü bölüm çıkışın ve çıkıştan bağımsız AB araçlarının kanallarını anlatıyor. Dördüncü bölüm duruşları, eylem paketini, müzakere pozisyonunu, erken uyarı göstergelerini ve bölgesel dağılımı bir araya getiriyor. Beşinci bölüm sınıfları tek tek inceliyor; etkisi sınırlı dokuz sınıf tek bir tabloda toplanıyor. Altıncı bölüm kümülatif defteri ve kalite kontrollerini veriyor. Ekler ana matris satırlarını (Ek A), yöntemdeki yenilikleri (Ek B), sınıf veri sayfalarını (Ek C), önceki parçalara ve fazlara düzeltmeleri (Ek D) ve varsayım ile veri boşluğu kayıtlarını (Ek E) içeriyor. On beş alanlı sınıf kartı yine üç yerde tamamlanıyor: nitel alanlar Bölüm 5'te, sayısal alanlar Ek C'de, eylemler ve tetikler Bölüm 4'te.

\hypertarget{nace_kartlar_16-18-2-sektuxf6ruxfcn-guxfcmruxfck-birliux11fi-iuxe7indeki-yeri}{%
\section{2. Sektörün Gümrük Birliği içindeki yeri}\label{nace_kartlar_16-18-2-sektuxf6ruxfcn-guxfcmruxfck-birliux11fi-iuxe7indeki-yeri}}

\Needspace{7\baselineskip}

\textbf{Tablo 2.1 --- Bölümlerin büyüklüğü ve dış ticareti (2024)}

\nopagebreak[4]

\begingroup\gbtablo\footnotesize

\begin{longtable}[]{@{}
  >{\raggedright\arraybackslash}p{(\columnwidth - 16\tabcolsep) * \real{0.2067}}
  >{\raggedleft\arraybackslash}p{(\columnwidth - 16\tabcolsep) * \real{0.0775}}
  >{\raggedleft\arraybackslash}p{(\columnwidth - 16\tabcolsep) * \real{0.0981}}
  >{\raggedleft\arraybackslash}p{(\columnwidth - 16\tabcolsep) * \real{0.1124}}
  >{\raggedleft\arraybackslash}p{(\columnwidth - 16\tabcolsep) * \real{0.1087}}
  >{\raggedleft\arraybackslash}p{(\columnwidth - 16\tabcolsep) * \real{0.1188}}
  >{\raggedleft\arraybackslash}p{(\columnwidth - 16\tabcolsep) * \real{0.0787}}
  >{\raggedleft\arraybackslash}p{(\columnwidth - 16\tabcolsep) * \real{0.1046}}
  >{\raggedleft\arraybackslash}p{(\columnwidth - 16\tabcolsep) * \real{0.0945}}@{}}
\toprule\noalign{}
\begin{minipage}[b]{\linewidth}\raggedright
\textbf{Bölüm}
\end{minipage} & \begin{minipage}[b]{\linewidth}\raggedleft
\textbf{Sınıf}
\end{minipage} & \begin{minipage}[b]{\linewidth}\raggedleft
\textbf{Çalışan (bin)}
\end{minipage} & \begin{minipage}[b]{\linewidth}\raggedleft
\textbf{Katma değer (\$ mlr)}
\end{minipage} & \begin{minipage}[b]{\linewidth}\raggedleft
\textbf{İhracat (\$ mlr)}
\end{minipage} & \begin{minipage}[b]{\linewidth}\raggedleft
\textbf{AB'ye ihracat (\$ mlr)}
\end{minipage} & \begin{minipage}[b]{\linewidth}\raggedleft
\textbf{AB payı}
\end{minipage} & \begin{minipage}[b]{\linewidth}\raggedleft
\textbf{İthalat (\$ mlr)}
\end{minipage} & \begin{minipage}[b]{\linewidth}\raggedleft
\textbf{Denge (\$ mlr)}
\end{minipage} \\
\midrule\noalign{}
\endhead
\bottomrule\noalign{}
\endlastfoot
16 Ağaç ve ağaç ürünleri & 6 & 91,9 & 1,93 & 1,88 & 0,35 & \%19 & 0,83 & +1,04 \\
17 Kâğıt ve kâğıt ürünleri & 7 & 98,8 & 4,18 & 3,89 & 1,02 & \%26 & 4,11 & −0,22 \\
18 Basım ve kayıt ortamlarının çoğaltılması & 5 & 54,7 & 1,18 & 0,02 & 0,01 & \%63 & 0,02 & 0,00 \\
\textbf{16--18 toplamı} & \textbf{18} & \textbf{245,5} & \textbf{7,30} & \textbf{5,78} & \textbf{1,38} & \textbf{\%24} & \textbf{4,96} & \textbf{+0,82} \\
İmalat (C) içindeki payı & 18 / 230 & \%4,9 & \%4,4 & \%2,3 & \%1,3 & --- & \%1,8 & --- \\
\end{longtable}

\endgroup

\emph{Kaynak:} TÜİK Yıllık Sanayi ve Hizmet İstatistikleri 2024 (çalışan, katma değer); BACI HS17 2024 (dış ticaret; Kıbrıs akımları AB dışı, ilk parçanın Ek D-2). \etiketV{}; HS6 → NACE eşlemesi \etiketT{}.

Üç bölüm istihdamda birbirine yakın, ama yapıları çok farklı. Ağaç ürünlerinde (16) 92 bin kişi çalışıyor ve bunların çoğu küçük işletmelerde. Doğramacılık ve ahşap yapı sınıfında (16.23) 11.871 girişim var ve büyük işletmelerin istihdam payı yalnızca \%1. Ahşap levha (16.21) ise tersine yoğunlaşmış bir sınıf: 22 bin çalışanın \%69 kadarı 250'den fazla kişi çalıştıran 15 işletmede. Kâğıtta (17) 99 bin kişi çalışıyor. Ağırlık kâğıdı işleyen sınıflarda: oluklu ambalajda 41 bin, hijyen kâğıtları ve bezlerde 25 bin kişi. Kâğıdın kendisini üreten sınıfta (17.12) çalışan sayısı 13,6 bin. Basımda (18) 55 bin kişi çalışıyor ve bunun neredeyse tamamı matbaacılıkta (18.12, 48 bin).

Verimlilik farkı sınıflar arasında büyük. AB27'deki aynı sınıf = 1 alındığında kişi başı katma değer ahşap levhada 0,55, kâğıt ve kartonda 0,52, hijyen kâğıtlarında 0,46, matbaacılıkta 0,42, doğramacılıkta ise yalnızca 0,17. Sermaye yoğun ve büyük ölçekli üretim yapan levha ve kâğıt sınıfları Avrupa'daki rakiplerine yaklaşıyor. Atölye düzeyinde kalan doğramacılık ve kereste ise çok geride.

Türkiye birkaç üründe dünya ihracatının ön sıralarında. 2024'te lif levhada (MDF ve HDF dahil, KN 4411) dünya ihracatının \%8,4 kadarını yaparak Almanya ve Çin'den sonra 3. sırada. Hijyen kâğıdı ana bobininde (4803) \%11,4 payla 2., bebek ve hasta bezi gibi hijyen ürünlerinde (9619) \%4,7 payla 7., oluklu mukavva kutuda (4819 10) \%3,2 payla 7. sırada. Bu üstünlüklerin hiçbiri ucuz hammaddeye dayanmıyor. Türkiye odunun bir kısmını, kâğıt hamurunun ise tamamına yakınını ithal ediyor. Üstünlük, büyük ve modern tesislerden, Orta Doğu, Kafkasya ve Kuzey Afrika pazarlarına yakınlıktan ve hızlı teslimattan geliyor \etiketH{}.

Pazar yapısı da bu yakınlığı gösteriyor. Ahşap levha ağır ve hacimli bir ürün; uzağa taşınması pahalı. Sınıfın AB dışı ihracatının \%14 kadarı Irak'a, \%12 kadarı Fas'a gidiyor. AB'deki ilk pazarlar komşular: Romanya (\%28), Bulgaristan (\%21) ve Yunanistan (\%13). Kâğıt ve kartonda da AB içinde ilk üç pazar Yunanistan, Bulgaristan ve Romanya. Oluklu ambalajda tablo farklı: AB'de ilk pazarlar Almanya (\%19) ve Fransa (\%18), AB dışında Birleşik Krallık (\%21) ve ABD (\%17). Hijyen ürünlerinde ilk pazar Irak (\%27). Sonuçta üç bölümün AB'ye ihracat payı \%24, imalat ortalamasının yarısından az.

İthalat tarafında tablo tersine dönüyor. Bölümlerin 5,0 milyar dolarlık ithalatının \%52 kadarı AB'den. Türkiye'nin kâğıt hamuru ithalatı 1.068 milyon dolar. İlk kaynaklar Brezilya (\%25), ABD (\%17), Finlandiya (\%16) ve Uruguay (\%11). Kâğıt ve karton ithalatının (2.526 milyon dolar) \%63 kadarı AB'den; Almanya, İsveç ve Finlandiya'nın yanında Çin de büyük bir kaynak. Bu ithalatın ağırlığı Türkiye'de üretilmeyen ya da az üretilen kalitelerde: kaplamalı ağaçsız kâğıt, kraftliner, sıvı ambalaj kartonu ve kaplamalı karton. Ağaçta ise ithalatın en büyük kaynağı Rusya. Ahşap levha sınıfının ithalatının \%39, kereste sınıfınınkinin \%13 kadarı Rusya'dan. Bölüm 3.2'de görüleceği gibi bu bağımlılık, AB ile ilişkinin en hassas noktası.

\Needspace{7\baselineskip}

\textbf{Tablo 2.2 --- AB'ye ihracatın rejimi ve S4'te eklenecek AB tarifesi (2025)}

\nopagebreak[4]

\begingroup\gbtablo\footnotesize

\begin{longtable}[]{@{}
  >{\raggedright\arraybackslash}p{(\columnwidth - 12\tabcolsep) * \real{0.1937}}
  >{\raggedleft\arraybackslash}p{(\columnwidth - 12\tabcolsep) * \real{0.1348}}
  >{\raggedleft\arraybackslash}p{(\columnwidth - 12\tabcolsep) * \real{0.1234}}
  >{\raggedleft\arraybackslash}p{(\columnwidth - 12\tabcolsep) * \real{0.1371}}
  >{\raggedleft\arraybackslash}p{(\columnwidth - 12\tabcolsep) * \real{0.1337}}
  >{\raggedleft\arraybackslash}p{(\columnwidth - 12\tabcolsep) * \real{0.1672}}
  >{\raggedleft\arraybackslash}p{(\columnwidth - 12\tabcolsep) * \real{0.1101}}@{}}
\toprule\noalign{}
\begin{minipage}[b]{\linewidth}\raggedright
\textbf{Bölüm}
\end{minipage} & \begin{minipage}[b]{\linewidth}\raggedleft
\textbf{AB'ye ihracat (€ mn)}
\end{minipage} & \begin{minipage}[b]{\linewidth}\raggedleft
\textbf{GB (sanayi) payı}
\end{minipage} & \begin{minipage}[b]{\linewidth}\raggedleft
\textbf{S4'te eklenecek tarife}
\end{minipage} & \begin{minipage}[b]{\linewidth}\raggedleft
\textbf{AB MFN ortalaması}
\end{minipage} & \begin{minipage}[b]{\linewidth}\raggedleft
\textbf{Tarifesi ≥ \%10 olan satırların payı}
\end{minipage} & \begin{minipage}[b]{\linewidth}\raggedleft
\textbf{KN8 satırı}
\end{minipage} \\
\midrule\noalign{}
\endhead
\bottomrule\noalign{}
\endlastfoot
16 Ağaç ve ağaç ürünleri & 318 & \%100 & \%4,6 & \%4,6 & \%2 & 173 \\
17 Kâğıt ve kâğıt ürünleri & 967 & \%100 & \%0,2 & \%0,2 & \%0 & 165 \\
18 Basım ve kayıt ortamlarının çoğaltılması & 9 & \%100 & \%0,0 & \%0,0 & \%0 & 1 \\
\textbf{16--18 toplamı} & \textbf{1.295} & \textbf{\%100} & \textbf{\%1,3} & \textbf{\%1,3} & \textbf{\%0} & \textbf{339} \\
İmalat (C) & 98.030 & \%91 & \%5,6 & \%6,1 & \%29 & 6.930 \\
\end{longtable}

\endgroup

\emph{Kaynak:} Eurostat Comext 2025 (TR→AB27, KN8), AB Kombine Nomenklatürü 2026; \kod{kod/\allowbreak{}f3d\_\allowbreak{}ab\_\allowbreak{}tarife.\allowbreak{}py}. Tarifeler TR→AB ihracatı ağırlıklı \etiketT{}. Bu parçanın sınıflarında tarım ve işlenmiş tarım ürünü rejimi yok; S4'te AB'nin MFN tarifesi bütün ihracata ekleniyor.

Tablo 2.2, AB'ye satılan 1.295 milyon avroluk ürünün tamamının Gümrük Birliği'nin sanayi rejiminde olduğunu gösteriyor. Yine de anlaşmasız çıkışta eklenecek tarife düşük: ortalama \%1,3. Kâğıt bölümünde AB'nin tarifesi neredeyse her satırda sıfır; ağaçta ortalama \%4,6. Tarifenin anlamlı olduğu tek sınıf ahşap levha. Yonga levhada ve lif levhada AB'nin MFN tarifesi \%7, kontrplakta \%10'a çıkıyor ve sınıfın ortalaması \%6,9. Basımda ise AB'ye ihracat yok denecek kadar az. Türkiye'nin basılı ürün ihracatı (Fasıl 49, 2024'te 209 milyon dolar, \%36 AB'ye) yayıncılık sınıflarına sayılıyor. Matbaacılığın kendisi bir hizmet sınıfı ve dış ticareti ölçülemiyor (VB-42).

\hypertarget{nace_kartlar_16-18-3-uxe7ux131kux131ux15f-bu-sektuxf6rleri-hangi-kanallardan-etkiliyor}{%
\section{3. Çıkış bu sektörleri hangi kanallardan etkiliyor?}\label{nace_kartlar_16-18-3-uxe7ux131kux131ux15f-bu-sektuxf6rleri-hangi-kanallardan-etkiliyor}}

Önceki iki parçada çıkışın ana kanalı AB'nin tarife duvarıydı. Bu parçada duvar alçak, bu yüzden kanalların sırası değişiyor. Bölüm 3.1 tarifenin neden küçük bir rol oynadığını gösteriyor. Bölüm 3.2 ve 3.3, Gümrük Birliği'nin dışında kalan ve hasmane bir AB'nin kullanabileceği iki araç ailesini ele alıyor: ticaret savunması ile Rusya yaptırımları, ve ormansızlaşma ile ambalaj düzenlemeleri. Bölüm 3.4 ithalat tarafına, yani c2 tarifesinin ve Türkiye'nin kendi korunma önlemlerinin kâğıt zincirindeki dağılım etkisine bakıyor. Bölüm 3.5 S2'nin menşe ve sınır maliyetlerini özetliyor.

\hypertarget{nace_kartlar_16-18-31-aluxe7ak-bir-tarife-duvarux131}{%
\subsection{3.1 Alçak bir tarife duvarı}\label{nace_kartlar_16-18-31-aluxe7ak-bir-tarife-duvarux131}}

AB'nin bu parçadaki ürünlere uyguladığı MFN tarifesi, imalatın geri kalanına göre çok düşük. Kâğıt hamuru, kâğıt, karton ve kâğıt ürünlerinin neredeyse bütün satırlarında oran sıfır. Ağaç ürünlerinde oranlar \%0 ile \%10 arasında ve ağırlıklı ortalama \%4,6. Tarifenin anlamlı olduğu tek büyük sınıf ahşap levha. Burada yonga ve lif levhada \%7, kontrplakta \%10'luk tarife, AB'ye ihracatı 216 milyon dolar olan sınıf için S4'te gerçek bir maliyet. Aşağıdaki tablo, S4'te tarifeyle karşılaşan sınıfları ve S2'deki karşılıklarını gösteriyor.

\Needspace{7\baselineskip}

\textbf{Tablo 3.1 --- GB ve İTÜ rejimindeki sınıflar: S4'te AB tarifesi ve S2'de menşe (T+5, P50)}

\nopagebreak[4]

\begingroup\gbtablo\footnotesize

\begin{longtable}[]{@{}
  >{\raggedright\arraybackslash}p{(\columnwidth - 16\tabcolsep) * \real{0.1355}}
  >{\raggedleft\arraybackslash}p{(\columnwidth - 16\tabcolsep) * \real{0.1242}}
  >{\raggedleft\arraybackslash}p{(\columnwidth - 16\tabcolsep) * \real{0.0962}}
  >{\raggedleft\arraybackslash}p{(\columnwidth - 16\tabcolsep) * \real{0.1077}}
  >{\raggedright\arraybackslash}p{(\columnwidth - 16\tabcolsep) * \real{0.1020}}
  >{\raggedleft\arraybackslash}p{(\columnwidth - 16\tabcolsep) * \real{0.1180}}
  >{\raggedleft\arraybackslash}p{(\columnwidth - 16\tabcolsep) * \real{0.1180}}
  >{\raggedleft\arraybackslash}p{(\columnwidth - 16\tabcolsep) * \real{0.0776}}
  >{\raggedleft\arraybackslash}p{(\columnwidth - 16\tabcolsep) * \real{0.1207}}@{}}
\toprule\noalign{}
\begin{minipage}[b]{\linewidth}\raggedright
\textbf{Sınıf}
\end{minipage} & \begin{minipage}[b]{\linewidth}\raggedleft
\textbf{AB'ye ihracat (€ mn, 2025)}
\end{minipage} & \begin{minipage}[b]{\linewidth}\raggedleft
\textbf{GB ve İTÜ payı}
\end{minipage} & \begin{minipage}[b]{\linewidth}\raggedleft
\textbf{S4'te AB tarifesi}
\end{minipage} & \begin{minipage}[b]{\linewidth}\raggedright
\textbf{Menşe riski S2/S3}
\end{minipage} & \begin{minipage}[b]{\linewidth}\raggedleft
\textbf{AB'ye ihracat S2 (\$ mn)}
\end{minipage} & \begin{minipage}[b]{\linewidth}\raggedleft
\textbf{AB'ye ihracat S4 (\$ mn)}
\end{minipage} & \begin{minipage}[b]{\linewidth}\raggedleft
\textbf{S2 ÷ S4}
\end{minipage} & \begin{minipage}[b]{\linewidth}\raggedleft
\textbf{İstihdam S4 (kişi)}
\end{minipage} \\
\midrule\noalign{}
\endhead
\bottomrule\noalign{}
\endlastfoot
16.21 Ahşap levha ve kontrplak & 194 & \%100 & \%6,9 & D/D & −22 & −56 & \%39 & −68 \\
16.23 Doğrama ve ahşap yapı & 57 & \%100 & \%1,2 & D/D & −4 & −5 & \%73 & −61 \\
16.24 Palet ve ahşap ambalaj & 12 & \%100 & \%3,3 & D/D & −1 & −2 & \%53 & +20 \\
\textbf{Toplam} & \textbf{263} & --- & --- & --- & \textbf{−27} & \textbf{−64} & \textbf{\%42} & \textbf{−110} \\
\end{longtable}

\endgroup

\emph{Kaynak:} \etiketT{} f3d, f3e. Menşe riski D/O/Y: düşük, orta, yüksek (\kod{veri/\allowbreak{}f3\_\allowbreak{}sinif\_\allowbreak{}varsayim.\allowbreak{}csv}, \etiketH{}). Seçim: AB'ye ihracatının en az yarısı GB veya İTÜ rejiminde, €10 mn üzerinde ve S4'te AB tarifesi en az \%1 olan sınıflar.

S4'te ahşap levhanın AB'ye ihracatı 56 milyon dolar azalıyor, ama sınıfın istihdam kaybı yalnızca 68 kişi. Bunun iki nedeni var. Birincisi, AB sınıfın ihracatının yalnızca \%16 kadarını alıyor; AB'den kaybedilen satışın bir kısmı Orta Doğu ve Kuzey Afrika'ya kayıyor. İkincisi, S4'te Türkiye'nin c2 tarifesi AB'den gelen levhayı pahalandırıyor ve iç pazarda yerli üreticiye pay kazandırıyor. AB pazarındaki rakiplerin konumu da önemli. AB'nin AB dışından levha tedarikinde Çin (\%24) ve Ukrayna (\%15) ilk iki sırada. Ukrayna, AB ile derin ve kapsamlı serbest ticaret anlaşması sayesinde sıfır tarifeyle giriyor. Anlaşmasız çıkışta Türk levhası, Romanya ve Bulgaristan gibi yakın pazarlarda Ukrayna'ya karşı \%7'lik bir dezavantajla rekabet eder.

Tarifesi sıfır olan sınıflar da S4'te AB'ye ihracat kaybediyor: kâğıt bölümünün kaybı 63 milyon dolar. Bu kaybın kaynağı sınır sürtünmesi. Gümrük Birliği'nde A.TR dolaşım belgesiyle yapılan sevkiyat, anlaşmasız çıkışta tam bir gümrük beyanına, menşe belgesine ve AB'nin üçüncü ülkelere uyguladığı belge ve risk denetimine tabi oluyor. Model bu maliyeti sevkiyat değerinin \%0,5--2'si olarak alıyor ve genel denge modelinin toplam kaybına göre ölçekliyor. Kâğıt ve ambalaj gibi değeri düşük, hacmi büyük ve tam zamanında teslim edilen ürünlerde bu maliyet görece daha ağır olabilir \etiketH{}.

\hypertarget{nace_kartlar_16-18-32-ticaret-savunmasux131-ve-rus-odunu}{%
\subsection{3.2 Ticaret savunması ve Rus odunu}\label{nace_kartlar_16-18-32-ticaret-savunmasux131-ve-rus-odunu}}

Bu parçanın AB ile ilişkideki en hassas noktası tarife değil, Rus odunu. Rusya ve Belarus, Türkiye'nin odun ve odun ürünleri ithalatında büyük bir yer tutuyor. 2024'te Türkiye'nin kontrplak ithalatının \%45 kadarı Rusya'dan geldi. Kaplamada (4408) Rusya'nın payı \%44, biçilmiş kerestede (4407) \%12. Bu ithalatın çıkış senaryolarıyla bir ilgisi yok; Türkiye Rusya'ya yaptırım uygulamıyor ve bu ticaret yasal. Ama AB tarafında üç ayrı düzenleme bu ticareti hedef alıyor.

Birincisi ticaret savunması. AB, Rus huş kontrplağına 2021'de \%14--16 dolayında anti-damping vergisi koydu. 2023'te Türkiye ve Kazakistan'dan gelen huş kontrplağı ithalatının hızla arttığını tespit ederek bir dolanma soruşturması açtı ve Mayıs 2024'te vergiyi bu iki ülkeden gönderilen huş kontrplağına genişletti (Uygulama Tüzüğü 2024/1287) \etiketV{}. Genişletme, ürünün beyan edilen menşeinden bağımsız olarak Türkiye'den gönderilen bütün huş kontrplağını kapsıyor; dolanmaya katılmadığını kanıtlayan üreticiler muafiyet isteyebiliyor. Önlemin beş yıllık süresi Kasım 2026'da doluyordu. AB'li üreticilerin başvurusu üzerine Komisyon 6 Kasım 2026'da süre sonu gözden geçirmesi başlattı ve vergiler 12--15 ay sürecek soruşturma boyunca yürürlükte kalıyor \etiketV{}. Yani bu genişletme, Gümrük Birliği'nin içinde ve sürerken yapıldı. Gümrük Birliği, AB'nin Türkiye'den gelen mallara ticaret savunması uygulamasını engellemiyor; çıkış bu açıdan yeni bir şey getirmiyor, yalnızca AB'nin elini serbestleştiriyor.

İkincisi yaptırımlar. AB 2022'den beri Rus odun ürünlerinin ithalatını yasaklıyor. Komisyon Mart 2025'te AB'li işletmelere, kontrplak ithalatının yaptırımları dolanmak için kullanılma riskinin yüksek olduğunu bildiren bir uyarı yayımladı \etiketV{}. Türkiye'de Rus odunundan yapılan bir levha, menşe kurallarına göre Türk malı ve yaptırımların doğrudan kapsamında değil. Ama AB, demir ve çelikte yaptığı gibi, üçüncü ülkelerde işlenmiş Rus girdisini de yasak kapsamına alan bir kural getirebilir \etiketH{}. Hasmane bir AB için bu, çıkıştan sonra Türk levha ve kontrplağına karşı kullanılabilecek ucuz bir araç.

Üçüncüsü ormansızlaşma tüzüğünün ülke sınıflandırması. Komisyon Rusya ve Belarus'u yüksek riskli ülkeler arasında saydı \etiketV{}. Türkiye'de işlenen bir ürünün içinde Rus odunu varsa, AB'li ithalatçı bu odun için de coğrafi konum verisi vermek ve yüksek riskli menşe için güçlendirilmiş özen yükümlülüğünü yerine getirmek zorunda. Pratikte bu, AB'li alıcıların Rus odunu içeren Türk ürününden kaçınması demek \etiketH{}.

Bu üç düzenleme birlikte okunduğunda risk açık. AB'ye satılan Türk levha, kontrplak ve kereste, tedarik zincirinde Rus odunu bulunmadığını belgeleyebildiği ölçüde güvende. Bunu belgeleyemeyen ihracatçı, hasmane bir AB'nin seçici ticaret savunması ve yaptırım araçlarına açık. Model bu riski ölçmüyor: ne genel denge modelinde ne de kısmi denge kanallarında ticaret savunması var. Riskin büyüklüğü, ihracatçıların odun tedarikinde Rus ve Belarus payına bağlı ve bu veri firma düzeyinde yok (VB-40). Önerilen yanıt iki parçalı. Birincisi, AB'ye ihracat yapan levha ve kontrplak üreticilerinin Rus ve Belarus odunundan ayrışmış, ithalat kayıtlarıyla doğrulanabilir bir tedarik zinciri kurması ve gözden geçirme soruşturmasında Türk üreticilerin muafiyet başvurularının desteklenmesi (E-16-02). İkincisi, S2 müzakeresinde ticaret savunması soruşturmalarından önce zorunlu istişare ve dolanma soruşturmalarında taraflar arasında bilgi paylaşımı (E-16-04).

AB'nin ticaret savunması pratiği Türk sektörü hakkında olumlu bir bilgi de içeriyor. Komisyon, Çin'den gelen çok katmanlı ahşap döşemeye yönelik anti-damping soruşturmasında Çin'in maliyetlerinin yerine Türkiye'nin maliyetlerini referans aldı ve Çinli ihracatçıların itirazına rağmen bu seçimi korudu (Uygulama Tüzüğü 2025/1342) \etiketV{}. Komisyon, Türk ahşap sektörünü piyasa koşullarında işleyen bir rekabetçi olarak görüyor. Hasmane bir ilişkide bile Türkiye'nin savunması bu görüntüyü korumaktan geçiyor.

\hypertarget{nace_kartlar_16-18-33-ormansux131zlaux15fma-ve-ambalaj-guxfcmruxfck-birliux11finden-baux11fux131msux131z-yuxfckuxfcmluxfcluxfckler}{%
\subsection{3.3 Ormansızlaşma ve ambalaj: Gümrük Birliği'nden bağımsız yükümlülükler}\label{nace_kartlar_16-18-33-ormansux131zlaux15fma-ve-ambalaj-guxfcmruxfck-birliux11finden-baux11fux131msux131z-yuxfckuxfcmluxfcluxfckler}}

AB'nin ormansızlaşma tüzüğü (EUDR, Tüzük 2023/1115), iki kez ertelendikten sonra 30 Aralık 2026'da büyük ve orta ölçekli işletmeler, 30 Haziran 2027'de mikro ve küçük işletmeler için uygulanmaya başlıyor \etiketV{}. Komisyon Mayıs 2026'daki sadeleştirme paketinde yeni bir erteleme önermedi \etiketV{}. Tüzük, belirli ürünleri AB pazarına süren işletmenin, ürünün ormansızlaşmaya yol açmamış bir arazide ve yasal olarak üretildiğini, üretim parselinin coğrafi konumuyla birlikte beyan etmesini istiyor. Bu parçada kereste, ahşap levha, palet, doğrama, kâğıt hamuru, kâğıt, karton ve kâğıt ürünleri kapsamda. Basılı ürünler (Fasıl 49), Aralık 2025'teki değişiklikle (Tüzük 2025/2650) kapsamdan çıkarıldı \etiketV{}. Tamamen geri dönüştürülmüş liften yapılan ürünler ve yalnızca başka bir ürünü taşımak ya da korumak için kullanılan ambalaj da kapsam dışında \etiketV{}.

AB'ye bu parçada satılan ürünlerin, kapsamdaki KN kodlarına göre ağırlıklı \%90 kadarı EUDR'ye giriyor. Bu bir üst sınır, çünkü Türk kâğıt ve oluklu ambalaj üretiminin büyük kısmı atık kâğıttan yapılıyor ve tamamen geri dönüştürülmüş liften yapılan ürün kapsam dışında \etiketH{}. Gerçek payı ölçmek için ihracatın lif bileşimine ilişkin veri gerekiyor (VB-41). Kapsamda kalan ürünler ağırlıkla üç grupta: Türk odunundan yapılan levha ve kereste, ithal kâğıt hamurundan yapılan hijyen kâğıtları ve yerli üretimi olmayan kâğıt kalitelerinden yapılan ürünler.

EUDR'nin Türkiye için iki önemli özelliği var. Birincisi, yükümlülük AB'li ithalatçıda; Türk ihracatçının hukuki bir yükümlülüğü yok, ama coğrafi konum verisini vermezse ürünü satamaz. İkincisi, Komisyon'un Mayıs 2025'te yayımladığı ülke sınıflandırmasında (Uygulama Tüzüğü 2025/1093) Türkiye düşük riskli ülkeler arasında \etiketV{}. Düşük riskli ülkeden yalnızca ürün alan ithalatçı risk değerlendirmesi ve risk azaltımı yapmak zorunda değil; yalnızca bilgiyi toplayıp beyan ediyor. Denetim oranı düşük riskli ülkelerde \%1, standart riskte \%3, yüksek riskte \%9 \etiketV{}. Avrupa Parlamentosu sınıflandırmaya Temmuz 2025'te itiraz etti, ama bu itiraz bağlayıcı değildi. Komisyon sınıflandırmayı 2026'da güncelleyeceğini açıkladı \etiketV{}. Hasmane bir AB için Türkiye'yi standart risk sınıfına almak, gerekçesi kolayca bulunabilecek, maliyeti düşük bir baskı aracı. Türkiye'nin buna karşı en güçlü kozu, verisinin güvenilirliği.

Bu noktada Türkiye'nin yapısal bir üstünlüğü var. Türkiye'de ormanların neredeyse tamamı devletin ve odun satışları Orman Genel Müdürlüğü (OGM) üzerinden yapılıyor \etiketV{}. OGM'nin amenajman planlarındaki bölme ve bölmecik sınırları, EUDR'nin istediği parsel verisinin hazır bir kaynağı. Her dikili ve istif satışında bu verinin alıcıya EUDR'ye uygun biçimde verilmesi, Türk odunundan yapılan bütün ürünler için coğrafi konum sorununu tek bir idari kararla çözüyor (E-16-01). İthal odun ve kâğıt hamurunda veri tedarikçiden gelmek zorunda; büyük hamur üreticileri bu veriyi zaten hazırlıyor \etiketH{}. İlk parçada gıda için önerilen EUDR uyum masası (E-10-03) bu parçada ahşap, kâğıt ve mobilyaya genişletiliyor.

Ambalaj tarafında AB'nin yeni Ambalaj ve Ambalaj Atıkları Tüzüğü (PPWR, Tüzük 2025/40) 12 Ağustos 2026'dan beri uygulanıyor \etiketV{}. Tüzük, AB pazarına ambalaj süren üreticiden uygunluk değerlendirmesi ve teknik dosya istiyor; ithal ambalajda bu dosyanın varlığından ithalatçı sorumlu. Gıdayla temas eden ambalajda PFAS sınırları da aynı tarihte başladı. Etiketleme, geri dönüştürülebilirlik ve geri dönüştürülmüş içerik yükümlülükleri 2028--2030 arasında devreye giriyor \etiketV{}. Türk oluklu ambalaj ve gıda ambalajı kâğıdı ihracatçıları için bunlar uyum maliyeti, ama karşılıklı tanıma gerektirmiyor: uygunluk beyanı üreticinin kendi beyanı. Bu yüzden PPWR'nin maliyeti senaryolar arasında değişmiyor. Önerilen yanıt, PFAS'sız kaplamaya geçiş ve uygunluk dosyası için ihracatçılara bir rehber ve Türkiye'nin ambalaj atığı mevzuatının PPWR ile uyumu (E-16-03).

Bu iki düzenleme senaryodan bağımsız. EUDR ve PPWR'nin yükü S2'de de S4'te de aynı. Çıkışın bunlarla ilişkisi iki noktada. Birincisi, hasmane bir AB'nin ülke sınıflandırması ve denetim sıklığı üzerindeki takdiri. İkincisi, S2'de bir düzenleyici işbirliği bölümünün sağlayabileceği kolaylık: örneğin OGM'nin veri sisteminin EUDR açısından tanınması. Komisyon üçüncü ülkelerle EUDR işbirliği ve ortaklık çerçeveleri kuruyor \etiketV{}. S2 müzakeresinde bu, Türkiye için düşük maliyetli ama değerli bir talep.

\hypertarget{nace_kartlar_16-18-34-iux307thalat-tarafux131-c2-tarifesi-kuxe2ux11fux131tta-iuxe7-ikame-ve-tuxfcrkiyenin-korunma-uxf6nlemleri}{%
\subsection{3.4 İthalat tarafı: c2 tarifesi, kâğıtta iç ikame ve Türkiye'nin korunma önlemleri}\label{nace_kartlar_16-18-34-iux307thalat-tarafux131-c2-tarifesi-kuxe2ux11fux131tta-iuxe7-ikame-ve-tuxfcrkiyenin-korunma-uxf6nlemleri}}

Bu parçanın senaryo sonuçlarını belirleyen kanal ithalat. Türkiye kâğıt ve kartonda net ithalatçı ve AB'den yılda 1.589 milyon dolarlık kâğıt ve karton alıyor. S4'te c2 rejimi, Faz 2A'da tanımlandığı gibi, yurt içi üretimin ithalatı karşılayabildiği satırlarda AB menşeli ürüne de MFN tarifesini ve İGV'yi uyguluyor; AB'nin baskın tedarikçi olduğu ve yurt içi kabiliyetin zayıf olduğu satırlarda İGV kalkıyor. Kâğıtta MFN tarifesi sıfır, ama bazı satırlarda İGV var. Bu yüzden AB'den gelen kâğıt ve kartonun ortalama tarifesi S4'te \%3,5 oluyor; oluklu kutuda \%24,9, ahşap levhada \%15,1.

Sonuç, AB'den kâğıt ithalatının azalması ve boşluğun bir kısmının yerli üretimle dolması. S4'te AB'den kâğıt ve karton ithalatı 239 milyon dolar azalıyor; yerli üretim 155 milyon dolar artıyor. İstihdam etkisi kâğıt ve kartonda +409, diğer kâğıt ürünlerinde +145 kişi. Bu artışlar, S4'te bu parçanın net istihdam etkisini artıya çeviriyor. Bedeli tüketici ödüyor, ama küçük: parçanın S4'teki TÜFE etkisi 0,9 baz puan. Aşağıdaki tablo bu etkinin sınıflara dağılımını veriyor.

\Needspace{7\baselineskip}

\textbf{Tablo 3.2 --- İthalat tarafı: TÜFE'ye birinci tur etkisi en büyük sınıflar (bp, P50)}

\nopagebreak[4]

\begingroup\gbtablo\footnotesize

\begin{longtable}[]{@{}
  >{\raggedright\arraybackslash}p{(\columnwidth - 16\tabcolsep) * \real{0.1731}}
  >{\raggedleft\arraybackslash}p{(\columnwidth - 16\tabcolsep) * \real{0.1152}}
  >{\raggedleft\arraybackslash}p{(\columnwidth - 16\tabcolsep) * \real{0.1442}}
  >{\raggedleft\arraybackslash}p{(\columnwidth - 16\tabcolsep) * \real{0.1369}}
  >{\raggedleft\arraybackslash}p{(\columnwidth - 16\tabcolsep) * \real{0.1444}}
  >{\raggedleft\arraybackslash}p{(\columnwidth - 16\tabcolsep) * \real{0.0481}}
  >{\raggedleft\arraybackslash}p{(\columnwidth - 16\tabcolsep) * \real{0.0481}}
  >{\raggedright\arraybackslash}p{(\columnwidth - 16\tabcolsep) * \real{0.0940}}
  >{\raggedleft\arraybackslash}p{(\columnwidth - 16\tabcolsep) * \real{0.0960}}@{}}
\toprule\noalign{}
\begin{minipage}[b]{\linewidth}\raggedright
\textbf{Sınıf}
\end{minipage} & \begin{minipage}[b]{\linewidth}\raggedleft
\textbf{TÜFE ağırlığı}
\end{minipage} & \begin{minipage}[b]{\linewidth}\raggedleft
\textbf{Tüketimde ithal pay}
\end{minipage} & \begin{minipage}[b]{\linewidth}\raggedleft
\textbf{İthalatta AB payı}
\end{minipage} & \begin{minipage}[b]{\linewidth}\raggedleft
\textbf{S4'te AB'ye tarife (c2)}
\end{minipage} & \begin{minipage}[b]{\linewidth}\raggedleft
\textbf{S2}
\end{minipage} & \begin{minipage}[b]{\linewidth}\raggedleft
\textbf{S4}
\end{minipage} & \begin{minipage}[b]{\linewidth}\raggedright
\textbf{Önlem (S4)}
\end{minipage} & \begin{minipage}[b]{\linewidth}\raggedleft
\textbf{S4, önlemle}
\end{minipage} \\
\midrule\noalign{}
\endhead
\bottomrule\noalign{}
\endlastfoot
17.12 Kâğıt ve karton & \%0,00 & \%3 & \%63 & \%3,5 & 0,13 & 0,48 & −0,20 & 0,28 \\
17.29 Diğer kâğıt ürünleri & \%0,00 & \%2 & \%77 & \%16,0 & 0,02 & 0,13 & --- & 0,13 \\
17.21 Oluklu ambalaj & \%0,00 & \%1 & \%76 & \%24,9 & 0,01 & 0,12 & --- & 0,12 \\
16.21 Ahşap levha ve kontrplak & \%0,00 & \%38 & \%36 & \%15,1 & 0,01 & 0,07 & --- & 0,07 \\
17.22 Hijyen kâğıtları ve bezler & \%0,41 & \%1 & \%51 & \%6,7 & 0,01 & 0,06 & --- & 0,06 \\
17.11 Kâğıt hamuru & \%0,00 & \%5 & \%39 & \%0,0 & 0,02 & 0,03 & --- & 0,03 \\
16.24 Palet ve ahşap ambalaj & \%0,00 & \%95 & \%62 & \%24,8 & 0,00 & 0,02 & --- & 0,02 \\
16.29 Diğer ağaç ürünleri & \%0,00 & \%95 & \%31 & \%4,8 & 0,00 & 0,01 & --- & 0,01 \\
Diğer 10 sınıf & \%0,09 & --- & --- & --- & 0,01 & 0,02 & --- & 0,02 \\
\textbf{16--18 toplamı} & \textbf{\%0,51} & --- & --- & --- & \textbf{0,21} & \textbf{0,93} & \textbf{−0,20} & \textbf{0,73} \\
\end{longtable}

\endgroup

\emph{Kaynak:} \etiketT{} f3e (doğrudan kanal COICOP 2018 sepeti, dolaylı kanal I-O 2023 Leontief fiyat modeli). c2: S4'te AB menşeli mala uygulanacak MFN + İGV (İGV, AB'nin baskın, yerli üretimin zayıf olduğu satırlarda herkes için kalkar; Faz 2A). Önlemler Bölüm 4'te.

Bu ikame tablosunun iki sınırı var. Birincisi aşağı akış. Kâğıt ve karton, oluklu ambalajın, matbaacılığın ve kâğıt ürünlerinin ana girdisi. Bu üç sınıf, kâğıt ve karton sınıfının (13,6 bin kişi) yedi katından fazla istihdam sağlıyor. Faz 2A'nın girdi-çıktı fiyat modeline göre, S4'te AB'den gelen bütün girdilere MFN uygulandığında matbaacılığın üretici fiyatı \%0,18, MFN ile birlikte İGV de uygulandığında \%0,58 artıyor; kâğıt ürünlerinde aralık \%0,15--0,52 \etiketT{}. c2 bu iki uç arasında. Bu, aşağı akış sınıflarının marjlarında bir sıkışma demek; model bu sıkışmanın istihdam etkisini ölçmüyor. Model, aşağı akışın kaybını görmezken yukarı akışın kazancını görüyor. Bu yüzden S4'teki net artı bir üst sınır.

İkincisi c2'nin kendi kuralı. Yurt içi kabiliyet bayrağı HS6 düzeyinde, üretimin ithalata oranıyla hesaplanıyor. Bazı HS6 satırları Türkiye'de üretilen ve üretilmeyen kaliteleri birlikte içeriyor. Örneğin sıvı ambalaj kartonu (KN 4811 59), Türkiye'de üretilen başka kaplamalı kâğıtlarla aynı satırda olduğu için ``güçlü'' sayılıyor ve S4'te AB menşelisi \%20 İGV ile karşılaşıyor (MFN tarifesi sıfır). Oysa süt ve meyve suyu kutusunun bu girdisi Türkiye'de üretilmiyor \etiketH{}. Ağartılmamış kraftliner (4804 11) de benzer: ihracatlık meyve ve sebze kutusunun girdisi, yerli üretimi sınırlı, ama S4'te AB menşelisine \%10 İGV geliyor. Bu iki satırda AB'den ithalat 2024'te yaklaşık 130 milyon dolar \etiketT{}. Bu kalitelerde İGV, yerli üreticiyi korumadan gıda ve içecek ambalajını pahalandırıyor. Önerilen düzeltme, yalnızca S4'te ve bu tür kalitelerde İGV'nin bütün menşeler için sıfırlanması ve c2 satır listesinin 8--12 haneli üretim verisiyle yeniden gözden geçirilmesi (E-17.12-01, VB-43). Düzeltme TÜFE etkisini yaklaşık 0,2 baz puan azaltıyor ve AB dışı menşelerden alınan İGV'nin kaybı nedeniyle yılda yaklaşık 0,7 milyar TL'ye mal oluyor \etiketT{}.

Çıkıştan bağımsız ama aynı soruyu soran bir başka tablo, Türkiye'nin kâğıtta bugün uyguladığı korunma önlemleri. Atık Kağıt Geri Dönüşümcüler Derneği'nin başvurusu üzerine Ocak 2024'te açılan soruşturma sonunda Türkiye, testliner ve fluting cinsi oluklu mukavva kâğıdı ithalatına üç yıl süreyle ton başına 87, 86 ve 85 dolarlık ek mali yükümlülük koydu (Cumhurbaşkanı Kararı 8651, Resmî Gazete 28 Haziran 2024) \etiketV{}. Önlem 28 Temmuz 2024'te yürürlüğe girdi ve Temmuz 2027'de, yani çıkıştan önce sona eriyor. Kararın tarife kontenjanıyla önlemden muaf tuttuğu ülkeler yalnızca gelişmekte olan ülkeler; AB üyeleri listede yok \etiketV{}. Oysa 2019--2023'te ithalatın yarıdan fazlası Almanya, Polonya, İsrail ve Finlandiya'dan geliyordu \etiketV{}. Yani Gümrük Birliği sürerken Türkiye, AB'den gelen bir sanayi ürününe DTÖ korunma önlemi uyguladı. Aralık 2025'te kaplamalı kâğıt ve çok katlı kaplamalı kartonda (4810 29 ve 4810 92) yeni bir korunma soruşturması açıldı ve süresi altı ay uzatıldı \etiketV{}. Karar 2027'nin ilk yarısında bekleniyor \etiketT{}.

Bu önlemler, c2 ile aynı dağılım sorusunu soruyor. Kaplamalı karton ve oluklu mukavva kâğıdı, ambalaj ve baskı sanayilerinin girdisi; bu sanayiler hem daha çok kişi çalıştırıyor hem de ihracatçıların ürünlerini taşıyor. Dahilde işleme rejimi ihracata giden ambalajı ek yükümlülükten kurtarıyor, ama iç pazara satılan ambalajın ve basılı ürünün maliyeti artıyor. Bu parçanın önerisi, önlem kararlarında aşağı akış kullanıcılarının ve tüketicinin çıkarının ayrı ve yayımlanan bir değerlendirmeyle tartılması (E-16-04). Gümrük Birliği'nden çıkış bu tartıyı değiştirmiyor; c2 ile S4'te yalnızca daha sık yapılması gereken bir tartı haline getiriyor.

Kâğıt hamurunda tablo farklı. Türkiye'nin hamur ithalatı 1.068 milyon dolar ve bunun \%39 kadarı AB'den. Hamurda MFN tarifesi sıfır ve c2 bunu değiştirmiyor. AB'den ithalat sınır sürtünmesiyle bir miktar azalıyor, ama Türkiye'de hamur üretimi yok denecek kadar az olduğu için boşluğu yerli üretim değil Brezilya, ABD ve Uruguay dolduruyor. Önceki parçalarda model, kâğıt hamurunda da yerli ikame varsayıyordu ve bu, S2'de 129, S4'te 98 kişilik anlamsız bir istihdam artışı üretiyordu. Bu parçada iç ikame tavanıyla bu artış sıfırlandı (V-44, Ek D).

\hypertarget{nace_kartlar_16-18-35-s2nin-sux131navux131-kolay-menux15fe-sux131nux131r-maliyeti}{%
\subsection{3.5 S2'nin sınavı: kolay menşe, sınır maliyeti}\label{nace_kartlar_16-18-35-s2nin-sux131navux131-kolay-menux15fe-sux131nux131r-maliyeti}}

S2'de bu parçanın sınavı kolay. Revize PEM kurallarında Fasıl 44, 47 ve 48'in ürünlerinin büyük çoğunluğu için menşe kuralı başlık değişikliği \etiketV{}. İthal tomruktan kereste, ithal yonga ve kaplamadan levha, ithal hamurdan kâğıt, ithal kartondan kutu Türk menşeli sayılıyor. Geri ödeme yasağı yalnızca tekstil fasıllarında (50--63) uygulanıyor; bu parçanın ürünlerinde dahilde işleme rejimiyle ithal edilen girdi sorun yaratmıyor. Bu yüzden menşe uyum oranı yüksek (0,95) ve uyumlu ihracatın menşe maliyeti düşük (\%1) varsayıldı \etiketH{}.

Kâğıtta ise menşe hiç sorun değil, çünkü AB'nin MFN tarifesi sıfır: ihracatçı tercih istemez, menşe ispatına da katlanmaz. Bu parçada model bu noktada düzeltildi (Ek B, Ek D). S2'deki kayıp bu yüzden çoğunlukla sınır maliyetinden geliyor. AB'ye ihracat 75 milyon dolar azalıyor; S4'teki kaybın \%56 kadarı. İki senaryo arasındaki fark, tarife duvarının alçaklığı nedeniyle dar. İthalat tarafında ise menşe önemli. Türkiye bazı kâğıt satırlarında İGV uyguladığı için AB'li ihracatçı, bu satırlarda sıfır tarifeden yararlanmak üzere menşe ispatlamak zorunda. Bu, S2'de de AB'den kâğıt ithalatını bir miktar azaltıyor ve küçük bir yerli ikame yaratıyor; kâğıt bölümünün S2'deki istihdam etkisi +230 kişi. Sonuçta bu parçada iyi bir anlaşmanın değeri, önceki parçalardaki gibi tarifede değil. Değer, ticaret savunması, düzenleyici işbirliği ve sınır kolaylaştırması bölümlerinde.

\hypertarget{nace_kartlar_16-18-4-strateji-neyi-korumalux131-neyi-geliux15ftirmeli-nasux131l-muxfczakere-etmeli}{%
\section{4. Strateji: neyi korumalı, neyi geliştirmeli, nasıl müzakere etmeli?}\label{nace_kartlar_16-18-4-strateji-neyi-korumalux131-neyi-geliux15ftirmeli-nasux131l-muxfczakere-etmeli}}

\hypertarget{nace_kartlar_16-18-41-duruux15flar}{%
\subsection{4.1 Duruşlar}\label{nace_kartlar_16-18-41-duruux15flar}}

Duruşlar önceki parçalardaki yöntemle atandı: önce kabiliyet, AB maruziyeti, ithalat bağımlılığı, enflasyon duyarlılığı ve stratejik önem skorlarına göre kurallarla, sonra her sınıf için ayrı ayrı gözden geçirilerek. Bu parçada gözden geçirme kural tabanlı sınıflamayı belirgin biçimde değiştirdi. Kurallar, toplam ihracatının en az \%8'i AB'ye giden ve ormansızlaşma tüzüğüne (EUDR) giren ürünlerden oluşan on sınıfı birincil arketip olarak ``Yeşil Mutabakat'a maruz'' (G) işaretliyor ve onlar için ``Dönüştür'' duruşunu öneriyordu. Bu işaretleme ilk iki parçada SKDM için anlamlıydı, çünkü SKDM ağır bir karbon maliyeti ve üretim teknolojisinin değişmesini gerektiriyor. EUDR ise bir bilgi yükümlülüğü: ürünü değil, ürünün kaynağına ilişkin veriyi değiştiriyor ve bütün senaryolarda, bütün menşelere aynı biçimde uygulanıyor. Bu yüzden G, AB'ye ihracatı küçük olan altı sınıfta kaldırıldı; kâğıt ve karton, oluklu ambalaj ve hijyen kâğıtlarında ikincil arketipe indi. Gıdayla temas eden kâğıdı ambalaj tüzüğüne (PPWR) uyum gerektiren diğer kâğıt ürünlerinde (17.29) birincil arketip olarak kaldı, ama duruşu ``Koru'' yapıldı.

\Needspace{7\baselineskip}

\textbf{Tablo 4.1 --- Sınıfların duruşa göre dağılımı}

\nopagebreak[4]

\begingroup\gbtablo\footnotesize

\begin{longtable}[]{@{}
  >{\raggedright\arraybackslash}p{(\columnwidth - 8\tabcolsep) * \real{0.0737}}
  >{\raggedleft\arraybackslash}p{(\columnwidth - 8\tabcolsep) * \real{0.0618}}
  >{\raggedright\arraybackslash}p{(\columnwidth - 8\tabcolsep) * \real{0.5349}}
  >{\raggedleft\arraybackslash}p{(\columnwidth - 8\tabcolsep) * \real{0.1197}}
  >{\raggedleft\arraybackslash}p{(\columnwidth - 8\tabcolsep) * \real{0.2099}}@{}}
\toprule\noalign{}
\begin{minipage}[b]{\linewidth}\raggedright
\textbf{Duruş}
\end{minipage} & \begin{minipage}[b]{\linewidth}\raggedleft
\textbf{Sınıf}
\end{minipage} & \begin{minipage}[b]{\linewidth}\raggedright
\textbf{Sınıflar}
\end{minipage} & \begin{minipage}[b]{\linewidth}\raggedleft
\textbf{Çalışan (bin)}
\end{minipage} & \begin{minipage}[b]{\linewidth}\raggedleft
\textbf{AB'ye ihracat (\$ mn, 2024)}
\end{minipage} \\
\midrule\noalign{}
\endhead
\bottomrule\noalign{}
\endlastfoot
Koru & 13 & 16.10, 16.23, 16.24, 16.29, 17.11, 17.21, 17.23, 17.24, 17.29, 18.11, 18.12, 18.14, 18.20 & 179,2 & 559 \\
Geliştir & 5 & 16.21, 16.22, 17.12, 17.22, 18.13 & 66,3 & 819 \\
\end{longtable}

\endgroup

\emph{Kaynak:} \kod{veri/\allowbreak{}nace\_\allowbreak{}matris.\allowbreak{}csv} (\kod{durus}; kural tabanlı öneri ve uzman düzeltmesi, \kod{uzman\_\allowbreak{}duzeltmesi} sütunu).

``Geliştir'' duruşundaki 5 sınıfın üçü büyük. Ahşap levha (16.21) küresel rekabetçi bir ihracatçı ve AB'ye bağımlılığı düşük. Strateji Orta Doğu, Kuzey Afrika ve Orta Asya pazarlarını büyütmek, AB pazarında ise EUDR verisiyle ve Rus odunundan ayrışmış tedarikle güvenilirliği korumak. Hijyen kâğıtları ve bezler (17.22) aynı profilde: dünyanın büyük bebek bezi ve hijyen kâğıdı ihracatçılarından biri, ilk pazarı Irak. Kâğıt ve karton (17.12) ise ithalata bağımlı bir sınıf ve kurallar onu yerlileştirme adayı (C) olarak işaretliyor. Buradaki ``Geliştir'', tarife korumasıyla değil, geri dönüştürülmüş lif temelinde yatırımla ve atık kâğıt toplamasının genişletilmesiyle geliştirmek anlamına geliyor. Sebebi Bölüm 3.4'te anlatıldı: kâğıdı koruyan her önlem, ondan yedi kat fazla kişi çalıştıran aşağı akışa maliyet yüklüyor. Ahşap döşeme (16.22) ve baskı öncesi hizmetler (18.13) küçük, ama ihracat profilleri rekabetçi.

``Koru'' duruşundaki 13 sınıf bu parçada iki farklı grubu içeriyor. Birinci grup, çıkıştan etkisi küçük ve çıkışa özgü bir önlem gerektirmeyen sınıflar: kereste, doğrama, palet, diğer ağaç ürünleri, kırtasiye, duvar kâğıdı ve basım sınıfları. Burada ``Koru'' bir izleme duruşu. İkinci grup, AB'ye entegre ama tarifeden etkilenmeyen sınıflar: oluklu ambalaj (17.21) ve diğer kâğıt ürünleri (17.29). Oluklu ambalajı kurallar ikincil olarak ``AB'ye entegre'' (A) diye işaretliyordu ve bu, uzman düzeltmesiyle birincil arketip yapıldı; AB'ye ihracatın payı \%39. Ama AB'nin oluklu kutudaki MFN tarifesi sıfır ve sınıfın AB'deki konumu yakınlıktan ve hızdan geliyor. Bu sınıf için asıl gündem ambalaj tüzüğüne (PPWR) uyum ve girdi kâğıdının maliyeti. Kâğıt hamuru (17.11) da bu grupta: tarife zaten sıfır ve yurt içi üretim yok; strateji tedarik çeşitlendirme.

Bu parçada ``Dönüştür'' ve ``Yönetilen uyum'' duruşu önerilmiyor. Bu bir olumsuz bulgu olarak da okunmalı. Kereste ve doğramacılık, AB'deki muadillerinin üçte biri ile altıda biri arasında bir verimlilikle çalışıyor ve atölye düzeyinde kalmış durumda. Bu sınıfların sorunu Gümrük Birliği değil ve çıkış politikası onlar için bir çözüm sunmuyor. Bunlar sanayi ve KOBİ politikasının konusu.

\hypertarget{nace_kartlar_16-18-42-uxf6nerilen-paket-ve-sux131ralamasux131}{%
\subsection{4.2 Önerilen paket ve sıralaması}\label{nace_kartlar_16-18-42-uxf6nerilen-paket-ve-sux131ralamasux131}}

Paket, bölüm düzeyinde 5 ortak eylemden ve sınıflara özgü 2 eylemden oluşuyor. Önceki parçalardan farklı olarak paketin ağırlığı S4 tamponlarında değil, senaryodan bağımsız hazırlıkta. Eylemler yine üç katmanda sıralanıyor.

İlk katman hemen başlayacak ve hiçbir senaryoda pişmanlık yaratmayacak hazırlık. Bu parçada bu katmanın bir takvimi var: EUDR 30 Aralık 2026'da uygulanmaya başlıyor. EUDR ahşap ve kâğıt veri hizmeti (E-16-01) bu yüzden paketin en acil eylemi. OGM'nin satış şartnamesine coğrafi konum verisinin eklenmesi bir genelgeyle yapılabilir ve ilk parçadaki EUDR uyum masası (E-10-03) ahşap, kâğıt ve mobilyaya genişletilir. Rus ve Belarus odunundan ayrışmış tedarik (E-16-02) aynı katmanda: AB'deki süre sonu gözden geçirmesi şimdi başladı ve Türk üreticilerin muafiyet başvuruları ancak soruşturma sürerken yapılabilir. Ambalajda PPWR uyumu (E-16-03), ticaret savunmasında karşılıklılık ve aşağı akışın gözetilmesi (E-16-04) ve levhada çevresel ürün beyanı (E-16.21-01) da bu katmanda.

İkinci katman uzun soluklu: odun ve lif hammaddesi güvenliği (E-16-05). Hızlı gelişen türlerle endüstriyel ağaçlandırma, Rus odununa bağımlılığı azaltmanın ve levha sanayisinin hammaddesini güvenceye almanın tek kalıcı yolu, ama sonuç vermesi on yılı aşıyor. Atık kâğıt toplamasının genişletilmesi daha hızlı sonuç verir ve kâğıt sanayisinin ithal hamura bağımlılığını azaltır. Türkiye'nin sıfır atık mevzuatı ve genişletilmiş üretici sorumluluğu bu konuda hazır bir çerçeve.

Üçüncü katman yalnızca S4'te ve çıkış gününde devreye giriyor. Bu parçada bu katmanda tek bir eylem var: ambalaj kâğıdında c2 düzeltmesi (E-17.12-01). Bu eylemin mantığı Faz 2A'nın c2 tasarımıyla aynı: yurt içi kabiliyeti olmayan girdilerde ek vergi yerli üreticiyi korumaz, yalnızca aşağı akışı ve tüketiciyi vergiler. Düzeltme, sıvı ambalaj kartonu ve kraftliner gibi kalitelerde İGV'yi bütün menşeler için sıfırlıyor. Bütün menşelere eşit uygulandığı için en çok kayrılan ülke kuralıyla uyumlu.

\begin{landscape}

\textbf{Tablo 4.2 --- Bölüm düzeyi ortak eylemler (16--18)}

\nopagebreak[4]

\begingroup\gbtablo\footnotesize

\begin{longtable}[]{@{}
  >{\raggedright\arraybackslash}p{(\columnwidth - 18\tabcolsep) * \real{0.0556}}
  >{\raggedright\arraybackslash}p{(\columnwidth - 18\tabcolsep) * \real{0.0967}}
  >{\raggedright\arraybackslash}p{(\columnwidth - 18\tabcolsep) * \real{0.1946}}
  >{\raggedright\arraybackslash}p{(\columnwidth - 18\tabcolsep) * \real{0.0990}}
  >{\raggedright\arraybackslash}p{(\columnwidth - 18\tabcolsep) * \real{0.1062}}
  >{\raggedright\arraybackslash}p{(\columnwidth - 18\tabcolsep) * \real{0.0845}}
  >{\raggedright\arraybackslash}p{(\columnwidth - 18\tabcolsep) * \real{0.0984}}
  >{\raggedright\arraybackslash}p{(\columnwidth - 18\tabcolsep) * \real{0.1039}}
  >{\raggedright\arraybackslash}p{(\columnwidth - 18\tabcolsep) * \real{0.0570}}
  >{\raggedright\arraybackslash}p{(\columnwidth - 18\tabcolsep) * \real{0.1040}}@{}}
\toprule\noalign{}
\begin{minipage}[b]{\linewidth}\raggedright
\textbf{Eylem}
\end{minipage} & \begin{minipage}[b]{\linewidth}\raggedright
\textbf{Araç}
\end{minipage} & \begin{minipage}[b]{\linewidth}\raggedright
\textbf{Parametre}
\end{minipage} & \begin{minipage}[b]{\linewidth}\raggedright
\textbf{Sorumlu}
\end{minipage} & \begin{minipage}[b]{\linewidth}\raggedright
\textbf{Hukuki araç}
\end{minipage} & \begin{minipage}[b]{\linewidth}\raggedright
\textbf{Bütçe}
\end{minipage} & \begin{minipage}[b]{\linewidth}\raggedright
\textbf{Başlangıç → bitiş}
\end{minipage} & \begin{minipage}[b]{\linewidth}\raggedright
\textbf{Performans koşulu}
\end{minipage} & \begin{minipage}[b]{\linewidth}\raggedright
\textbf{Sunset}
\end{minipage} & \begin{minipage}[b]{\linewidth}\raggedright
\textbf{KPI}
\end{minipage} \\
\midrule\noalign{}
\endhead
\bottomrule\noalign{}
\endlastfoot
E-16-01 & EUDR ahşap ve kâğıt veri hizmeti & OGM'nin amenajman planlarındaki bölme ve bölmecik sınırlarının, her dikili ve istif satışında alıcıya EUDR'ye uygun coğrafi konum ve yasallık verisi olarak verilmesi; ithal odun ve kâğıt hamurunda tedarikçiden coğrafi konum verisinin ithalat beyannamesine eklenmesi; ilk parçadaki EUDR uyum masasının (E-10-03) ahşap, kâğıt ve mobilyaya genişletilmesi & Tarım ve Orman Bakanlığı (OGM), Ticaret Bakanlığı (Gümrükler GM, İhracat GM), TİM & OGM genelgesi (satış şartnamesi); Ticaret Bakanlığı genelgesi & ≈ 0,1 mlr TL/yıl \etiketT{} & Hemen → T−12 (EUDR 30.12.2026'da uygulanmaya başlıyor) & OGM satışlarının EUDR verisiyle yapılan payı & Yok (kalıcı altyapı) & EUDR gerekçesiyle AB'de reddedilen Türk orman ürünü sevkiyatı sıfır; Türkiye'nin düşük risk sınıfında kalması \\
E-16-02 & Rus ve Belarus odunundan ayrışmış tedarik & Levha, kontrplak ve kereste ihracatçıları için firma düzeyinde odun menşe beyanı (ithalat kayıtlarıyla doğrulanan); AB'ye ihracat hatlarında Rus ve Belarus odunu kullanılmaması; huş kontrplağındaki süre sonu gözden geçirmesine (AB, 6.11.2026) katılım ve Türk üreticilerin muafiyet başvurularına teknik destek & Ticaret Bakanlığı (İthalat GM, Ticaret Savunması), TİM (orman ürünleri ihracatçı birlikleri) & Bakanlık genelgesi; AB Temel Anti-Damping Tüzüğü md. 11(2) ve 13(4) çerçevesinde başvuru & ≈ 0,02 mlr TL/yıl \etiketT{} & Hemen → T+24 & Menşe beyanı veren ihracatçıların AB'ye levha ve kontrplak ihracatındaki payı & T+24 & Türkiye'den gönderilen ürüne yeni anti-damping genişletmesi sıfır; gözden geçirmede Türk üreticiler için muafiyet \\
E-16-03 & Ambalajda PPWR uyumu & AB'ye ambalaj ve gıdayla temas eden kâğıt satan ihracatçılar için PFAS'sız kaplamaya geçiş, uygunluk beyanı ve teknik dosya şablonu; üye devletlerde genişletilmiş üretici sorumluluğu kaydı için rehber; Türkiye'nin ambalaj atığı mevzuatının PPWR ile uyumu (geri dönüştürülebilirlik ve etiketleme) & Ticaret Bakanlığı, Çevre, Şehircilik ve İklim Değişikliği Bakanlığı, TSE, TİM & Bakanlık rehberi; Ambalaj Atıklarının Kontrolü Yönetmeliği değişikliği & ≈ 0,05 mlr TL/yıl \etiketT{} & Hemen → T & Uygunluk dosyası hazır ihracatçı sayısı & T+12 & PPWR gerekçesiyle AB pazar gözetiminde Türk ambalajına ilişkin bildirim sıfır \\
E-16-04 & Ticaret savunmasında karşılıklılık ve aşağı akışın gözetilmesi & (i) S2 metninde ticaret savunması soruşturmasından önce zorunlu istişare ve anti-damping genişletmelerinde (md. 13) taraflar arası bilgi paylaşımı; (ii) Türkiye'nin kâğıt korunma önlemlerinde (testliner ve fluting, 2027'de sona eriyor; kaplamalı kâğıt soruşturması, Tebliğ 2026/3) ambalaj ve matbaacılık kullanıcılarının ve tüketicinin çıkarının kararda ayrıca tartılması & Ticaret Bakanlığı (İthalat GM) & S2 anlaşma metni (ticaret savunması bölümü); İthalatta Korunma Önlemleri Yönetmeliği uygulaması & İdari & Hemen → T & S2 metninde istişare maddesi & --- & Önlem kararlarında aşağı akış etkisinin yayımlanan soruşturma raporunda yer alması \\
E-16-05 & Odun ve lif hammaddesi güvenliği & Hızlı gelişen türlerle (kavak, okaliptüs, kızılçam) endüstriyel ağaçlandırma ve özel ağaçlandırma kredisi; belediye atık kâğıt toplamasının genişletilmesi (depozito ve genişletilmiş üretici sorumluluğu); kâğıt hamuru ve kerestede tedarik kaynaklarının çeşitlendirilmesi & Tarım ve Orman Bakanlığı (OGM), Çevre, Şehircilik ve İklim Değişikliği Bakanlığı, belediyeler & OGM ağaçlandırma programı; Sıfır Atık mevzuatı & ≈ 0,2 mlr TL/yıl (mevcut programlara ek) \etiketT{} & Hemen → T+60 & Yıllık endüstriyel ağaçlandırma alanı; atık kâğıt toplama oranı & T+60 & Odun ithalatında Rusya ve Belarus payı T+3'te ≤ \%10; atık kâğıt toplama oranı T+5'te +10 puan \\
\end{longtable}

\endgroup

\end{landscape}

\emph{Kaynak:} Bütçeler \etiketT{}, 2026 fiyatlarıyla; kurumsal sorumluluk mevcut görev dağılımına göre. Mali yükler ve cari denge katkıları kümülatif defterde (Bölüm 6) toplanıyor.

\begin{landscape}

\textbf{Tablo 4.3 --- Sınıf eylemleri}

\nopagebreak[4]

\begingroup\gbtablo\footnotesize

\begin{longtable}[]{@{}
  >{\raggedright\arraybackslash}p{(\columnwidth - 18\tabcolsep) * \real{0.0728}}
  >{\raggedright\arraybackslash}p{(\columnwidth - 18\tabcolsep) * \real{0.0847}}
  >{\raggedright\arraybackslash}p{(\columnwidth - 18\tabcolsep) * \real{0.1971}}
  >{\raggedright\arraybackslash}p{(\columnwidth - 18\tabcolsep) * \real{0.0845}}
  >{\raggedright\arraybackslash}p{(\columnwidth - 18\tabcolsep) * \real{0.1175}}
  >{\raggedright\arraybackslash}p{(\columnwidth - 18\tabcolsep) * \real{0.0894}}
  >{\raggedright\arraybackslash}p{(\columnwidth - 18\tabcolsep) * \real{0.0717}}
  >{\raggedright\arraybackslash}p{(\columnwidth - 18\tabcolsep) * \real{0.0911}}
  >{\raggedright\arraybackslash}p{(\columnwidth - 18\tabcolsep) * \real{0.0786}}
  >{\raggedright\arraybackslash}p{(\columnwidth - 18\tabcolsep) * \real{0.1126}}@{}}
\toprule\noalign{}
\begin{minipage}[b]{\linewidth}\raggedright
\textbf{Eylem}
\end{minipage} & \begin{minipage}[b]{\linewidth}\raggedright
\textbf{Araç}
\end{minipage} & \begin{minipage}[b]{\linewidth}\raggedright
\textbf{Parametre}
\end{minipage} & \begin{minipage}[b]{\linewidth}\raggedright
\textbf{Sorumlu}
\end{minipage} & \begin{minipage}[b]{\linewidth}\raggedright
\textbf{Hukuki araç}
\end{minipage} & \begin{minipage}[b]{\linewidth}\raggedright
\textbf{Bütçe}
\end{minipage} & \begin{minipage}[b]{\linewidth}\raggedright
\textbf{Başlangıç → bitiş}
\end{minipage} & \begin{minipage}[b]{\linewidth}\raggedright
\textbf{Performans koşulu}
\end{minipage} & \begin{minipage}[b]{\linewidth}\raggedright
\textbf{Sunset}
\end{minipage} & \begin{minipage}[b]{\linewidth}\raggedright
\textbf{KPI}
\end{minipage} \\
\midrule\noalign{}
\endhead
\bottomrule\noalign{}
\endlastfoot
E-16.21-01 & Çevresel ürün beyanı ve karbon verisi & MDF, yonga levha ve kontrplakta EN 15804 uyumlu çevresel ürün beyanı (EPD); yeni Yapı Ürünleri Tüzüğü'nün (2024/3110) çevresel beyan yükümlülüğüne hazırlık & Sanayi ve Teknoloji Bakanlığı, TSE, TİM & Bakanlık programı & ≈ 0,02 mlr TL/yıl \etiketT{} & Hemen → T+36 & EPD'si olan ürün sayısı & T+36 & AB'ye levha ihracatının EPD'li payı ≥ \%80 \\
E-17.12-01 & Ambalaj kâğıdında c2 düzeltmesi (yalnız S4) & Yurt içi üretimi olmayan ya da sınırlı ambalaj girdisi kalitelerinde (sıvı ambalaj kartonu 4811 59, ağartılmamış kraftliner 4804 11) İGV bütün menşeler için sıfırlanır; c2 satır listesi 8--12 haneli üretim verisiyle gözden geçirilir (VB-43) & Ticaret Bakanlığı, Hazine ve Maliye Bakanlığı & İthalat Rejimi Kararı eki (Cumhurbaşkanı Kararı) & ≈ 0,7 mlr TL/yıl (AB dışı menşelerin İGV kaybı) \etiketT{} & T → T+36 & Bu kalitelerde ithalat birim fiyatı & T+36 (yurt içi kapasite oluşursa gözden geçirilir) & Oluklu ambalaj ve sıvı gıda ambalajı üreticilerinin girdi maliyeti artışı ≤ \%1 \\
\end{longtable}

\endgroup

\end{landscape}

\emph{Kaynak:} Bütçeler \etiketT{}, 2026 fiyatlarıyla. Tarihler T'ye göre (T2028'de T = 1.1.2028, T2029'da 1.1.2029).

Paketin yıllık mali maliyeti S2 yolunda 0,39 milyar TL, S4 yolunda 1,09 milyar TL. Maliyetin çoğu hammadde güvenliği programının mevcut ağaçlandırma bütçesine eki ve EUDR veri hizmeti. S4'te buna c2 düzeltmesinin İGV kaybı ekleniyor. Paketin cari dengeye ölçülebilir bir katkısı yok. EUDR ve ticaret savunması önlemlerinin değeri, olası bir kaybı önlemekte. Model bu kaybı ölçmediği için deftere yazılmadı.

\hypertarget{nace_kartlar_16-18-43-muxfczakere-pozisyonu}{%
\subsection{4.3 Müzakere pozisyonu}\label{nace_kartlar_16-18-43-muxfczakere-pozisyonu}}

AB ile müzakerede bu sektörün tutumu çoğunlukla nötr; yalnızca ahşap levhada savunma. Tarife tarafında pazarlık konusu yok. Kâğıtta AB'nin tarifesi zaten sıfır; ahşap levhada S2 sıfır tarifeyi geri getiriyor. Menşe kuralları kolay. Bu parçanın S2'den beklentisi tarife dışı üç bölümde.

Birincisi ticaret savunması. S2 metni, taraflardan birinin diğerine karşı anti-damping, telafi edici vergi ya da dolanma soruşturması açmadan önce zorunlu istişare yapmasını ve dolanma soruşturmalarında ihracatçı tarafın idaresiyle bilgi paylaşılmasını içermeli (E-16-04). AB'nin son serbest ticaret anlaşmalarında ticaret savunması bölümleri bu tür usul hükümleri içeriyor \etiketD{}. Talep AB için düşük maliyetli ve Türkiye için huş kontrplağı tecrübesinin tekrarlanmaması açısından değerli. İkincisi düzenleyici işbirliği: OGM'nin orman verisinin EUDR açısından tanınması ve EUDR ile PPWR'nin uygulanmasında bir teknik işbirliği mekanizması. Üçüncüsü sınır kolaylaştırması: yetkilendirilmiş yükümlü statüsünün karşılıklı tanınması, kâğıt ve ambalaj gibi hacimli ve düşük değerli ürünlerde sınır maliyetinin en etkili ilacı.

AB'nin bu sektörde Türkiye'den isteyeceği şey, S2'de Türkiye'nin kâğıt ve karton pazarına sıfır tarifeyle ve korunma önlemi olmadan erişim. Finlandiya, İsveç ve Almanya'nın kâğıt üreticileri için Türkiye büyük bir pazar. Türkiye'nin oluklu mukavva kâğıdındaki korunma önlemi bu pazara erişimi zaten kısıtladı. Bu, müzakerede bir pazarlık konusu: önlemin Temmuz 2027'de uzatılmaması ya da S2'de AB'ye karşı ikili bir korunma mekanizmasıyla değiştirilmesi, AB'nin ticaret savunması usulleri konusundaki tavizine karşılık sunulabilir \etiketH{}. Bu taviz, Türkiye'nin aşağı akış sanayileri için zaten yararlı.

Üçüncü ülkelerle önceliklerin başında Körfez ve Irak geliyor. Irak levhanın ve hijyen ürünlerinin ilk pazarı. Körfez İşbirliği Konseyi ile STA müzakereleri (Faz 2C) levha, kâğıt ürünleri ve hijyen ürünlerinde gümrük vergisini kaldırarak bu pazarı derinleştirebilir. Birleşik Krallık, kâğıt ve oluklu ambalajın en büyük AB dışı pazarlarından biri ve Türkiye--Birleşik Krallık STA'sı bütün senaryolarda sürüyor. Odun ve hamur tedarikinde ise amaç çeşitlendirme: Brezilya, Şili ve Uruguay'dan hamur, Gürcistan, Romanya ve Kuzey Amerika'dan odun.

\hypertarget{nace_kartlar_16-18-44-erken-uyarux131-ve-tetikler}{%
\subsection{4.4 Erken uyarı ve tetikler}\label{nace_kartlar_16-18-44-erken-uyarux131-ve-tetikler}}

Bu parçanın tetikleri ağırlıkla senaryodan bağımsız. En kritik iki gösterge AB'nin ticaret savunması ve EUDR uygulaması. AB'nin Türk levha ya da kontrplağına yönelik yeni bir soruşturma açması, Rus odunundan ayrışma önlemini (E-16-02) hızlandırıyor ve S2 müzakeresinde istişare maddesini öne alıyor. EUDR kontrollerinde Türk orman ürününe ilişkin bir ret ya da yüksek risk bildirimi, veri hizmetinin genişletilmesini ve Komisyon'la teknik görüşmeyi başlatıyor. S4'e özgü tek tetik AB'den kâğıt ithalatının birim fiyatı: yıllık \%15'lik bir artış, c2 düzeltmesinin kapsamının genişletilmesini gerektiriyor.

\Needspace{7\baselineskip}

\textbf{Tablo 4.4 --- Erken uyarı göstergeleri ve tetikler}

\nopagebreak[4]

\begingroup\gbtablo\small

\begin{longtable}[]{@{}
  >{\raggedright\arraybackslash}p{(\columnwidth - 6\tabcolsep) * \real{0.0679}}
  >{\raggedright\arraybackslash}p{(\columnwidth - 6\tabcolsep) * \real{0.4173}}
  >{\raggedright\arraybackslash}p{(\columnwidth - 6\tabcolsep) * \real{0.1499}}
  >{\raggedright\arraybackslash}p{(\columnwidth - 6\tabcolsep) * \real{0.3649}}@{}}
\toprule\noalign{}
\begin{minipage}[b]{\linewidth}\raggedright
\textbf{Sınıf}
\end{minipage} & \begin{minipage}[b]{\linewidth}\raggedright
\textbf{Gösterge}
\end{minipage} & \begin{minipage}[b]{\linewidth}\raggedright
\textbf{Eşik}
\end{minipage} & \begin{minipage}[b]{\linewidth}\raggedright
\textbf{Tetiklenecek önlem}
\end{minipage} \\
\midrule\noalign{}
\endhead
\bottomrule\noalign{}
\endlastfoot
16.21 & AB'nin Türk levha ya da kontrplağına yönelik yeni ticaret savunması soruşturması & Soruşturma açılışı & E-16-02 hızlandırılır; S2 müzakeresinde istişare maddesi öne alınır \\
16.21 & EUDR kontrollerinde Türk orman ürününe ilişkin ret ya da yüksek risk bildirimi & İlk bildirim & E-16-01 veri hizmetinin genişletilmesi; Komisyonla teknik görüşme \\
17.12 & AB'den kâğıt ve karton ithalatının birim fiyatı (aylık) & Yıllık +\%15 (S4) & E-17.12-01 kapsamının genişletilmesi \\
17.21 & AB'ye oluklu ambalaj ihracatı (aylık) & Yıllık −\%15 & E-16-03 hızlandırılır; Balkan ve Kafkasya pazarları \\
17.22 & Kâğıt hamuru ithal birim fiyatı (aylık) & 3 ayda +\%25 & İthalatçı stok bildirimi; tedarik çeşitlendirme (E-16-05) \\
\end{longtable}

\endgroup

\emph{Kaynak:} Göstergeler aylık izlenir: TÜİK dış ticaret ve TÜFE alt endeksleri, Eurostat Comext, AB Resmî Gazetesi.

\hypertarget{nace_kartlar_16-18-45-kazananlar-kaybedenler-ve-buxf6lgeler}{%
\subsection{4.5 Kazananlar, kaybedenler ve bölgeler}\label{nace_kartlar_16-18-45-kazananlar-kaybedenler-ve-buxf6lgeler}}

Çıkışın bu parçadaki dağılım etkileri küçük, ama yönleri belli. S4'te kazananlar iç pazara satan kâğıt ve karton üreticileri ve c2 tarifesiyle AB ürününe karşı iç pazarda pay kazanan oluklu kutu ve kâğıt ürünleri üreticileri. Kaybedenler AB'ye levha satan büyük entegre üreticiler ve kâğıdı girdi olarak kullanan matbaalar ve ambalajcılar. Matbaacılar kaybını model göremiyor. Tüketicinin payı küçük: S4'te TÜFE etkisi 0,9 baz puan.

\Needspace{7\baselineskip}

\textbf{Tablo 4.5 --- Bölgeler: 16--18 istihdamının en yoğun olduğu 10 İBBS-2 bölgesi (T+5, P50)}

\nopagebreak[4]

\begingroup\gbtablo\footnotesize

\begin{longtable}[]{@{}
  >{\raggedright\arraybackslash}p{(\columnwidth - 12\tabcolsep) * \real{0.2024}}
  >{\raggedleft\arraybackslash}p{(\columnwidth - 12\tabcolsep) * \real{0.1467}}
  >{\raggedleft\arraybackslash}p{(\columnwidth - 12\tabcolsep) * \real{0.1662}}
  >{\raggedleft\arraybackslash}p{(\columnwidth - 12\tabcolsep) * \real{0.0921}}
  >{\raggedleft\arraybackslash}p{(\columnwidth - 12\tabcolsep) * \real{0.0921}}
  >{\raggedleft\arraybackslash}p{(\columnwidth - 12\tabcolsep) * \real{0.1141}}
  >{\raggedleft\arraybackslash}p{(\columnwidth - 12\tabcolsep) * \real{0.1864}}@{}}
\toprule\noalign{}
\begin{minipage}[b]{\linewidth}\raggedright
\textbf{Bölge (İBBS-2)}
\end{minipage} & \begin{minipage}[b]{\linewidth}\raggedleft
\textbf{16--18 çalışanı (bin)}
\end{minipage} & \begin{minipage}[b]{\linewidth}\raggedleft
\textbf{Bölge imalatındaki payı}
\end{minipage} & \begin{minipage}[b]{\linewidth}\raggedleft
\textbf{S2 (kişi)}
\end{minipage} & \begin{minipage}[b]{\linewidth}\raggedleft
\textbf{S4 (kişi)}
\end{minipage} & \begin{minipage}[b]{\linewidth}\raggedleft
\textbf{S4 ÷ çalışan}
\end{minipage} & \begin{minipage}[b]{\linewidth}\raggedleft
\textbf{En çok etkilenen sınıflar (S4)}
\end{minipage} \\
\midrule\noalign{}
\endhead
\bottomrule\noalign{}
\endlastfoot
İstanbul & 60,7 & \%6 & +61 & +158 & \%0,3 & 16.21 (−12), 16.23 (−4) \\
Kocaeli, Sakarya, Düzce, Bolu, Yalova & 26,1 & \%5 & +5 & +24 & \%0,1 & 16.21 (−13), 16.23 (−4) \\
İzmir & 16,9 & \%6 & +21 & +55 & \%0,3 & 16.23 (−2), 16.21 (−1) \\
Bursa, Eskişehir, Bilecik & 14,5 & \%3 & +3 & +13 & \%0,1 & 16.21 (−10), 16.23 (−3) \\
Gaziantep, Adıyaman, Kilis & 12,4 & \%6 & +5 & +19 & \%0,1 & 16.23 (−3), 17.22 (0) \\
Adana, Mersin & 11,6 & \%6 & +16 & +41 & \%0,4 & 16.23 (−4), 16.21 (−1) \\
Tekirdağ, Edirne, Kırklareli & 10,9 & \%5 & +15 & +38 & \%0,3 & 16.23 (−2), 16.21 (−1) \\
Ankara & 10,9 & \%3 & −1 & +3 & \%0,0 & 16.21 (−2), 16.23 (−2) \\
Aydın, Denizli, Muğla & 9,4 & \%6 & +2 & +10 & \%0,1 & 16.23 (−5), 16.21 (0) \\
Manisa, Afyon, Kütahya, Uşak & 9,3 & \%4 & +3 & +12 & \%0,1 & 16.23 (−4), 16.21 (−2) \\
\textbf{16--18 toplamı (bütün bölgeler)} & \textbf{245,5} & \textbf{\%5} & \textbf{+170} & \textbf{+476} & \textbf{\%0,2} & --- \\
\end{longtable}

\endgroup

\emph{Kaynak:} \etiketT{} f3b (YSHİ 2024 sınıf çalışanı, SGK 2025 bölüm × il sigortalısıyla RAS) ve f3e. Bölge değerleri sınıf P50'lerinin bölge paylarıyla dağıtılmasından; toplam, sınıf P50'lerinin toplamıdır (grup yüzdeliği değil).

Bölgesel tablo, bu parçanın istihdamının büyük şehirlerde toplandığını gösteriyor. İstanbul tek başına 60.669 kişiyle ilk sırada; bunun büyük kısmı matbaacılık ve ambalaj. Kocaeli--Sakarya--Düzce--Bolu--Yalova bölgesi (26.133 kişi) levha ve kâğıdın, İzmir (16.883) kâğıt ve ambalajın merkezi. Bu bölgelerin sanayisi çeşitli ve parçanın bölge imalatındaki payı \%5--6 dolayında. Ağaç ve kâğıt sanayisinin görece en ağır olduğu bölge Kastamonu--Çankırı--Sinop (TR82): burada üç bölüm bölge imalat istihdamının \%13 kadarını oluşturuyor. Ama bu bölgede de senaryoların istihdam etkisi birkaç kişiyle sınırlı (−7 kişi). Bu parçada bölgesel bir işgücü önlemi gerekmiyor.

\hypertarget{nace_kartlar_16-18-5-sux131nux131f-incelemeleri}{%
\section{5. Sınıf incelemeleri}\label{nace_kartlar_16-18-5-sux131nux131f-incelemeleri}}

Bu bölüm dokuz sınıfı tek tek inceliyor. Etkisi sınırlı dokuz sınıf Bölüm 5.10'daki tabloda toplanıyor. Her incelemenin başındaki künye sınıf kartının özet alanlarını veriyor; senaryo tablosu çıkışın etkisinin büyük olduğu ya da yön değiştirdiği dört sınıf için burada, diğerleri için Ek C'de.

\hypertarget{nace_kartlar_16-18-51-1621-ahux15fap-levha-ve-kontrplak}{%
\subsection{5.1 16.21 Ahşap levha ve kontrplak}\label{nace_kartlar_16-18-51-1621-ahux15fap-levha-ve-kontrplak}}

\emph{GB statüsü: GB · Arketip: B (küresel rekabetçi, AB bağımlılığı düşük) / G (SKDM ve Yeşil Mutabakat maruz) · Duruş: Geliştir · AB ile müzakere: savunma · Güven: Orta · Veri sayfası: Ek C.2}

Ahşap levha bu parçanın en büyük ihracatçısı ve Türkiye'nin dünya pazarındaki en güçlü ürünlerinden biri. Sınıfta 22.174 kişi çalışıyor; üretim 5,0 milyar dolar ve ihracat 1.350 milyon dolar. 2024'te lif levha ihracatında Türkiye, Almanya ve Çin'den sonra dünyada 3. sırada; yonga levhada payı \%2,5. Sınıf az sayıda büyük, entegre grubun elinde: istihdamın \%69 kadarı 250'den fazla kişi çalıştıran işletmelerde. Bu gruplardan Kastamonu Entegre, kendi açıklamasına göre MDF ve yonga levhada Avrupa'nın üçüncü, dünyanın dördüncü büyük üreticisi \etiketD{}. Verimlilik AB'deki muadilinin 0,55 katı; bu parçada AB'ye en yakın sınıflardan biri.

AB sınıfın ihracatının \%16 kadarını alıyor ve AB'deki ilk pazarlar komşular: Romanya, Bulgaristan ve Yunanistan. AB dışında Irak (\%14) ve Fas (\%12) öne çıkıyor. S4'te AB tarifesi ortalama \%6,9; yapı ürünü olarak satılan levhanın bir kısmı için AB'de tanınan bir kuruluşun deneyi de gerekiyor. AB'ye ihracat 56 milyon dolar azalıyor, ama istihdam kaybı 68 kişiyle sınırlı. Kaybedilen AB satışının bir kısmı başka pazarlara kayıyor, bir kısmını da c2 tarifesiyle pahalanan AB levhasının iç pazarda bıraktığı boşluk karşılıyor. S2'de kayıp 23 kişi.

Sınıfın asıl riski Bölüm 3.2'de anlatılan Rus odunu. Sınıfın ithalatının \%39 kadarı Rusya'dan ve bunun ağırlığı kontrplak ve kaplama. AB, Rus huş kontrplağına uyguladığı anti-damping vergisini 2024'te Türkiye'den gönderilen huş kontrplağına genişletti ve bu genişletme, 6 Kasım 2026'da açılan süre sonu gözden geçirmesi boyunca yürürlükte \etiketV{}. Türkiye'nin kontrplak ihracatı küçük (69 milyon dolar), ama aynı mantık lif ve yonga levhaya da uygulanabilir. Sınıfın AB pazarındaki geleceği, ürünün Rus odunu içermediğini belgeleyebilmesine ve EUDR verisine bağlı.

Önerilen duruş ``Geliştir''. Strateji, AB dışı pazarlarda büyümek ve AB pazarında güvenilirliği belgelemek. Ortak eylemler bu sınıf için tasarlandı: EUDR veri hizmeti (E-16-01) ve Rus ve Belarus odunundan ayrışmış tedarik (E-16-02). Sınıfa özgü eylem, AB'nin yeni Yapı Ürünleri Tüzüğü'nün çevresel beyan yükümlülüğüne hazırlık olarak MDF, yonga levha ve kontrplakta çevresel ürün beyanı (E-16.21-01). Girdi tarafında yonga, tomruk ve kaplamanın tarifesi zaten sıfıra yakın; S4'te bunlar sıfırda kalmalı (P1). Değerlendirmenin güveni orta, çünkü Rus odununun sınıfın tedarikindeki gerçek payı bilinmiyor (VB-40).

\Needspace{7\baselineskip}

\textbf{Tablo 5.1 --- 16.21 Ahşap levha ve kontrplak: senaryo etkileri (T+5'te yıllık; P50, parantez içinde P10 / P90)}

\nopagebreak[4]

\begingroup\gbtablo\small

\begin{longtable}[]{@{}lrrr@{}}
\toprule\noalign{}
\textbf{Değişken} & \textbf{S2} & \textbf{S3} & \textbf{S4} \\
\midrule\noalign{}
\endhead
\bottomrule\noalign{}
\endlastfoot
AB'ye ihracat (\$ mn) & −22 (−51 / +4) & −31 (−61 / −5) & −56 (−93 / −24) \\
AB'ye ihracat (\%) & −10,1 (−23,7 / +2,0) & −14,3 (−28,3 / −2,2) & −26,1 (−42,9 / −11,2) \\
Toplam ihracat (\$ mn) & −16 (−37 / +3) & −22 (−44 / −3) & −45 (−74 / −19) \\
Nihai mal ithalatı (\$ mn) & −8 (−26 / +7) & −11 (−28 / +4) & −22 (−41 / −5) \\
Üretim (\$ mn) & −7 (−15 / −3) & −11 (−19 / −6) & −22 (−35 / −14) \\
İstihdam (kişi) & −23 (−45 / −8) & −33 (−58 / −18) & −68 (−113 / −41) \\
İstihdam (\%) & −0,1 (−0,2 / 0,0) & −0,1 (−0,3 / −0,1) & −0,3 (−0,5 / −0,2) \\
Cari denge (\$ mn) & −6 (−11 / −2) & −8 (−14 / −5) & −17 (−28 / −11) \\
TÜFE (bp) & +0,01 (+0,01 / +0,01) & +0,02 (+0,02 / +0,02) & +0,07 (+0,06 / +0,08) \\
\end{longtable}

\endgroup

\emph{Kaynak:} \etiketT{} f3e; toplamlar Faz 2B genel denge dağılımına kalibre. Yüzdeler 2024 düzeyine göre.

\hypertarget{nace_kartlar_16-18-52-1610-kereste}{%
\subsection{5.2 16.10 Kereste}\label{nace_kartlar_16-18-52-1610-kereste}}

\emph{GB statüsü: GB · Arketip: --- (belirgin arketip yok) · Duruş: Koru · AB ile müzakere: nötr · Güven: Orta · Veri sayfası: Ek C.1}

Kereste sınıfında 15.956 kişi ve 3.425 girişim var. Çoğu, OGM'den tomruk alan küçük bıçkıhaneler; istihdamın \%18 kadarı Kocaeli--Sakarya--Düzce--Bolu--Yalova bölgesinde, kalanı Karadeniz'e dağılmış durumda. Verimlilik AB'deki muadilinin 0,31 katı. Türkiye kerestede net ithalatçı. İhracat 101 milyon dolar ve ilk pazarlar İran ile Irak; ithalat 327 milyon dolar ve ilk kaynaklar Ukrayna (\%24), ABD (\%24) ve Rusya (\%13).

Çıkışın sınıfa etkisi ihmal edilebilir: AB'ye ihracat yalnızca 24 milyon dolar ve istihdam etkisi her iki senaryoda da sıfıra yakın. Kurallar sınıfı EUDR nedeniyle ``Yeşil Mutabakat'a maruz'' olarak işaretliyordu; bu işaret kaldırıldı, çünkü AB'ye ihracat küçük ve EUDR'nin coğrafi konum sorunu OGM verisiyle çözülebiliyor. Önerilen duruş ``Koru''. Sınıfın gerçek sorunu hammadde: yerli tomruk arzı sınırlı ve ithal kerestenin önemli bir kısmı Rusya'dan. Bu sorun ortak eylemlerle, özellikle hammadde güvenliği programıyla (E-16-05) ele alınıyor. Değerlendirmenin güveni orta.

\hypertarget{nace_kartlar_16-18-53-1623-doux11frama-ve-ahux15fap-yapux131}{%
\subsection{5.3 16.23 Doğrama ve ahşap yapı}\label{nace_kartlar_16-18-53-1623-doux11frama-ve-ahux15fap-yapux131}}

\emph{GB statüsü: GB · Arketip: --- (belirgin arketip yok) · Duruş: Koru · AB ile müzakere: nötr · Güven: Orta · Veri sayfası: Ek C.4}

Doğrama bu parçanın en büyük işveren sınıflarından biri: 30.303 kişi ve 11.871 girişim. Ahşap kapı, pencere, merdiven, parke dışındaki ahşap yapı elemanları ve prefabrik ahşap yapılar bu sınıfta. Üretim büyük ölçüde iç pazara, inşaata dönük ve atölye düzeyinde: kişi başı katma değer AB'deki muadilinin 0,17 katı, parçanın en düşüğü. İhracat 227 milyon dolar; AB'nin payı \%22 ve ilk AB pazarları Bulgaristan, Romanya, Almanya ve Fransa. İthalatın ilk kaynağı Avusturya (\%34), yani yapı sistemleri ve lamine ahşap.

S4'te istihdam kaybı 61 kişi, S2'de 37. Kayıp küçük, ama oransal olarak parçanın en büyüklerinden. Bunun bir nedeni, sınıfın ürünlerinin AB'de yapı ürünü olması. Pencere ve dış kapının AB'de satılabilmesi için CE işareti ve AB'nin tanıdığı bir laboratuvarın deney raporu gerekiyor. S4'te Türk kuruluşlarının raporları tanınmıyor ve bu, küçük ihracatçılar için tarifeden daha ağır bir engel olabilir \etiketH{}. S2'de uygunluk değerlendirmesinin karşılıklı tanınması (ACAA) bu maliyeti büyük ölçüde kaldırıyor.

Önerilen duruş ``Koru''. Sınıfın çıkıştan bağımsız ve daha büyük sorunu verimlilik. Bu sorun ticaret politikasıyla değil, KOBİ ve sanayi politikasıyla çözülebilir: makineleşme, standart ürün ve yapı sistemleri, ahşap yapının yönetmeliklerde yer alması. Deprem sonrası yeniden yapılanma ve ahşap yapının Avrupa'da yaygınlaşması sınıfa bir büyüme alanı açıyor \etiketH{}. Değerlendirmenin güveni orta.

\hypertarget{nace_kartlar_16-18-54-1711-kuxe2ux11fux131t-hamuru}{%
\subsection{5.4 17.11 Kâğıt hamuru}\label{nace_kartlar_16-18-54-1711-kuxe2ux11fux131t-hamuru}}

\emph{GB statüsü: GB · Arketip: --- (belirgin arketip yok) · Duruş: Koru · AB ile müzakere: nötr · Güven: Düşük · Veri sayfası: Ek C.7}

Türkiye'de kâğıt hamuru üretimi yok denecek kadar az. Sınıfın üretim ve istihdam verisi gizli; bu, sınıfta çok az sayıda işletme olduğunu gösteriyor (VB-44). Türkiye'nin hamur ithalatı ise büyük: 1.068 milyon dolarlık kimyasal hamur ve 289 milyon dolarlık atık kâğıt. Hamurun ilk kaynakları Brezilya (\%25), ABD (\%17), Finlandiya (\%16), Uruguay (\%11) ve Portekiz. Sınıfın 117 milyon dolarlık ihracatının \%77 kadarı İran'a gidiyor; bu, büyük olasılıkla ithal hamurun yeniden ihracatı \etiketH{}.

Hamurda Türkiye'nin MFN tarifesi sıfır ve c2 bunu değiştirmiyor. S4'te AB'den hamur ithalatı sınır sürtünmesiyle bir miktar azalıyor, boşluğu başka menşeler dolduruyor. Önceki parçalarda model bu boşluğun bir kısmının yerli üretimle dolacağını varsayıyordu ve bu, S2'de 129, S4'te 98 kişilik bir istihdam artışı üretiyordu. Bu artış anlamsızdı, çünkü ikame edecek bir yerli üretim yok. Bu parçada sınıfın iç ikame tavanı sıfıra çekildi ve artış kalktı (V-44). Sınıfın senaryo etkisi artık sıfıra yakın.

Önerilen duruş ``Koru''. Hamur, hijyen kâğıtlarının ve kâğıdın tamamen ithal hammaddesi ve bu yüzden kritik hammadde boyutunda sınıfın stratejik skoru 31. Ama yerli bir hamur sanayisi kurmanın odun arzı açısından gerçekçi bir yolu yok: Türkiye'nin endüstriyel odunu levha sanayisine bile yetmiyor. Strateji tedarik çeşitlendirme, atık kâğıt toplamasının genişletilmesi ve uzun vadede hızlı gelişen türlerle ağaçlandırma (E-16-05). Değerlendirmenin güveni düşük, çünkü sınıfın yurt içi verisi yok.

\hypertarget{nace_kartlar_16-18-55-1712-kuxe2ux11fux131t-ve-karton}{%
\subsection{5.5 17.12 Kâğıt ve karton}\label{nace_kartlar_16-18-55-1712-kuxe2ux11fux131t-ve-karton}}

\emph{GB statüsü: GB · Arketip: C (ithalata bağımlı, yerlileştirme potansiyeli) / G (SKDM ve Yeşil Mutabakat maruz) · Duruş: Geliştir · AB ile müzakere: nötr · Güven: Orta · Veri sayfası: Ek C.8}

Kâğıt ve karton sınıfı bu parçanın kilit noktası. Sınıfta 13.598 kişi çalışıyor ve üretim 3,1 milyar dolar. Sınıfın iki yüzü var. Bir yüzü güçlü bir ihracatçı: Türkiye hijyen kâğıdı ana bobininde dünyanın 2. büyük ihracatçısı ve sınıfın ihracatı 1.372 milyon dolar. AB'deki ilk pazarlar Yunanistan, Bulgaristan ve Romanya; AB dışında Birleşik Krallık (\%20) ve Irak. Diğer yüzü büyük bir ithalatçı: ithalat 2.526 milyon dolar ve \%63 kadarı AB'den. Türkiye geri dönüştürülmüş liften yapılan oluklu mukavva kâğıdında ve hijyen kâğıdında güçlü; kaplamalı ağaçsız kâğıt, kraftliner, sıvı ambalaj kartonu ve kaplamalı kartonda ise ithalata bağımlı.

Çıkışın sınıfa etkisi ithalat tarafından geliyor ve olumlu. S4'te c2 tarifesi AB'den gelen kâğıt ve kartonu ortalama \%3,5 pahalandırıyor; AB'den ithalat 239 milyon dolar azalıyor ve sınıfın istihdam etkisi +409 kişi. Cari dengeye katkı +98 milyon dolar, TÜFE etkisi 0,48 baz puan. S2'de de AB'li ihracatçının menşe maliyeti nedeniyle küçük bir ikame var (+199 kişi). Bu kazanç iki nedenle abartılı (Bölüm 3.4). Birincisi, aşağı akıştaki matbaa ve ambalajcıların maliyet artışı modelde yok. İkincisi, c2'nin HS6 düzeyindeki kuralı, Türkiye'de üretilmeyen bazı kalitelere de İGV koyuyor ve bu kalitelerde ithalatın yerini yerli üretim alamaz.

Sınıf, Türkiye'nin ticaret politikasının bugün de en aktif olduğu alanlardan biri. Oluklu mukavva kâğıdında 2024'ten beri AB menşelisini de kapsayan bir korunma önlemi var; kaplamalı kâğıtta yeni bir soruşturma sürüyor \etiketV{}. Kurallar sınıfı ``Yeşil Mutabakat'a maruz'' (G) olarak işaretliyordu; bu işaret ikincil arketipe indirildi, birincil arketip ``ithalata bağımlı, yerlileştirme potansiyeli olan'' (C) yapıldı. Önerilen duruş ``Geliştir'', ama korumayla değil yatırımla. Türkiye'nin rekabet üstünlüğü geri dönüştürülmüş lifte. Kaplamalı karton ve testliner gibi bu lifle üretilebilen kalitelerde kapasite yatırımı, yatırım teşvik sisteminin araçlarıyla desteklenebilir. Lif arzı için de atık kâğıt toplamasının genişletilmesi gerekiyor (E-16-05). Korunma önlemlerinde aşağı akış kullanıcılarının çıkarı ayrıca tartılmalı (E-16-04). S4'te ise c2 düzeltmesi (E-17.12-01), yerli üretimi olmayan ambalaj kalitelerinde İGV'yi kaldırıyor. Değerlendirmenin güveni orta.

\Needspace{7\baselineskip}

\textbf{Tablo 5.2 --- 17.12 Kâğıt ve karton: senaryo etkileri (T+5'te yıllık; P50, parantez içinde P10 / P90)}

\nopagebreak[4]

\begingroup\gbtablo\small

\begin{longtable}[]{@{}lrrr@{}}
\toprule\noalign{}
\textbf{Değişken} & \textbf{S2} & \textbf{S3} & \textbf{S4} \\
\midrule\noalign{}
\endhead
\bottomrule\noalign{}
\endlastfoot
AB'ye ihracat (\$ mn) & −16 (−39 / +7) & −28 (−57 / −5) & −26 (−47 / −11) \\
AB'ye ihracat (\%) & −4,8 (−11,5 / +2,0) & −8,3 (−16,7 / −1,3) & −7,6 (−13,8 / −3,2) \\
Toplam ihracat (\$ mn) & −12 (−29 / +5) & −20 (−41 / −3) & −21 (−38 / −9) \\
Nihai mal ithalatı (\$ mn) & −79 (−254 / +75) & −112 (−282 / +45) & −155 (−281 / −34) \\
Üretim (\$ mn) & +66 (−70 / +225) & +92 (−48 / +241) & +135 (+23 / +245) \\
İstihdam (kişi) & +199 (−211 / +694) & +277 (−148 / +745) & +409 (+72 / +753) \\
İstihdam (\%) & +1,5 (−1,6 / +5,1) & +2,0 (−1,1 / +5,5) & +3,0 (+0,5 / +5,5) \\
Cari denge (\$ mn) & +48 (−51 / +163) & +67 (−35 / +175) & +98 (+17 / +177) \\
TÜFE (bp) & +0,13 (+0,11 / +0,16) & +0,26 (+0,21 / +0,32) & +0,48 (+0,41 / +0,56) \\
\end{longtable}

\endgroup

\emph{Kaynak:} \etiketT{} f3e; toplamlar Faz 2B genel denge dağılımına kalibre. Yüzdeler 2024 düzeyine göre.

\hypertarget{nace_kartlar_16-18-56-1721-oluklu-ambalaj}{%
\subsection{5.6 17.21 Oluklu ambalaj}\label{nace_kartlar_16-18-56-1721-oluklu-ambalaj}}

\emph{GB statüsü: GB · Arketip: A (AB'ye entegre ihracat şampiyonu) / G (SKDM ve Yeşil Mutabakat maruz) · Duruş: Koru · AB ile müzakere: nötr · Güven: Orta · Veri sayfası: Ek C.9}

Oluklu mukavva ve kâğıt ambalaj, bu parçanın en büyük işveren sınıfı: 40.967 kişi ve 2.426 girişim. Sınıf, Türkiye'nin bütün ihracatçı sanayilerinin ürününü taşıyan ambalajı üretiyor; gıda, içecek ve ilaç tedarik zincirinin vazgeçilmez bir halkası. Bu yüzden stratejik skoru gıda güvenliği boyutundan geliyor (SS 23). Sınıfın kendi ihracatı da büyük: 834 milyon dolar ve oluklu kutuda Türkiye dünyanın 7. büyük ihracatçısı. AB'nin payı \%39; ilk AB pazarları Almanya (\%19) ve Fransa (\%18), AB dışında Birleşik Krallık ve ABD.

AB'nin oluklu kutudaki MFN tarifesi sıfır. Bu yüzden S4'te AB'ye ihracat yalnızca sınır sürtünmesiyle azalıyor (16 milyon dolar). Buna karşılık c2, AB'den gelen kutuyu \%25 pahalandırıyor ve iç pazarda yerli üreticiye pay kazandırıyor; sınıfın S4'teki net istihdam etkisi +20 kişi. S2'de etki sıfıra yakın (+1 kişi): sınır maliyetiyle kaybedilen AB satışını, AB'li rakiplerin menşe maliyetiyle kazanılan iç pazar payı karşılıyor. Model, sınıfın en önemli S4 etkisini görmüyor: girdi maliyeti. Sınıfın ana girdileri kraftliner, testliner ve fluting; bunların bir kısmı AB'den geliyor ve testliner ile fluting 2024'ten beri korunma önlemine tabi. S4'te c2, AB menşeli kraftlinere de \%10 İGV ekliyor \etiketT{}. İhracatlık meyve ve sebze kutusu gibi yüksek dayanımlı ambalaj bu girdiye bağımlı.

Kurallar sınıfı birincil olarak ``Yeşil Mutabakat'a maruz'' (G), ikincil olarak ``AB'ye entegre ihracat şampiyonu'' (A) diye işaretliyordu. Uzman düzeltmesiyle A birincil, G ikincil arketip yapıldı ve duruş ``Koru'' olarak belirlendi. Sınıfın AB'deki konumu tarifeye değil, yakınlığa ve hıza dayanıyor; S4 bu konumu ciddi biçimde sarsmıyor. Asıl gündem AB'nin yeni ambalaj tüzüğü (PPWR): uygunluk dosyası, 2028'den itibaren etiketleme ve 2030'dan itibaren geri dönüştürülebilirlik şartları. Türk oluklu sanayisinin atık kâğıt temeli bu şartlar karşısında bir üstünlük. Ormansızlaşma tüzüğü ise tamamen geri dönüştürülmüş liften yapılan kutuyu ve yalnızca başka bir ürünü taşıyan ambalajı kapsamıyor \etiketV{}. Önerilen eylemler ortak paketin içinde: PPWR uyumu (E-16-03) ve girdide c2 düzeltmesi (E-17.12-01). Değerlendirmenin güveni orta.

\Needspace{7\baselineskip}

\textbf{Tablo 5.3 --- 17.21 Oluklu ambalaj: senaryo etkileri (T+5'te yıllık; P50, parantez içinde P10 / P90)}

\nopagebreak[4]

\begingroup\gbtablo\small

\begin{longtable}[]{@{}lrrr@{}}
\toprule\noalign{}
\textbf{Değişken} & \textbf{S2} & \textbf{S3} & \textbf{S4} \\
\midrule\noalign{}
\endhead
\bottomrule\noalign{}
\endlastfoot
AB'ye ihracat (\$ mn) & −10 (−24 / +7) & −17 (−36 / −3) & −16 (−29 / −7) \\
AB'ye ihracat (\%) & −3,0 (−7,3 / +2,0) & −5,2 (−10,9 / −0,9) & −4,9 (−8,9 / −2,1) \\
Toplam ihracat (\$ mn) & −7 (−18 / +4) & −12 (−25 / −2) & −13 (−24 / −5) \\
Nihai mal ithalatı (\$ mn) & −7 (−24 / +5) & −10 (−25 / +3) & −19 (−34 / −4) \\
Üretim (\$ mn) & 0 (−4 / +9) & −3 (−7 / +5) & +4 (−3 / +15) \\
İstihdam (kişi) & +1 (−19 / +45) & −16 (−36 / +22) & +20 (−15 / +72) \\
İstihdam (\%) & 0,0 (0,0 / +0,1) & 0,0 (−0,1 / +0,1) & 0,0 (0,0 / +0,2) \\
Cari denge (\$ mn) & 0 (−3 / +7) & −2 (−5 / +3) & +3 (−2 / +11) \\
TÜFE (bp) & +0,01 (+0,01 / +0,02) & +0,03 (+0,02 / +0,03) & +0,12 (+0,10 / +0,13) \\
\end{longtable}

\endgroup

\emph{Kaynak:} \etiketT{} f3e; toplamlar Faz 2B genel denge dağılımına kalibre. Yüzdeler 2024 düzeyine göre.

\hypertarget{nace_kartlar_16-18-57-1722-hijyen-kuxe2ux11fux131tlarux131-ve-bezler}{%
\subsection{5.7 17.22 Hijyen kâğıtları ve bezler}\label{nace_kartlar_16-18-57-1722-hijyen-kuxe2ux11fux131tlarux131-ve-bezler}}

\emph{GB statüsü: GB · Arketip: B (küresel rekabetçi, AB bağımlılığı düşük) / G (SKDM ve Yeşil Mutabakat maruz) · Duruş: Geliştir · AB ile müzakere: nötr · Güven: Orta · Veri sayfası: Ek C.10}

Bu sınıf tuvalet kâğıdı, kâğıt havlu ve peçetenin yanında bebek bezi, hasta bezi ve hijyenik pedi de kapsıyor. Sınıfta 24.641 kişi çalışıyor ve üretim 5,9 milyar dolar. İhracat 1.284 milyon dolar; bebek ve hasta bezi gibi hijyen ürünlerinde Türkiye dünyanın 7. büyük ihracatçısı. İlk pazar Irak (\%27); AB'nin payı yalnızca \%18. Sınıf, birkaç büyük entegre grubun elinde ve bu gruplar kâğıt hamurundan ana bobine, ana bobinden son ürüne kadar zincirin tamamını yönetiyor.

Çıkışın sınıfa etkisi küçük. S4'te AB tarifesi \%0,8 ve istihdam etkisi −2 kişi; S2'de −15. Ormansızlaşma tüzüğü sınıfın ürünlerinin bir kısmını kapsıyor: tuvalet kâğıdı ve havlu kapsamda, bebek ve hasta bezi değil. Kapsamdaki ürünlerde hamur ithal olduğu için coğrafi konum verisi hamur tedarikçisinden gelmek zorunda. Büyük hamur üreticileri bu veriyi hazırlıyor, ama Türk üreticinin bu veriyi sözleşmelerine yazdırması gerekiyor \etiketH{}.

Kurallar sınıfı ``Yeşil Mutabakat'a maruz'' (G) olarak işaretliyordu; uzman düzeltmesiyle birincil arketip ``küresel rekabetçi, AB bağımlılığı düşük'' (B) yapıldı ve duruş ``Geliştir'' oldu. Strateji Orta Doğu, Afrika ve Orta Asya'da büyümek; Körfez ülkeleriyle bir STA bu sınıf için değerli. Sınıfın stratejik önemi sağlıktan geliyor (SS 30): salgın ve afet dönemlerinde tuvalet kâğıdı, bebek ve hasta bezi temel ihtiyaç. Bu önem, hamur tedarikinin çeşitlendirilmesini ve hamur fiyatının izlenmesini gerektiriyor; bunun için bir tetik tanımlandı. Değerlendirmenin güveni orta.

\Needspace{7\baselineskip}

\textbf{Tablo 5.4 --- 17.22 Hijyen kâğıtları ve bezler: senaryo etkileri (T+5'te yıllık; P50, parantez içinde P10 / P90)}

\nopagebreak[4]

\begingroup\gbtablo\small

\begin{longtable}[]{@{}lrrr@{}}
\toprule\noalign{}
\textbf{Değişken} & \textbf{S2} & \textbf{S3} & \textbf{S4} \\
\midrule\noalign{}
\endhead
\bottomrule\noalign{}
\endlastfoot
AB'ye ihracat (\$ mn) & −13 (−30 / +5) & −16 (−32 / −3) & −15 (−26 / −7) \\
AB'ye ihracat (\%) & −5,6 (−12,5 / +2,0) & −7,0 (−13,4 / −1,1) & −6,5 (−11,0 / −2,8) \\
Toplam ihracat (\$ mn) & −10 (−22 / +3) & −12 (−23 / −2) & −12 (−21 / −5) \\
Nihai mal ithalatı (\$ mn) & −4 (−14 / +7) & −6 (−16 / +4) & −12 (−22 / −2) \\
Üretim (\$ mn) & −5 (−8 / −3) & −6 (−8 / −3) & −1 (−4 / +3) \\
İstihdam (kişi) & −15 (−25 / −8) & −16 (−25 / −9) & −2 (−10 / +9) \\
İstihdam (\%) & −0,1 (−0,1 / 0,0) & −0,1 (−0,1 / 0,0) & 0,0 (0,0 / 0,0) \\
Cari denge (\$ mn) & −4 (−6 / −2) & −4 (−6 / −2) & −1 (−3 / +2) \\
TÜFE (bp) & +0,01 (+0,01 / +0,01) & +0,02 (+0,02 / +0,03) & +0,06 (+0,05 / +0,07) \\
\end{longtable}

\endgroup

\emph{Kaynak:} \etiketT{} f3e; toplamlar Faz 2B genel denge dağılımına kalibre. Yüzdeler 2024 düzeyine göre.

\hypertarget{nace_kartlar_16-18-58-1729-diux11fer-kuxe2ux11fux131t-uxfcruxfcnleri}{%
\subsection{5.8 17.29 Diğer kâğıt ürünleri}\label{nace_kartlar_16-18-58-1729-diux11fer-kuxe2ux11fux131t-uxfcruxfcnleri}}

\emph{GB statüsü: GB · Arketip: G (SKDM ve Yeşil Mutabakat maruz) · Duruş: Koru · AB ile müzakere: nötr · Güven: Orta · Veri sayfası: Ek C.13}

Bu sınıfta etiketler, filtre kâğıtları, kâğıt bardak ve tabaklar, gıda ambalajı kâğıtları ve kesilmiş kâğıt ürünleri var. 14.587 kişi çalışıyor. İhracat 187 milyon dolar ve \%35 kadarı AB'ye gidiyor; ithalatın \%77 kadarı da AB'den, ağırlıkla Almanya ve İtalya'dan. S4'te c2, AB'den gelen bu ürünleri ortalama \%16 pahalandırıyor ve sınıfın istihdam etkisi +145 kişi. Bu, bu parçada kâğıt ve kartondan sonraki ikinci büyük S4 kazancı.

Sınıf için asıl gündem ambalaj tüzüğü. Gıdayla temas eden kâğıt ve kartonda PFAS sınırları 12 Ağustos 2026'da başladı \etiketV{}. Yağa ve suya dayanıklılık için PFAS içeren kaplamalar kullanan üreticiler AB pazarında kaplamayı değiştirmek zorunda. Bu yüzden sınıfta ``Yeşil Mutabakat'a maruz'' (G) arketipi korundu. Duruş ise ``Koru'' yapıldı, çünkü uyum bir kaplama değişikliği; üretim teknolojisinin dönüşümünü gerektirmiyor. Önerilen eylem ortak paketteki PPWR uyumu (E-16-03). Değerlendirmenin güveni orta.

\hypertarget{nace_kartlar_16-18-59-1812-matbaacux131lux131k}{%
\subsection{5.9 18.12 Matbaacılık}\label{nace_kartlar_16-18-59-1812-matbaacux131lux131k}}

\emph{GB statüsü: GB · Arketip: --- (belirgin arketip yok) · Duruş: Koru · AB ile müzakere: nötr · Güven: Düşük · Veri sayfası: Ek C.15}

Matbaacılık bu parçanın en büyük ikinci işveren sınıfı: 47.966 kişi ve 10.768 girişim. İstihdamın \%42 kadarı İstanbul'da. Sınıf ticari baskı, ambalaj baskısı, kitap ve dergi baskısı ile güvenlikli basımı (banknot, değerli kâğıt, kimlik belgesi, oy pusulası) kapsıyor. Güvenlikli basım sınıfa küçük bir stratejik önem veriyor (SS 15).

Matbaacılığın dış ticareti bu çalışmanın verisiyle ölçülemiyor. Türkiye'nin basılı ürün ihracatı (209 milyon dolar) yayıncılık sınıflarına sayılıyor; ambalaj baskısı ise oluklu ambalajın ve kâğıt ürünlerinin içinde (VB-42). Bu yüzden sınıfın senaryo etkisi sıfır görünüyor. Gerçek etki iki kanaldan geliyor ve ikisi de girdi tarafında. Birincisi kâğıt fiyatı. Matbaaların ana girdisi kaplamalı ve kaplamasız baskı kâğıdı ve bunun büyük kısmı ithal; AB'nin payı yüksek. S4'te c2, Türkiye'de üretilmeyen kalitelerde İGV koymuyor, ama Faz 2A'nın fiyat modeline göre bütün girdilerin pahalanmasıyla matbaacılığın üretici fiyatı \%0,18--0,58 artıyor \etiketT{}. İkincisi Türkiye'nin kendi korunma önlemleri. Kaplamalı kâğıtta Aralık 2025'te açılan soruşturma, bir önlemle sonuçlanırsa sınıfın girdi maliyetini çıkıştan bağımsız olarak artıracak. Basılı ürünlerin EUDR kapsamından çıkarılması ise sınıfın AB'ye ihracatını kolaylaştırıyor \etiketV{}.

Önerilen duruş ``Koru''. Sınıfın gündemi çıkıştan çok sektörün yapısal dönüşümü: gazete ve dergi baskısının dijitalleşmeyle küçülmesi ve ambalaj baskısının büyümesi. Ticaret politikasında önerilen tek şey, kâğıt korunma önlemlerinde matbaacılığın çıkarının tartılması (E-16-04). Değerlendirmenin güveni düşük, çünkü sınıfın dış ticareti ölçülemiyor.

\hypertarget{nace_kartlar_16-18-510-etkisi-sux131nux131rlux131-sux131nux131flar}{%
\subsection{5.10 Etkisi sınırlı sınıflar}\label{nace_kartlar_16-18-510-etkisi-sux131nux131rlux131-sux131nux131flar}}

Dokuz sınıfta çıkışın etkisi küçük ya da ölçülemiyor. Hepsinin veri sayfası ve senaryo tablosu Ek C'de.

\Needspace{7\baselineskip}

\textbf{Tablo 5.5 --- Etkisi sınırlı sınıflar: temel göstergeler}

\nopagebreak[4]

\begingroup\gbtablo\footnotesize

\begin{longtable}[]{@{}
  >{\raggedright\arraybackslash}p{(\columnwidth - 16\tabcolsep) * \real{0.1439}}
  >{\raggedleft\arraybackslash}p{(\columnwidth - 16\tabcolsep) * \real{0.0830}}
  >{\raggedleft\arraybackslash}p{(\columnwidth - 16\tabcolsep) * \real{0.1066}}
  >{\raggedleft\arraybackslash}p{(\columnwidth - 16\tabcolsep) * \real{0.0787}}
  >{\raggedright\arraybackslash}p{(\columnwidth - 16\tabcolsep) * \real{0.1144}}
  >{\raggedright\arraybackslash}p{(\columnwidth - 16\tabcolsep) * \real{0.1034}}
  >{\raggedleft\arraybackslash}p{(\columnwidth - 16\tabcolsep) * \real{0.1374}}
  >{\raggedleft\arraybackslash}p{(\columnwidth - 16\tabcolsep) * \real{0.1368}}
  >{\raggedleft\arraybackslash}p{(\columnwidth - 16\tabcolsep) * \real{0.0959}}@{}}
\toprule\noalign{}
\begin{minipage}[b]{\linewidth}\raggedright
\textbf{Sınıf}
\end{minipage} & \begin{minipage}[b]{\linewidth}\raggedleft
\textbf{Çalışan}
\end{minipage} & \begin{minipage}[b]{\linewidth}\raggedleft
\textbf{İhracat (\$ mn)}
\end{minipage} & \begin{minipage}[b]{\linewidth}\raggedleft
\textbf{AB payı}
\end{minipage} & \begin{minipage}[b]{\linewidth}\raggedright
\textbf{GB statüsü}
\end{minipage} & \begin{minipage}[b]{\linewidth}\raggedright
\textbf{S4 AB tarifesi}
\end{minipage} & \begin{minipage}[b]{\linewidth}\raggedleft
\textbf{AB'ye ihracat S4 P50 (\$ mn)}
\end{minipage} & \begin{minipage}[b]{\linewidth}\raggedleft
\textbf{İstihdam S4 P50 (kişi)}
\end{minipage} & \begin{minipage}[b]{\linewidth}\raggedleft
\textbf{TÜFE S4 (bp)}
\end{minipage} \\
\midrule\noalign{}
\endhead
\bottomrule\noalign{}
\endlastfoot
16.22 Ahşap döşeme (parke) & 2.756 & 82 & \%17 & GB & \%0,0 & −1,0 & +6 & 0,01 \\
16.24 Palet ve ahşap ambalaj & 14.721 & 29 & \%45 & GB & \%3,3 & −2,4 & +20 & 0,02 \\
16.29 Diğer ağaç ürünleri & 6.039 & 89 & \%41 & GB & \%0,7 & −1,9 & +4 & 0,01 \\
17.23 Kâğıt kırtasiye & 4.538 & 69 & \%57 & GB & \%0,0 & −1,5 & 0 & 0,01 \\
17.24 Duvar kâğıdı & 495 & 24 & \%12 & GB & \%0,4 & −0,2 & +1 & 0,00 \\
18.11 Gazete basımı & 1.965 & 0 & \%0 & GB & --- & 0,0 & 0 & 0,00 \\
18.13 Baskı öncesi hizmetler & 3.170 & 16 & \%63 & GB & \%0,0 & −0,3 & +1 & 0,00 \\
18.14 Ciltçilik & 1.422 & 0 & \%0 & KAPSAM DIŞI & --- & 0,0 & 0 & 0,00 \\
18.20 Kayıtlı ortam çoğaltma & 199 & 0 & \%0 & GB & --- & 0,0 & 0 & 0,00 \\
\end{longtable}

\endgroup

\emph{Kaynak:} TÜİK YSHİ 2024, BACI 2024 \etiketV{}; senaryolar \etiketT{} f3e. Veri sayfaları Ek C'de.

Ahşap döşeme (16.22) küçük ama rekabetçi bir sınıf; AB'ye ihracatı yapı ürünü olarak CE işaretine ve AB'de tanınan deneye bağlı. Palet ve ahşap ambalaj (16.24) ağırlıkla ihracat ambalajı; başka bir ürünü taşımak için kullanıldığında ormansızlaşma tüzüğünün kapsamı dışında. S4'te c2 AB'den gelen paleti pahalandırıyor ve sınıfın istihdam etkisi +20 kişi. Diğer ağaç ürünlerinin (16.29) ve kırtasiyenin (17.23) AB'ye ihracatı küçük. Duvar kâğıdı (17.24) çok küçük bir sınıf. Gazete basımı (18.11), ciltçilik (18.14) ve kayıtlı ortamların çoğaltılması (18.20) hizmet sınıfları; dış ticaretleri ölçülemiyor ve çıkıştan etkilenmeleri beklenmiyor. Baskı öncesi hizmetler (18.13) baskı kalıbı ve klişe ihraç ediyor; ilk pazarı Almanya (\%43).

\hypertarget{nace_kartlar_16-18-6-defter-kux131sux131tlar-ve-kalite-kontroluxfc}{%
\section{6. Defter, kısıtlar ve kalite kontrolü}\label{nace_kartlar_16-18-6-defter-kux131sux131tlar-ve-kalite-kontroluxfc}}

Kümülatif defter, tamamlanan parçaların şoklarını ve önerilen önlemlerini bir araya getiriyor ve Faz 0'da konan makro kısıtlara karşı bir ara hesap tutuyor. Kısıtlar şunlar: birinci tur TÜFE etkisi yılda 30, kümülatif olarak 50 baz puanı aşmamalı; cari denge T+5'te en az sıfır olmalı; önlemlerin mali maliyeti yılda GSYH'nin \%0,5'ini geçmemeli; istihdam ve kabiliyet kayıpları sınırlanmalı. Aşağıdaki tablo 10--18 bölümlerinin, yani ilk üç parçanın 72 sınıfının toplamı. Toplamların yüzdelikleri her çekilişteki 10--18 toplamının dağılımından alınıyor. Tablonun ``bu parçanın katkısı'' satırı 16--18'i ayrıca gösteriyor.

\Needspace{7\baselineskip}

\textbf{Tablo 6.1 --- Kümülatif defter (Faz 3, parça 3 sonrası)}

\nopagebreak[4]

\begingroup\gbtablo\small

\begin{longtable}[]{@{}
  >{\raggedright\arraybackslash}p{(\columnwidth - 2\tabcolsep) * \real{0.1246}}
  >{\raggedright\arraybackslash}p{(\columnwidth - 2\tabcolsep) * \real{0.8754}}@{}}
\toprule\noalign{}
\begin{minipage}[b]{\linewidth}\raggedright
\textbf{Kalem}
\end{minipage} & \begin{minipage}[b]{\linewidth}\raggedright
\textbf{Değer}
\end{minipage} \\
\midrule\noalign{}
\endhead
\bottomrule\noalign{}
\endlastfoot
Tamamlanan / kalan sınıf & 72 / 158 C sınıfı; ek kapsam (B, D, E, hizmet) 0 / 29 \\
Kullanılan / kalan enflasyon bütçesi (bp; birinci tur, P50) & S2 yolu: şok 2,8 + önlem −0,11 = 2,7 bp; S4 yolu: şok 22,7 + önlem −13,76 = 8,9 bp. Kümülatif 50 bp bütçeden kalan: S2 47,3, S4 41,1 bp \\
Kümülatif cari denge katkısı (\$ mlr/yıl, T+5) & Senaryo P50: S2 −0,952, S4 −2,648 (S4 P10 −3,96; 1/98 askısında ortalama −3,36); önlemler: S2 +0,221, S4 +0,010 \etiketT{}/\etiketH{} \\
Kümülatif mali maliyet (mlr TL/yıl; GSYH \%) & S2 yolu 2,04 (\%0,003); S4 yolu 23,85 (\%0,032) \\
Riskteki istihdam (kişi, T+5) & S2 P50 −15.498; S4 P50 −39.733, P10 −61.720; 1/98 askısında ortalama −42.507. S4 P10'da kaybı \%5'i aşan sınıf: 13.20 (−\%5,4), 13.91 (−\%5,9), 13.92 (−\%5,3), 14.13 (−\%5,5), 14.14 (−\%8,5), 14.19 (−\%5,4), 14.31 (−\%14,9), 14.39 (−\%19,2) \\
Bu parçanın (16--18) katkısı & TÜFE şoku S2 0,2, S4 0,9 bp; önlemler S2 0,00, S4 −0,20 bp. İstihdam S2 P50 +163, S4 P50 +471 (P10 −8). Mali maliyet S2 0,39, S4 1,09 mlr TL/yıl. Cari denge önlem katkısı S2 0,000, S4 0,000 \$ mlr \\
Açık kalan veri boşlukları & VB-30: sınıf düzeyinde yabancı sermaye payı (TÜİK yabancı kontrollü girişim istatistikleri yalnız bölüm düzeyinde); VB-31: Türkiye'ye özgü sınıf düzeyi Ar-Ge yoğunluğu (OECD sınıflaması vekil); VB-32 (sürüyor): DTÖ bağlı oranları, satır düzeyi. 13--18'de bağlayıcı değil: c2 uygulanan MFN oranlarını değiştirmiyor, yalnız AB tercihini kaldırıyor; 16--18'deki tek tarife önerisi (E-17.12-01) bir indirim; VB-33: TMO yerli buğday fiyatı ile ithal paritesi farkı (menşe maliyeti c\_RoO'nun kalibrasyonu); VB-34: EUDR kapsamındaki ihracatçıların tedarik zinciri verisi; VB-35: YSHİ ile BACI sınıf eşlemesindeki tutarsızlık (10.86; 11.03, 11.04, 11.06'da YSHİ verisi yok); VB-06 (sürüyor): 1/98 tercih oranları satır düzeyinde; 1/98 askısı varyantı bu yüzden üst sınır; VB-36: firma düzeyinde kumaş menşei (giyim ihracatında PEM dışı kumaşın payı; p\_uyum'un kalibrasyonu); VB-37: PEM ortaklarında Türk kumaşının AB'ye giden üretimdeki gerçek payı (kümülasyon maruziyeti üst sınır; κ ve λ \etiketH{}); VB-38: tekstil ve giyimde 2025 istihdam daralmasının sınıf düzeyinde dağılımı (SGK yalnız bölüm düzeyinde); VB-39: deri eşya ve çantada KDV oranı (takasın kapsamı); hızlı modada sınır gecikmesinin maliyeti; VB-40: levha, kontrplak ve kâğıtta ithal odun ve kâğıt hamurunun firma düzeyinde menşe dağılımı (Rusya ve Belarus payı; EUDR ve anti-damping riski); VB-41: tamamen geri dönüştürülmüş lifle üretilen kâğıt ve ambalajın ihracattaki payı (EUDR kapsamı dışında; sınıfların EUDR payı bu yüzden üst sınır); VB-42: basılı ürün ve baskı hizmeti ihracatının sınıf eşlemesi (Fasıl 49 yayıncılık sınıflarına, 58.xx, gidiyor; 18.11, 18.12, 18.14 ve 18.20'nin dış ticareti ölçülemiyor); VB-43: c2 rejiminde yurt içi kabiliyet bayrağı HS6 düzeyinde hesaplanıyor; sıvı ambalaj kartonu (4811 59) ve kraftliner (4804 11) gibi kalitelerde yanlış sınıflama, 8--12 haneli üretim verisi gerekir; VB-44: kâğıt hamuru (17.11) YSHİ verisi gizli; sınıfın çalışanı ve üretimi bilinmiyor \\
Varsayım kaydına eklenen yeni maddeler & V-30: Faz 3 senaryoları --- toplam genel dengeden, sınıf dağılımı kısmi dengeden (ortak z, N = 2.000); V-31: S4'te 1/98'in askıya alınma olasılığı 0,25 (hasmane AB) \etiketH{}; V-32: ticaret esnekliği σ --- ProTEE HS6, sınıf içi AB'ye ihracat ağırlıklı medyan, {[}2, 15{]} aralığında; V-33: sınır sürtünmesi \%0,5--2 (yalnız GB ve İTÜ), S2'de × 0,5; AB ihracatçısının menşe maliyeti S2'de \%0,5, S3'te \%1; V-34: arketip eşikleri (KS, AME, ithalat penetrasyonu, EDE ≥ 0,3 bp/puan, SS ≥ 60) \etiketH{}; V-35 (güncellendi): SS boyut puanları ve ağırlıkları \etiketH{} (veri/f3\_ss\_boyutlar.csv; 10--18 gözden geçirildi, diğerleri ilk tahmin); V-36: CKE tanımı --- 1 puan tarife ya da \%1 üretim desteği, 5 yıl, kapasite yatırımının ithalatı (sermaye/hasıla 1, ithal payı 0,6) düşülü; V-37: KN8 → NACE eşlemesinde CPA 2.2 → 2.1 kimlik bağlantısı önceliği (D-1); V-38: Grup toplamlarının yüzdelikleri çekilişlerin toplamından hesaplanır (sınıf yüzdeliklerinin toplamı değil); 1/98 askısının etkisi askı çekilişlerinin koşullu ortalamasıyla raporlanır; V-39: Eylem takvimi T'ye göre; T2028'de T−18 geçmişte kaldığından hazırlık eylemleri 2026'nın son çeyreğinde başlar; V-40: PEM kümülasyonu kanalı yalnız S4'te; kayıp = λ · κ · maruziyet, λ \textasciitilde{} üçgen(0,3; 0,6; 0,9), κ: kumaş 1, teknik ve diğer tekstil 0,3, diğerleri 0 \etiketH{}; genel denge toplamına eklenir; V-41: tekstil, giyim ve deride menşe parametreleri revize PEM kurallarından (Karar 1/2023): giyimde p\_uyum 0,85--0,88 (Asya kumaşı payı \%10--15), kumaşta 0,92 \etiketH{}; V-42: geri ödeme yasağı (PEM Ek I md. 16) bugün AB'ye A.TR ihracatında telafi edici vergi olarak uygulandığından S2'de ek maliyet sayılmadı; V-43: KDV takası (E-13-04) giyim ve ayakkabının S4 TÜFE şokunu tasarım gereği siler; net mali maliyet Faz 2A merkez değeri (laspeyres, c2); V-44: iç ikame tavanı (veri/f3\_sinif\_varsayim.csv, ikame\_tavani): yerli üretimi AB ithalatının yerini alamayan sınıflarda iç ikame payı sınırlanır; kâğıt hamurunda (17.11) 0 \etiketH{}; V-45: CE ve onaylanmış kuruluş bağımlılığı (ce\_payi): ahşap levhada 0,15, ahşap döşemede 0,5, doğramada 0,3 (AB Yapı Ürünleri Tüzüğü; AB'ye ihracatta üçüncü taraf deneyi gereken pay) \etiketH{}; V-46: ağaç, kâğıt ve basımda menşe: Fasıl 44, 47--49 ve 4818'de başlık değişikliği; S2'de risk düşük (c\_RoO \%1, p\_uyum 0,95), S3'te \%1,5 ve 0,92 \etiketH{}; kurallar \etiketV{}; V-47: EUDR payı KN koduna göre (f3d); tamamen geri dönüştürülmüş ürün ve yalnız taşıma amaçlı ambalaj istisnası payı hesaba katılmadı (üst sınır, VB-41); V-48: E-17.12-01'in (ambalaj kâğıdında c2 düzeltmesi) etkisi: TÜFE −0,2 bp, mali maliyet AB dışı menşelerin İGV kaybı (≈ \$16 mn ≈ 0,7 mlr TL/yıl); AB menşeli ürünün c2 geliri bugün olmadığı için maliyet sayılmadı \etiketT{}; V-49: S2 ve S3'te menşe maliyeti yalnız tercihin kurtardığı tarife kadar: τ\_etkin = (1 − p\_uyum)·τ + q·p\_uyum·min(c\_RoO, τ); ithalatta min(c\_AB, τ\_TR,c2). MFN'si sıfır üründe ihracatçı tercih istemez. Sınıf ortalaması düzeyinde yaklaşık \etiketT{} \\
\end{longtable}

\endgroup

\emph{Kaynak:} \kod{kod/\allowbreak{}f3g\_\allowbreak{}kartlar.\allowbreak{}py}; matris, grup senaryoları (\kod{veri/\allowbreak{}f3\_\allowbreak{}senaryo\_\allowbreak{}grup.\allowbreak{}csv}) ve eylem kayıtlarından. Senaryo şokları politika seçimi değil, çıkışın kendisi; önlemler Bölüm 4'teki eylemlerden. Makro mutabakat Faz 6'da.

Defterin söylediği şu. Bu parça kümülatif hesabı neredeyse değiştirmiyor. Enflasyon tarafında 16--18'in S4'teki TÜFE şoku 0,9 baz puan; c2 düzeltmesiyle 0,7 baz puana iniyor. İlk üç parçanın S4 yolundaki kümülatif kullanımı önlemlerle 8,9, S2 yolunda 2,7 baz puan. 50 baz puanlık bütçenin çoğu hâlâ kullanılmamış duruyor. Kalan on üç parçada asıl sınav otomotiv, beyaz eşya, kimya ve ilaç gibi AB'nin baskın tedarikçi olduğu tüketim ve ara mallarında. Bu parçanın bu açıdan dersi c2'nin kendisiyle ilgili: kural tabanlı c2, satır düzeyinde yanlış sınıflamalarla yerli üretimi olmayan girdileri de vergilendirebiliyor. Sonraki parçalarda her sınıfın c2 satırları bu gözle denetlenecek.

Mali tarafta bu parçanın önlemleri S2 yolunda yılda 0,39, S4 yolunda 1,09 milyar TL. İlk üç parçanın kümülatif mali yükü S2 yolunda 2,04, S4 yolunda 23,85 milyar TL ve yıllık tavanın çok altında. Cari denge tarafında bu parça küçük bir olumlu katkı yapıyor: S4'te +101, S2'de +41 milyon dolar. Bu katkının kaynağı kâğıtta ithal ikamesi ve Bölüm 3.4'te anlatıldığı gibi bir üst sınır. İlk üç parçanın S4'teki kümülatif cari denge etkisi −2.648 milyon dolar; kısıtın imalatın bütününde değerlendirilmesi Faz 6'da.

İstihdam tarafında bu parça kümülatif kaybı biraz azaltıyor. S4'te ilk üç parçanın toplam istihdam etkisi −39.733 kişi (P10'da −61.720). Bu parçanın katkısı +471 kişi. S2'de kümülatif etki −15.498 kişi. Kaybın neredeyse tamamı hâlâ giyimde.

Belgenin kalite kontrolü şu noktaları kapsıyor. Her sayının birimi ve dönemi belirtildi: sınıf verileri 2024, AB'ye ihracatın rejim dağılımı ve AB tarifeleri 2025--2026, senaryolar T+5'teki yıllık fark. Modelde üç değişiklik yapıldı ve hepsi Ek B'de belgelendi: kâğıt hamurunda iç ikame tavanı (V-44), yapı ürünlerinde CE bağımlılığı (V-45) ve S2'de menşe maliyetinin yalnız tarifeli ürünlerde sayılması (V-49). Üçüncüsü önceki parçaların S2 sonuçlarını değiştirdi; etkiler Ek D'de. Bütün tarife önerileri MFN artışı içermiyor; c2 uygulanan MFN oranlarını yükseltmiyor, yalnızca AB'nin tercihini kaldırıyor. E-17.12-01 ise bir tarife indirimi ve bütün menşelere eşit uygulanıyor. Paket sanayi ürünlerinde ihracata bağlı hiçbir destek içermiyor. EUDR veri hizmeti ve hammadde programı ihracattan bağımsız, bütün kullanıcılara açık idari hizmetler.

Önceki parçalarda tanımlanan iki tipik hatadan bu parçada da kaçınıldı. Gümrük Birliği dışında kalan bir tarifeye çıkışa bağlı bir artış eklenmedi; bu parçada tarım rejimi ya da işlenmiş tarım ürünü yok. Olumsuz bulgular gizlenmedi. S4'te kâğıtta bulunan istihdam artışının abartılı olduğu ve parçanın asıl risklerinin, yani ticaret savunması, Rus odunu ve EUDR uyumunun, modelin dışında kaldığı yönetici özetinde ve ilgili sınıf incelemelerinde açıkça yazıldı. Belirsiz olgular \etiketD{} olarak işaretlendi ve veri boşlukları Ek E'de listelendi.

\hypertarget{nace_kartlar_16-18-ek-a--ana-matris-satux131rlarux131}{%
\section{Ek A --- Ana matris satırları}\label{nace_kartlar_16-18-ek-a--ana-matris-satux131rlarux131}}

Ana matrisin bütün sütunları ve 230 imalat sınıfının tamamı \kod{veri/\allowbreak{}nace\_\allowbreak{}matris.\allowbreak{}csv} dosyasında. Aşağıdaki iki tablo çalışma çerçevesinin öngördüğü matris sütunlarını 16--18 sınıfları için veriyor. Skorlar 0--100 arasında, C sınıfları içindeki yüzdelik sıra olarak ifade ediliyor.

\begin{landscape}

\textbf{Tablo A.1 --- Ana matris: kimlik, ticaret ve maruziyet}

\nopagebreak[4]

\begingroup\gbtablo\footnotesize

\begin{longtable}[]{@{}
  >{\raggedright\arraybackslash}p{(\columnwidth - 26\tabcolsep) * \real{0.0402}}
  >{\raggedright\arraybackslash}p{(\columnwidth - 26\tabcolsep) * \real{0.1530}}
  >{\raggedright\arraybackslash}p{(\columnwidth - 26\tabcolsep) * \real{0.1033}}
  >{\raggedleft\arraybackslash}p{(\columnwidth - 26\tabcolsep) * \real{0.0402}}
  >{\raggedleft\arraybackslash}p{(\columnwidth - 26\tabcolsep) * \real{0.0684}}
  >{\raggedleft\arraybackslash}p{(\columnwidth - 26\tabcolsep) * \real{0.0600}}
  >{\raggedleft\arraybackslash}p{(\columnwidth - 26\tabcolsep) * \real{0.0590}}
  >{\raggedleft\arraybackslash}p{(\columnwidth - 26\tabcolsep) * \real{0.0572}}
  >{\raggedleft\arraybackslash}p{(\columnwidth - 26\tabcolsep) * \real{0.1005}}
  >{\raggedleft\arraybackslash}p{(\columnwidth - 26\tabcolsep) * \real{0.0747}}
  >{\raggedright\arraybackslash}p{(\columnwidth - 26\tabcolsep) * \real{0.0903}}
  >{\raggedleft\arraybackslash}p{(\columnwidth - 26\tabcolsep) * \real{0.0734}}
  >{\raggedleft\arraybackslash}p{(\columnwidth - 26\tabcolsep) * \real{0.0383}}
  >{\raggedleft\arraybackslash}p{(\columnwidth - 26\tabcolsep) * \real{0.0415}}@{}}
\toprule\noalign{}
\begin{minipage}[b]{\linewidth}\raggedright
\textbf{Kod}
\end{minipage} & \begin{minipage}[b]{\linewidth}\raggedright
\textbf{Ad}
\end{minipage} & \begin{minipage}[b]{\linewidth}\raggedright
\textbf{GB statüsü}
\end{minipage} & \begin{minipage}[b]{\linewidth}\raggedleft
\textbf{KD}
\end{minipage} & \begin{minipage}[b]{\linewidth}\raggedleft
\textbf{İstihdam}
\end{minipage} & \begin{minipage}[b]{\linewidth}\raggedleft
\textbf{İhracat}
\end{minipage} & \begin{minipage}[b]{\linewidth}\raggedleft
\textbf{AB payı}
\end{minipage} & \begin{minipage}[b]{\linewidth}\raggedleft
\textbf{İthalat}
\end{minipage} & \begin{minipage}[b]{\linewidth}\raggedleft
\textbf{AB ithalat payı}
\end{minipage} & \begin{minipage}[b]{\linewidth}\raggedleft
\textbf{τ\_AB (S4)}
\end{minipage} & \begin{minipage}[b]{\linewidth}\raggedright
\textbf{Menşe S2/S3}
\end{minipage} & \begin{minipage}[b]{\linewidth}\raggedleft
\textbf{İthal girdi}
\end{minipage} & \begin{minipage}[b]{\linewidth}\raggedleft
\textbf{RCA}
\end{minipage} & \begin{minipage}[b]{\linewidth}\raggedleft
\textbf{PCI}
\end{minipage} \\
\midrule\noalign{}
\endhead
\bottomrule\noalign{}
\endlastfoot
16.10 & Kereste & GB & 246 & 15.956 & 101 & \%24 & 327 & \%13 & \%0,0 & D/D & \%22 & 0,17 & −0,33 \\
16.21 & Ahşap levha ve kontrplak & GB & 1.010 & 22.174 & 1.350 & \%16 & 287 & \%36 & \%6,9 & D/D & \%22 & 2,61 & −0,25 \\
16.22 & Ahşap döşeme (parke) & GB & 115 & 2.756 & 82 & \%17 & 26 & \%53 & \%0,0 & D/D & \%22 & 1,55 & −0,30 \\
16.23 & Doğrama ve ahşap yapı & GB & 248 & 30.303 & 227 & \%22 & 49 & \%66 & \%1,2 & D/D & \%22 & 1,19 & −0,16 \\
16.24 & Palet ve ahşap ambalaj & GB & 266 & 14.721 & 29 & \%45 & 28 & \%62 & \%3,3 & D/D & \%22 & 0,39 & 0,16 \\
16.29 & Diğer ağaç ürünleri & GB & 46 & 6.039 & 89 & \%41 & 116 & \%31 & \%0,7 & D/D & \%22 & 0,30 & −0,57 \\
17.11 & Kâğıt hamuru & GB & --- & --- & 117 & \%0 & 1.131 & \%39 & \%0,0 & D/D & \%28 & 0,19 & 0,24 \\
17.12 & Kâğıt ve karton & GB & 759 & 13.598 & 1.372 & \%25 & 2.526 & \%63 & \%0,0 & D/D & \%28 & 1,02 & −0,10 \\
17.21 & Oluklu ambalaj & GB & 1.435 & 40.967 & 834 & \%39 & 90 & \%76 & \%0,0 & D/D & \%28 & 2,24 & −0,39 \\
17.22 & Hijyen kâğıtları ve bezler & GB & 1.281 & 24.641 & 1.284 & \%18 & 186 & \%51 & \%0,8 & D/D & \%28 & 2,77 & −0,45 \\
17.23 & Kâğıt kırtasiye & GB & 163 & 4.538 & 69 & \%57 & 23 & \%32 & \%0,0 & D/D & \%28 & 0,86 & −0,67 \\
17.24 & Duvar kâğıdı & GB & 27 & 495 & 24 & \%12 & 8 & \%63 & \%0,4 & D/D & \%28 & 1,39 & 0,35 \\
17.29 & Diğer kâğıt ürünleri & GB & 520 & 14.587 & 187 & \%35 & 147 & \%77 & \%0,0 & D/D & \%28 & 0,95 & 0,05 \\
18.11 & Gazete basımı & GB & 25 & 1.965 & 0 & \%0 & 0 & \%0 & --- & D/D & \%26 & --- & --- \\
18.12 & Matbaacılık & GB & 1.069 & 47.966 & 0 & \%0 & 0 & \%0 & --- & D/D & \%26 & --- & --- \\
18.13 & Baskı öncesi hizmetler & GB & 60 & 3.170 & 16 & \%63 & 17 & \%52 & \%0,0 & D/D & \%26 & 1,90 & 0,40 \\
18.14 & Ciltçilik & KAPSAM DIŞI & 25 & 1.422 & 0 & \%0 & 0 & \%0 & --- & D/D & \%26 & --- & --- \\
18.20 & Kayıtlı ortam çoğaltma & GB & 5 & 199 & 0 & \%0 & 0 & \%0 & --- & D/D & \%26 & --- & --- \\
\end{longtable}

\endgroup

\end{landscape}

\emph{Kaynak:} \kod{veri/\allowbreak{}nace\_\allowbreak{}matris.\allowbreak{}csv}. KD, ihracat ve ithalat \$ mn (2024). τ\_AB: S4'te eklenecek AB tarifesi (1/98 sürer varsayımı).

\begin{landscape}

\textbf{Tablo A.2 --- Ana matris: skorlar, strateji ve senaryo}

\nopagebreak[4]

\begingroup\gbtablo\footnotesize

\begin{longtable}[]{@{}
  >{\raggedright\arraybackslash}p{(\columnwidth - 28\tabcolsep) * \real{0.0408}}
  >{\raggedleft\arraybackslash}p{(\columnwidth - 28\tabcolsep) * \real{0.0285}}
  >{\raggedleft\arraybackslash}p{(\columnwidth - 28\tabcolsep) * \real{0.0285}}
  >{\raggedleft\arraybackslash}p{(\columnwidth - 28\tabcolsep) * \real{0.0420}}
  >{\raggedleft\arraybackslash}p{(\columnwidth - 28\tabcolsep) * \real{0.0388}}
  >{\raggedleft\arraybackslash}p{(\columnwidth - 28\tabcolsep) * \real{0.0388}}
  >{\raggedright\arraybackslash}p{(\columnwidth - 28\tabcolsep) * \real{0.0580}}
  >{\raggedright\arraybackslash}p{(\columnwidth - 28\tabcolsep) * \real{0.0556}}
  >{\raggedright\arraybackslash}p{(\columnwidth - 28\tabcolsep) * \real{0.1724}}
  >{\raggedright\arraybackslash}p{(\columnwidth - 28\tabcolsep) * \real{0.1885}}
  >{\raggedright\arraybackslash}p{(\columnwidth - 28\tabcolsep) * \real{0.0715}}
  >{\raggedleft\arraybackslash}p{(\columnwidth - 28\tabcolsep) * \real{0.0560}}
  >{\raggedleft\arraybackslash}p{(\columnwidth - 28\tabcolsep) * \real{0.0576}}
  >{\raggedright\arraybackslash}p{(\columnwidth - 28\tabcolsep) * \real{0.0728}}
  >{\raggedright\arraybackslash}p{(\columnwidth - 28\tabcolsep) * \real{0.0503}}@{}}
\toprule\noalign{}
\begin{minipage}[b]{\linewidth}\raggedright
\textbf{Kod}
\end{minipage} & \begin{minipage}[b]{\linewidth}\raggedleft
\textbf{KS}
\end{minipage} & \begin{minipage}[b]{\linewidth}\raggedleft
\textbf{SS}
\end{minipage} & \begin{minipage}[b]{\linewidth}\raggedleft
\textbf{AME}
\end{minipage} & \begin{minipage}[b]{\linewidth}\raggedleft
\textbf{EDE}
\end{minipage} & \begin{minipage}[b]{\linewidth}\raggedleft
\textbf{CKE}
\end{minipage} & \begin{minipage}[b]{\linewidth}\raggedright
\textbf{Arketip}
\end{minipage} & \begin{minipage}[b]{\linewidth}\raggedright
\textbf{Duruş}
\end{minipage} & \begin{minipage}[b]{\linewidth}\raggedright
\textbf{Tarife yolu}
\end{minipage} & \begin{minipage}[b]{\linewidth}\raggedright
\textbf{Ana araçlar}
\end{minipage} & \begin{minipage}[b]{\linewidth}\raggedright
\textbf{AB müzakere}
\end{minipage} & \begin{minipage}[b]{\linewidth}\raggedleft
\textbf{S2 cari P50}
\end{minipage} & \begin{minipage}[b]{\linewidth}\raggedleft
\textbf{S2 TÜFE P50}
\end{minipage} & \begin{minipage}[b]{\linewidth}\raggedright
\textbf{İstihdam riski}
\end{minipage} & \begin{minipage}[b]{\linewidth}\raggedright
\textbf{Güven}
\end{minipage} \\
\midrule\noalign{}
\endhead
\bottomrule\noalign{}
\endlastfoot
16.10 & 39 & 0 & 22 & 0,02 & 65 & --- & Koru & AB ile karşılıklı \%0 (S2); S4'te AB'ye c2; tomruk ve kerestede MFN \%0 sürer (P1) & EUDR veri hizmeti (E-16-01); Rus ve Belarus odunundan ayrışma (E-16-02); hammadde güvenliği (E-16-05) & nötr & +0,1 & 0,00 & Düşük & Orta \\
16.21 & 51 & 0 & 53 & 0,01 & 53 & B/G & Geliştir & AB ile karşılıklı \%0 (S2); S4'te AB'ye c2; yonga, lif, tomruk ve kaplamada \%0 (P1) & EUDR veri hizmeti (E-16-01); Rus ve Belarus odunundan ayrışma ve ticaret savunması hazırlığı (E-16-02); Körfez, Orta Asya ve Afrika pazarları & savunma & −5,6 & 0,01 & Düşük & Orta \\
16.22 & 50 & 0 & 19 & 0,00 & 46 & B & Geliştir & AB ile karşılıklı \%0 (S2); S4'te AB'ye c2 & CE işareti için AB'de tanınan deney (S2'de ACAA); yeni pazarlar & nötr & +0,3 & 0,00 & Düşük & Yüksek \\
16.23 & 43 & 0 & 27 & 0,00 & 50 & --- & Koru & AB ile karşılıklı \%0 (S2); S4'te AB'ye c2 & CE işareti (pencere ve dış kapı) için AB'de tanınan deney; EUDR verisi (E-16-01) & nötr & −1,3 & 0,00 & Düşük & Orta \\
16.24 & 43 & 7 & 70 & 0,00 & 48 & --- & Koru & Değişiklik yok (çoğunlukla ihracat ambalajı) & ISPM 15 ısıl işlem sertifikası; taşıma ambalajı olarak kullanıldığında EUDR kapsamı dışında & nötr & +0,8 & 0,00 & Düşük & Yüksek \\
16.29 & 38 & 0 & 40 & 0,01 & 65 & --- & Koru & AB ile karşılıklı \%0 (S2); S4'te AB'ye c2 & İzleme; EUDR verisi (E-16-01) & nötr & −0,4 & 0,00 & Düşük & Yüksek \\
17.11 & 48 & 31 & 10 & 0,09 & 3 & --- & Koru & Bugün de bütün menşelere \%0 (P1); değişiklik yok & Tedarik çeşitlendirme ve EUDR verisi (E-16-01); lif hammaddesi güvenliği (E-16-05) & nötr & 0,0 & 0,02 & Düşük & Düşük \\
17.12 & 47 & 23 & 23 & 0,21 & 74 & C/G & Geliştir & AB ile karşılıklı \%0 (S2); S4'te c2, yurt içi üretimi olmayan ambalaj girdisi kalitelerinde İGV yok (E-17.12-01); hamur ve atık kâğıtta \%0 & Geri dönüştürülmüş lifle kaplamalı karton ve testliner kapasitesi; atık kâğıt toplama (E-16-05); korunma önlemlerinde aşağı akışın gözetilmesi (E-16-04) & nötr & +48,1 & 0,13 & Düşük & Orta \\
17.21 & 48 & 23 & 31 & 0,01 & 37 & A/G & Koru & AB ile karşılıklı \%0 (S2); AB'nin MFN'si zaten \%0; S4'te AB'ye c2 & PPWR uyumu (E-16-03); girdide c2 düzeltmesi (E-17.12-01) & nötr & +0,2 & 0,01 & Düşük & Orta \\
17.22 & 44 & 30 & 24 & 0,02 & 37 & B/G & Geliştir & AB ile karşılıklı \%0 (S2); S4'te AB'ye c2; kâğıt hamuru ve süper emicide \%0 (P1) & Kâğıt hamurunda tedarik çeşitlendirme ve EUDR verisi (E-16-01); Orta Doğu, Afrika ve Orta Asya pazarları & nötr & −3,7 & 0,01 & Düşük & Orta \\
17.23 & 40 & 0 & 42 & 0,00 & 38 & --- & Koru & AB ile karşılıklı \%0 (S2); S4'te AB'ye c2 & İzleme & nötr & −0,2 & 0,00 & Düşük & Yüksek \\
17.24 & 61 & 0 & 18 & 0,00 & 57 & --- & Koru & AB ile karşılıklı \%0 (S2); S4'te AB'ye c2 & İzleme & nötr & +0,1 & 0,00 & Düşük & Yüksek \\
17.29 & 52 & 8 & 28 & 0,01 & 47 & G & Koru & AB ile karşılıklı \%0 (S2); S4'te AB'ye c2 & PPWR uyumu, gıdayla temas eden kâğıtta PFAS'sız kaplama (E-16-03) & nötr & +5,5 & 0,02 & Düşük & Orta \\
18.11 & 42 & 0 & 5 & 0,00 & 5 & --- & Koru & Değişiklik yok (hizmet; mal ticareti yok) & İzleme & nötr & 0,0 & 0,00 & Düşük & Düşük \\
18.12 & 48 & 15 & 5 & 0,00 & 5 & --- & Koru & Değişiklik yok (hizmet); kâğıt girdisinde E-17.12-01 & Kaplamalı kâğıt korunma önleminde aşağı akış değerlendirmesi (E-16-04); güvenlikli basım kapasitesi & nötr & 0,0 & 0,00 & Düşük & Düşük \\
18.13 & 58 & 0 & 47 & 0,00 & 37 & B & Geliştir & AB ile karşılıklı \%0 (S2); S4'te AB'ye c2 & İzleme & nötr & +0,1 & 0,00 & Düşük & Yüksek \\
18.14 & 52 & 0 & 5 & 0,00 & 5 & --- & Koru & Değişiklik yok (hizmet) & İzleme & nötr & 0,0 & 0,00 & Düşük & Düşük \\
18.20 & 51 & 0 & 5 & 0,00 & 5 & --- & Koru & Değişiklik yok & İzleme & nötr & 0,0 & 0,00 & Düşük & Düşük \\
\end{longtable}

\endgroup

\end{landscape}

\emph{Kaynak:} Skorlar 0--100 (C sınıfları içinde yüzdelik); EDE bp/puan; S2 cari denge \$ mn; TÜFE bp. Arketip ve duruş kural tabanlı öneri ve uzman düzeltmesi.

\hypertarget{nace_kartlar_16-18-ek-b--yuxf6ntem-bu-paruxe7adaki-yenilikler}{%
\section{Ek B --- Yöntem: bu parçadaki yenilikler}\label{nace_kartlar_16-18-ek-b--yuxf6ntem-bu-paruxe7adaki-yenilikler}}

Hesap zinciri ilk parçanın Ek B'sinde anlatıldı ve değişmedi: f3a--f3f bütün 230 sınıfı birlikte hesaplıyor, f3g her parçanın belgesini üretiyor. Bu parçada zincire iki ekleme ve bir düzeltme yapıldı, parçaya özgü varsayımlar derlendi. Tablo B.1 değişen betikleri ve tabloları özetliyor.

\Needspace{7\baselineskip}

\textbf{Tablo B.1 --- Bu parçada değişen betikler ve tablolar}

\nopagebreak[4]

\begingroup\gbtablo\small

\begin{longtable}[]{@{}
  >{\raggedright\arraybackslash}p{(\columnwidth - 4\tabcolsep) * \real{0.1559}}
  >{\raggedright\arraybackslash}p{(\columnwidth - 4\tabcolsep) * \real{0.6914}}
  >{\raggedright\arraybackslash}p{(\columnwidth - 4\tabcolsep) * \real{0.1527}}@{}}
\toprule\noalign{}
\begin{minipage}[b]{\linewidth}\raggedright
\textbf{Betik ya da tablo}
\end{minipage} & \begin{minipage}[b]{\linewidth}\raggedright
\textbf{Değişiklik}
\end{minipage} & \begin{minipage}[b]{\linewidth}\raggedright
\textbf{Etiket}
\end{minipage} \\
\midrule\noalign{}
\endhead
\bottomrule\noalign{}
\endlastfoot
\kod{kod/\allowbreak{}f3e\_\allowbreak{}senaryo.\allowbreak{}py} & İç ikame tavanı (\kod{ikame\_\allowbreak{}tavani}) ve CE ve onaylanmış kuruluş bağımlılığı (\kod{ce\_\allowbreak{}payi}) sınıf tablosundan okunuyor; S2 ve S3'te menşe maliyeti yalnız tercihin kurtardığı tarife kadar sayılıyor (ihracatta ve ithalatta) & \etiketT{}; parametreler \etiketH{} \\
\kod{kod/\allowbreak{}f3g\_\allowbreak{}kartlar.\allowbreak{}py} & Bölge tablosunun çalışana göre sıralanan biçimi (\kod{bolge\_\allowbreak{}l}) & --- \\
\kod{kod/\allowbreak{}requirements.\allowbreak{}txt} & pyfixest 0.60.0'a sabitlendi (0.70, f2f'nin kullandığı bir parametreyi kaldırdı); PyYAML ve matplotlib eklendi & --- \\
\kod{veri/\allowbreak{}f3\_\allowbreak{}sinif\_\allowbreak{}varsayim.\allowbreak{}csv} & 16--18 için menşe kuralları, c\_RoO, p\_uyum, değer zinciri rolü ve DİR bağımlılığı; iki yeni sütun & \etiketH{}; kurallar \etiketV{} \\
\kod{veri/\allowbreak{}f3\_\allowbreak{}ss\_\allowbreak{}boyutlar.\allowbreak{}csv} & 16--18 için stratejik önem boyutları gözden geçirildi & \etiketH{} \\
\end{longtable}

\endgroup

\textbf{İç ikame tavanı.} Kısmi denge modelinde AB'den ithalattaki düşüşün bir kısmını yerli üretim karşılıyor: ΔD = −ΔM\_AB · D/(D + M\_AB-dışı), D sınıfın iç satışı. Bu formül, yerli üretimin aynı ürünü üretebildiğini varsayıyor. Kâğıt hamurunda bu varsayım geçerli değil: yurt içi üretim yok denecek kadar az ve sınıfın verisi gizli olduğu için formül, ithal paylarına dayanan bir yedek değer kullanıyordu. Sınıf tablosuna eklenen \kod{ikame\_\allowbreak{}tavani} sütunu, iç ikame payını bir üst sınırla kısıtlıyor; boş bırakılan sınıflarda bir değişiklik yok. Bu parçada yalnızca kâğıt hamurunda tavan tanımlandı ve değeri sıfır (V-44). Kâğıt ve kartonda tavan gerekmedi. AB'den ithalattaki düşüş, c2'nin İGV koyduğu ve yerli üretimin güçlü olduğu satırlarda yoğunlaşıyor; sınıfın ortalama iç ikame payı (\%65) bu yoğunlaşmayla tutarlı. Yanlış sınıflanan satırlar (VB-43) E-17.12-01 ile ele alınıyor.

\textbf{CE ve onaylanmış kuruluş bağımlılığı.} Model, AB'ye satışında AB'nin tanıdığı bir kuruluşun deneyini ya da belgesini gerektiren ihracat payına, S4'te değerin \%1--4'ü, S2'de bunun dörtte biri kadar bir maliyet ekliyor. Bu kanal ilk parçadan beri modeldeydi, ama bütün sınıflarda pay sıfırdı. Bu parçada pay, sınıf tablosundaki \kod{ce\_\allowbreak{}payi} sütunundan okunuyor ve yapı ürünleri için tanımlandı: ahşap levhada 0,15 (taşıyıcı levha ve yangın tepkisi deneyi gereken ürünler), ahşap döşemede 0,5, doğramada 0,3 (pencere ve dış kapı) (V-45). Elektrikli ekipman, makine, otomotiv, tıbbi cihaz ve oyuncakta bu pay sonraki parçalarda belirlenecek ve orada çok daha belirleyici olacak.

\textbf{S2 ve S3'te menşe maliyeti.} Serbest ticaret anlaşmasında ihracatçı, tercihli tarifeden yararlanmak için menşe ispatlamak zorunda; bunun maliyeti modelde c\_RoO. Ama tercihin bir değeri ancak MFN tarifesi pozitifse var. AB'nin MFN tarifesi sıfır olan bir üründe ihracatçı tercih istemez, menşe ispatına da katlanmaz; MFN'yi öder ve bu sıfırdır. Model şimdiye kadar menşe maliyetini tarifeden bağımsız olarak, menşe kuralına uyan bütün ihracata ekliyordu. Bu parçada fiyat kaması τ\_etkin = (1 − p\_uyum)·τ + q·p\_uyum·min(c\_RoO, τ) + NTM + CE biçimine getirildi: uyumlu ihracatçı, menşe maliyeti ile tarifenin küçüğünü öder. Aynı düzeltme ithalat tarafında AB'li ihracatçının menşe maliyetine de uygulandı. Düzeltme sınıf ortalaması düzeyinde yaklaşık, çünkü bir sınıfın içinde tarifeli ve tarifesiz satırlar bir arada olabilir (V-49). S2'nin toplam AB ihracatı kaybı genel denge modelinden geldiği için düzeltme bu toplamı değiştirmiyor; kaybın sınıflara dağılımını değiştiriyor. Tarifesi sıfır ya da çok düşük sınıfların (kâğıt, bazı makine, mobilya, gemi ve tıbbi cihaz sınıfları) S2 kaybı küçülüyor, tarifeli ve emek yoğun sınıflarınki büyüyor (Ek D).

\textbf{Menşe parametreleri.} Ağaç, kâğıt ve basım sınıflarının menşe kuralı revize PEM kurallarından alındı; Fasıl 44, 47 ve 48'in ürünlerinin büyük çoğunluğunda kural başlık değişikliği \etiketV{}. Bu yüzden bütün sınıflarda S2 menşe riski düşük: uyum oranı (p\_uyum) 0,95, uyumlu ihracatın menşe maliyeti (c\_RoO) \%1; S3'te 0,92 ve \%1,5 (V-46). Geri ödeme yasağı bu fasıllarda uygulanmıyor.

\textbf{EUDR payı.} Sınıfların EUDR payı (f3d), AB'ye ihracatın EUDR kapsamındaki KN kodlarına giren kısmı. Tüzük 2025/2650 ile basılı ürünlerin (Fasıl 49) kapsamdan çıkarılması bu parçanın sınıflarını etkilemedi, çünkü Fasıl 49 yayıncılık sınıflarına sayılıyor. Tamamen geri dönüştürülmüş liften yapılan ürünlerin ve yalnızca taşıma amaçlı kullanılan ambalajın kapsam dışı olması ise KN koduyla ayrılamıyor; bu yüzden pay bir üst sınır (V-47, VB-41).

\textbf{Yeni ortamda yeniden üretim.} Bu parça yeni bir çalışma ortamında hazırlandı. Paketin bütün betikleri çevrimdışı kipte, yalnızca paketteki girdilerle yeniden çalıştırıldı. Model değişikliklerinden önce, ikinci parçanın çıktılarının tamamı bayt bayt aynı üretildi; tek fark, pyfixest'in yeni sürümünün f2f'nin kullandığı bir parametreyi kaldırmasıydı ve sürüm sabitlendi. Bu parçanın sonundaki yeniden üretim sınaması README'de.

\hypertarget{nace_kartlar_16-18-ek-c--sux131nux131f-veri-sayfalarux131}{%
\section{Ek C --- Sınıf veri sayfaları}\label{nace_kartlar_16-18-ek-c--sux131nux131f-veri-sayfalarux131}}

Her veri sayfası sınıf kartının sayısal alanlarını bir arada veriyor. Senaryo tablosu Bölüm 5'te verilen sınıflar için sayfada ilgili tabloya atıf var; diğer sınıfların senaryo tablosu sayfanın hemen altında. Veriler \etiketV{}; eşleme, bölge payları ve skorlar \etiketT{}; menşe parametreleri, CE payı, iç ikame tavanı ve stratejik skor \etiketH{}.

\hypertarget{nace_kartlar_16-18-c1-1610-aux11fauxe7larux131n-biuxe7ilmesi-ve-planyalanmasux131}{%
\subsection{C.1 16.10 Ağaçların biçilmesi ve planyalanması}\label{nace_kartlar_16-18-c1-1610-aux11fauxe7larux131n-biuxe7ilmesi-ve-planyalanmasux131}}

\Needspace{7\baselineskip}

\textbf{Tablo C.1 --- 16.10: veri sayfası}

\nopagebreak[4]

\begingroup\gbtablo\small

\begin{longtable}[]{@{}
  >{\raggedright\arraybackslash}p{(\columnwidth - 4\tabcolsep) * \real{0.1345}}
  >{\raggedright\arraybackslash}p{(\columnwidth - 4\tabcolsep) * \real{0.3484}}
  >{\raggedright\arraybackslash}p{(\columnwidth - 4\tabcolsep) * \real{0.5171}}@{}}
\toprule\noalign{}
\begin{minipage}[b]{\linewidth}\raggedright
\textbf{Alan}
\end{minipage} & \begin{minipage}[b]{\linewidth}\raggedright
\textbf{Gösterge}
\end{minipage} & \begin{minipage}[b]{\linewidth}\raggedright
\textbf{Değer}
\end{minipage} \\
\midrule\noalign{}
\endhead
\bottomrule\noalign{}
\endlastfoot
1 Kimlik & Rev. 2.1 karşılığı; CPA & 16.11, 16.12 (1:2); CPA 16.10 \\
& Başlıca HS4 (ihracat, 2024) & 4407 \%50, 4409 \%40, 4403 \%7, 4404 \%2 \\
& BEC payları: ara / sermaye / tüketim & \%100 / \%0 / \%0 \\
2 Ölçek & Üretim; katma değer (GSYH payı) & \$1.088 mn; \$246 mn (\%0,02) \\
& Çalışan; girişim & 15.956; 3.425 \\
& Girişim ölçeği: 0 / 1--9 / 10--49 / 50--249 / 250+ çalışan & 84 / 3.312 / 277 / 7 / 1 \\
& 250+ girişimlerin istihdam payı; kişi başı KD (AB27 = 1) & \%2; 0,31 \\
& İstihdam değişimi, 2015--2024 & \%16 \\
& İlk üç bölge (İBBS-2) & Kocaeli, Sakarya, Düzce, Bolu, Yalova \%18; Samsun, Tokat, Çorum, Amasya \%9; Kastamonu, Çankırı, Sinop \%8 \\
& Çapa yatırımcılar & OGM'den tomruk alan küçük ve orta ölçekli bıçkıhaneler; Karadeniz (Kastamonu, Bolu, Düzce, Samsun) ve Kocaeli--Sakarya kümeleri \etiketH{} \\
3 Dış ticaret & İhracat 2024; 2020--2024 değişimi & \$101 mn; \%42 \\
& AB payı; ilk beş AB pazarı & \%24; Romanya \%28; Bulgaristan \%20; İtalya \%9; Portekiz \%9; İspanya \%6 \\
& İlk beş AB dışı pazar & İran \%30; Irak \%15; ABD \%9; Kıbrıs \%4; Libya \%4 \\
& İthalat 2024; AB ve Çin payı & \$327 mn; AB \%13, Çin \%1 \\
& İlk beş ithalat kaynağı & Ukrayna \%24; ABD \%24; Rusya \%13; Kamerun \%9; Belarus \%8 \\
& Dış ticaret dengesi & −\$227 mn \\
& İthal girdi içeriği (I-O); KD ÷ üretim & \%22; \%23 \\
& RCA: dünya 2020 → 2024; AB pazarı 2024 & 0,13 → 0,17; 0,24 \\
& Birim değer ÷ AB pazarı ortalaması & 1,92 \\
& AB pazarındaki pay (AB dışı tedarik içinde) 2020 → 2024; Türkiye'nin sırası & \%0,2 → \%0,9; 19. \\
4 AB tarafı & AB'ye ihracat 2025; rejim payları (GB / İTÜ / tarım) & €21,7 mn; \%100 / \%0 / \%0 \\
& S4'te AB tarifesi: 1/98 sürerse / askıda / tam MFN ortalaması & \%0,0 / \%0,0 / \%0,0 \\
& Tepe satırlar (S4 artışı ≥ \%10) & yok \\
& AB pazarında AB dışı rakipler (AB tercihi) & Norveç \%13 (AEA/STA); Ukrayna \%12 (STA/EPA); Brezilya \%10 (MFN/GSP); Kamerun \%9 (STA/EPA); ABD \%8 (MFN/GSP) \\
& Tercihli rakiplerin payı; AME (puan; ham) & \%53; 22 (0,48) \\
& Menşe kuralı (PEM) & Fasıl 44 (4406--4407): başlık değişikliği; ithal tomruktan biçilmiş kereste menşe kazanır \\
& Menşe riski S2 / S3; c\_RoO (\%); p\_uyum & D/D; 1,0 / 1,5; 0,95 / 0,92 \etiketH{} \\
& Değer zincirindeki rol; DİR bağımlılığı & ara; düşük \\
& EUDR ve SKDM kapsamındaki AB'ye ihracat payı & \%100; \%0 \\
& PEM ortaklarına ihracat; kümülasyon maruziyeti; κ & \$6 mn; \$3 mn; 0,0 \etiketH{} \\
5 Türkiye tarafı & Uygulanan tarife 2026 (ithalat ağırlıklı): MFN; İGV dahil & \%0,2; \%0,2 \\
& S4'te AB menşeline tarife (c2) & \%0,0 \\
& Kritik girdiler: 2024 ithalatı, AB payı, ilk kaynaklar & HS4 4403, 4407: \$252 mn (AB \%17; Ukrayna \%30, Kamerun \%12, Rusya \%11) \\
6 Kabiliyet & KS; sıra (ağırlık duyarlılığında P10--P90) & 38,7; 209 / 230 (172--216) \\
& KS bileşenleri (yüzdelik): RCA, PCI, yurt içi KD, Ar-Ge, tedarikçi, verimlilik, firma & 29, 31, 49, 25, 62, 17, 78 \\
& PCI (sınıf) & −0,33 \\
& Atlanabilir ürünler (yoğunluk üst çeyrek, PCI \textgreater{} sınıf) & 440791 (PCI +0,68, RCA 0,11); 440140 (PCI +0,58, RCA 0,17); 440410 (PCI −0,30, RCA 0,16) \\
7 Stratejik önem & SS; en yüksek boyut; ağırlıklı ortalama & 0,0; --- 0; 0,0 \etiketH{} \\
8 Makro duyarlılık & TÜFE ağırlığı (doğrudan); tüketimde ithal pay & \%0,00; \%95 \\
& ÜFE vekili (imalat üretimindeki pay) & \%0,16 \\
& EDE, 1 puan tarife: tüm menşeler / yalnız AB (bp; ρ = 1) & 0,017 / 0,002 \\
& CKE: tarife (\$ mn / bp); mali (\$ mn / mlr TL); puan & 4.293,5; 26,31; 65 \\
9 Senaryolar & Senaryo tablosu & Tablo C.1b (aşağıda) \\
\end{longtable}

\endgroup

\Needspace{7\baselineskip}

\textbf{Tablo C.1b --- 16.10 Kereste: senaryo etkileri (T+5'te yıllık; P50, parantez içinde P10 / P90)}

\nopagebreak[4]

\begingroup\gbtablo\small

\begin{longtable}[]{@{}lrrr@{}}
\toprule\noalign{}
\textbf{Değişken} & \textbf{S2} & \textbf{S3} & \textbf{S4} \\
\midrule\noalign{}
\endhead
\bottomrule\noalign{}
\endlastfoot
AB'ye ihracat (\$ mn) & −1 (−2 / 0) & −1 (−3 / 0) & −1 (−2 / 0) \\
AB'ye ihracat (\%) & −3,0 (−7,2 / +2,0) & −5,1 (−10,5 / −0,8) & −4,8 (−8,6 / −2,0) \\
Toplam ihracat (\$ mn) & −1 (−1 / 0) & −1 (−2 / 0) & −1 (−2 / 0) \\
Nihai mal ithalatı (\$ mn) & −1 (−2 / +2) & −1 (−3 / +1) & −1 (−2 / 0) \\
Üretim (\$ mn) & 0 (−2 / +1) & 0 (−2 / +1) & 0 (0 / 0) \\
İstihdam (kişi) & +2 (−21 / +13) & +2 (−17 / +13) & 0 (−2 / +4) \\
İstihdam (\%) & 0,0 (−0,1 / +0,1) & 0,0 (−0,1 / +0,1) & 0,0 (0,0 / 0,0) \\
Cari denge (\$ mn) & 0 (−2 / +1) & 0 (−1 / +1) & 0 (0 / 0) \\
TÜFE (bp) & 0,00 (0,00 / 0,00) & 0,00 (0,00 / 0,00) & 0,00 (0,00 / 0,00) \\
\end{longtable}

\endgroup

\emph{Kaynak:} \etiketT{} f3e; toplamlar Faz 2B genel denge dağılımına kalibre. Yüzdeler 2024 düzeyine göre.

\hypertarget{nace_kartlar_16-18-c2-1621-ahux15fap-kaplama-paneli-ve-aux11fauxe7-esaslux131-panel-imalatux131}{%
\subsection{C.2 16.21 Ahşap kaplama paneli ve ağaç esaslı panel imalatı}\label{nace_kartlar_16-18-c2-1621-ahux15fap-kaplama-paneli-ve-aux11fauxe7-esaslux131-panel-imalatux131}}

\Needspace{7\baselineskip}

\textbf{Tablo C.2 --- 16.21: veri sayfası}

\nopagebreak[4]

\begingroup\gbtablo\small

\begin{longtable}[]{@{}
  >{\raggedright\arraybackslash}p{(\columnwidth - 4\tabcolsep) * \real{0.1350}}
  >{\raggedright\arraybackslash}p{(\columnwidth - 4\tabcolsep) * \real{0.3528}}
  >{\raggedright\arraybackslash}p{(\columnwidth - 4\tabcolsep) * \real{0.5123}}@{}}
\toprule\noalign{}
\begin{minipage}[b]{\linewidth}\raggedright
\textbf{Alan}
\end{minipage} & \begin{minipage}[b]{\linewidth}\raggedright
\textbf{Gösterge}
\end{minipage} & \begin{minipage}[b]{\linewidth}\raggedright
\textbf{Değer}
\end{minipage} \\
\midrule\noalign{}
\endhead
\bottomrule\noalign{}
\endlastfoot
1 Kimlik & Rev. 2.1 karşılığı; CPA & 16.21, 16.27 (1:2); CPA 16.21 \\
& Başlıca HS4 (ihracat, 2024) & 4411 \%72, 4410 \%20, 4412 \%5, 4408 \%3 \\
& BEC payları: ara / sermaye / tüketim & \%100 / \%0 / \%0 \\
2 Ölçek & Üretim; katma değer (GSYH payı) & \$4.974 mn; \$1.010 mn (\%0,07) \\
& Çalışan; girişim & 22.174; 417 \\
& Girişim ölçeği: 0 / 1--9 / 10--49 / 50--249 / 250+ çalışan & 23 / 280 / 83 / 39 / 15 \\
& 250+ girişimlerin istihdam payı; kişi başı KD (AB27 = 1) & \%69; 0,55 \\
& İstihdam değişimi, 2015--2024 & \%30 \\
& İlk üç bölge (İBBS-2) & Kocaeli, Sakarya, Düzce, Bolu, Yalova \%20; İstanbul \%17; Bursa, Eskişehir, Bilecik \%15 \\
& Çapa yatırımcılar & Kastamonu Entegre, AGT, Yıldız Entegre, Starwood ve Çamsan (MDF ve yonga levha); Kocaeli, Kastamonu, Antalya, Ordu ve İnegöl \etiketH{} \\
3 Dış ticaret & İhracat 2024; 2020--2024 değişimi & \$1.350 mn; \%83 \\
& AB payı; ilk beş AB pazarı & \%16; Romanya \%28; Bulgaristan \%21; Yunanistan \%13; İspanya \%7; İtalya \%6 \\
& İlk beş AB dışı pazar & Irak \%14; Fas \%12; Libya \%6; Gürcistan \%5; Kanada \%5 \\
& İthalat 2024; AB ve Çin payı & \$287 mn; AB \%36, Çin \%4 \\
& İlk beş ithalat kaynağı & Rusya \%39; Letonya \%8; İtalya \%5; Romanya \%5; Çin \%4 \\
& Dış ticaret dengesi & +\$1.062 mn \\
& İthal girdi içeriği (I-O); KD ÷ üretim & \%22; \%20 \\
& RCA: dünya 2020 → 2024; AB pazarı 2024 & 1,87 → 2,61; 1,96 \\
& Birim değer ÷ AB pazarı ortalaması & 0,96 \\
& AB pazarındaki pay (AB dışı tedarik içinde) 2020 → 2024; Türkiye'nin sırası & \%2,9 → \%7,8; 4. \\
4 AB tarafı & AB'ye ihracat 2025; rejim payları (GB / İTÜ / tarım) & €194,2 mn; \%100 / \%0 / \%0 \\
& S4'te AB tarifesi: 1/98 sürerse / askıda / tam MFN ortalaması & \%6,9 / \%6,9 / \%6,9 \\
& Tepe satırlar (S4 artışı ≥ \%10) & 11 satır, AB'ye ihracatın \%3'i: KN 4412 99 90 \%10,0 (€3,7 mn); KN 4412 41 99 \%10,0 (€0,9 mn); KN 4412 92 90 \%10,0 (€0,2 mn) \\
& AB pazarında AB dışı rakipler (AB tercihi) & Çin \%24 (MFN/GSP); Ukrayna \%15 (STA/EPA); Brezilya \%10 (MFN/GSP); İsviçre \%7 (AEA/STA); Gabon \%5 (MFN/GSP) \\
& Tercihli rakiplerin payı; AME (puan; ham) & \%41; 53 (1,43) \\
& Menşe kuralı (PEM) & 4410--4412: başlık değişikliği; ithal yonga, lif ve kaplamadan levha ve kontrplak menşe kazanır \\
& Menşe riski S2 / S3; c\_RoO (\%); p\_uyum & D/D; 1,0 / 1,5; 0,95 / 0,92 \etiketH{} \\
& Değer zincirindeki rol; DİR bağımlılığı & ara; düşük \\
& EUDR ve SKDM kapsamındaki AB'ye ihracat payı & \%100; \%0 \\
& PEM ortaklarına ihracat; kümülasyon maruziyeti; κ & \$438 mn; \$224 mn; 0,0 \etiketH{} \\
5 Türkiye tarafı & Uygulanan tarife 2026 (ithalat ağırlıklı): MFN; İGV dahil & \%5,4; \%14,4 \\
& S4'te AB menşeline tarife (c2) & \%15,1 \\
& Kritik girdiler: 2024 ithalatı, AB payı, ilk kaynaklar & HS4 4401, 4403, 4408, 3909: \$944 mn (AB \%44; Çin \%19, ABD \%9, Almanya \%9) \\
6 Kabiliyet & KS; sıra (ağırlık duyarlılığında P10--P90) & 50,8; 109 / 230 (58--153) \\
& KS bileşenleri (yüzdelik): RCA, PCI, yurt içi KD, Ar-Ge, tedarikçi, verimlilik, firma & 83, 35, 45, 25, 62, 71, 35 \\
& PCI (sınıf) & −0,25 \\
& Atlanabilir ürünler (yoğunluk üst çeyrek, PCI \textgreater{} sınıf) & yok \\
7 Stratejik önem & SS; en yüksek boyut; ağırlıklı ortalama & 0,0; --- 0; 0,0 \etiketH{} \\
8 Makro duyarlılık & TÜFE ağırlığı (doğrudan); tüketimde ithal pay & \%0,00; \%38 \\
& ÜFE vekili (imalat üretimindeki pay) & \%0,75 \\
& EDE, 1 puan tarife: tüm menşeler / yalnız AB (bp; ρ = 1) & 0,015 / 0,005 \\
& CKE: tarife (\$ mn / bp); mali (\$ mn / mlr TL); puan & 5.315,1; 7,85; 53 \\
9 Senaryolar & Senaryo tablosu & Tablo 5.1 \\
\end{longtable}

\endgroup

\hypertarget{nace_kartlar_16-18-c3-1622-birleux15ftirilmiux15f-parke-yer-duxf6ux15femelerinin-imalatux131}{%
\subsection{C.3 16.22 Birleştirilmiş parke yer döşemelerinin imalatı}\label{nace_kartlar_16-18-c3-1622-birleux15ftirilmiux15f-parke-yer-duxf6ux15femelerinin-imalatux131}}

\Needspace{7\baselineskip}

\textbf{Tablo C.3 --- 16.22: veri sayfası}

\nopagebreak[4]

\begingroup\gbtablo\small

\begin{longtable}[]{@{}
  >{\raggedright\arraybackslash}p{(\columnwidth - 4\tabcolsep) * \real{0.1357}}
  >{\raggedright\arraybackslash}p{(\columnwidth - 4\tabcolsep) * \real{0.3598}}
  >{\raggedright\arraybackslash}p{(\columnwidth - 4\tabcolsep) * \real{0.5045}}@{}}
\toprule\noalign{}
\begin{minipage}[b]{\linewidth}\raggedright
\textbf{Alan}
\end{minipage} & \begin{minipage}[b]{\linewidth}\raggedright
\textbf{Gösterge}
\end{minipage} & \begin{minipage}[b]{\linewidth}\raggedright
\textbf{Değer}
\end{minipage} \\
\midrule\noalign{}
\endhead
\bottomrule\noalign{}
\endlastfoot
1 Kimlik & Rev. 2.1 karşılığı; CPA & 16.22, 16.27 (1:2); CPA 16.22 \\
& Başlıca HS4 (ihracat, 2024) & 4418 \%100 \\
& BEC payları: ara / sermaye / tüketim & \%100 / \%0 / \%0 \\
2 Ölçek & Üretim; katma değer (GSYH payı) & \$464 mn; \$115 mn (\%0,01) \\
& Çalışan; girişim & 2.756; 49 \\
& Girişim ölçeği: 0 / 1--9 / 10--49 / 50--249 / 250+ çalışan & 6 / 24 / 15 / 3 / 4 \\
& 250+ girişimlerin istihdam payı; kişi başı KD (AB27 = 1) & \%70; 0,93 \\
& İstihdam değişimi, 2015--2024 & \%26 \\
& İlk üç bölge (İBBS-2) & Samsun, Tokat, Çorum, Amasya \%31; Kocaeli, Sakarya, Düzce, Bolu, Yalova \%31; Tekirdağ, Edirne, Kırklareli \%24 \\
& Çapa yatırımcılar & Masif ve çok katmanlı parke üreticileri; levha gruplarının döşeme hatları \etiketH{} \\
3 Dış ticaret & İhracat 2024; 2020--2024 değişimi & \$82 mn; \%70 \\
& AB payı; ilk beş AB pazarı & \%17; Bulgaristan \%15; Belçika \%14; Almanya \%13; Romanya \%11; Fransa \%10 \\
& İlk beş AB dışı pazar & ABD \%16; Irak \%12; Libya \%8; Azerbaycan \%6; Kanada \%4 \\
& İthalat 2024; AB ve Çin payı & \$26 mn; AB \%53, Çin \%17 \\
& İlk beş ithalat kaynağı & İtalya \%19; Çin \%17; Belarus \%10; Avusturya \%6; Polonya \%6 \\
& Dış ticaret dengesi & +\$56 mn \\
& İthal girdi içeriği (I-O); KD ÷ üretim & \%22; \%25 \\
& RCA: dünya 2020 → 2024; AB pazarı 2024 & 1,05 → 1,55; 0,67 \\
& Birim değer ÷ AB pazarı ortalaması & 1,16 \\
& AB pazarındaki pay (AB dışı tedarik içinde) 2020 → 2024; Türkiye'nin sırası & \%1,7 → \%2,7; 7. \\
4 AB tarafı & AB'ye ihracat 2025; rejim payları (GB / İTÜ / tarım) & €7,2 mn; \%100 / \%0 / \%0 \\
& S4'te AB tarifesi: 1/98 sürerse / askıda / tam MFN ortalaması & \%0,0 / \%0,0 / \%0,0 \\
& Tepe satırlar (S4 artışı ≥ \%10) & yok \\
& AB pazarında AB dışı rakipler (AB tercihi) & Çin \%52 (MFN/GSP); Ukrayna \%14 (STA/EPA); İsviçre \%8 (AEA/STA); Malezya \%5 (MFN/GSP); Endonezya \%4 (MFN/GSP) \\
& Tercihli rakiplerin payı; AME (puan; ham) & \%27; 19 (0,34) \\
& Menşe kuralı (PEM) & 4418 73--79: başlık değişikliği \\
& Menşe riski S2 / S3; c\_RoO (\%); p\_uyum & D/D; 1,0 / 1,5; 0,95 / 0,92 \etiketH{} \\
& Değer zincirindeki rol; DİR bağımlılığı & nihai; düşük \\
& EUDR ve SKDM kapsamındaki AB'ye ihracat payı & \%100; \%0 \\
& PEM ortaklarına ihracat; kümülasyon maruziyeti; κ & \$10 mn; \$5 mn; 0,0 \etiketH{} \\
5 Türkiye tarafı & Uygulanan tarife 2026 (ithalat ağırlıklı): MFN; İGV dahil & \%0,1; \%12,2 \\
& S4'te AB menşeline tarife (c2) & \%12,8 \\
& Kritik girdiler: 2024 ithalatı, AB payı, ilk kaynaklar & --- \\
6 Kabiliyet & KS; sıra (ağırlık duyarlılığında P10--P90) & 49,6; 121 / 230 (63--170) \\
& KS bileşenleri (yüzdelik): RCA, PCI, yurt içi KD, Ar-Ge, tedarikçi, verimlilik, firma & 70, 32, 53, 25, 62, 95, 17 \\
& PCI (sınıf) & −0,30 \\
& Atlanabilir ürünler (yoğunluk üst çeyrek, PCI \textgreater{} sınıf) & yok \\
7 Stratejik önem & SS; en yüksek boyut; ağırlıklı ortalama & 0,0; --- 0; 0,0 \etiketH{} \\
8 Makro duyarlılık & TÜFE ağırlığı (doğrudan); tüketimde ithal pay & \%0,00; \%53 \\
& ÜFE vekili (imalat üretimindeki pay) & \%0,07 \\
& EDE, 1 puan tarife: tüm menşeler / yalnız AB (bp; ρ = 1) & 0,001 / 0,001 \\
& CKE: tarife (\$ mn / bp); mali (\$ mn / mlr TL); puan & 4.007,1; 5,02; 46 \\
9 Senaryolar & Senaryo tablosu & Tablo C.3b (aşağıda) \\
\end{longtable}

\endgroup

\Needspace{7\baselineskip}

\textbf{Tablo C.3b --- 16.22 Ahşap döşeme (parke): senaryo etkileri (T+5'te yıllık; P50, parantez içinde P10 / P90)}

\nopagebreak[4]

\begingroup\gbtablo\small

\begin{longtable}[]{@{}lrrr@{}}
\toprule\noalign{}
\textbf{Değişken} & \textbf{S2} & \textbf{S3} & \textbf{S4} \\
\midrule\noalign{}
\endhead
\bottomrule\noalign{}
\endlastfoot
AB'ye ihracat (\$ mn) & 0 (−1 / 0) & −1 (−2 / 0) & −1 (−2 / 0) \\
AB'ye ihracat (\%) & −3,3 (−7,7 / +2,0) & −7,5 (−14,9 / −1,1) & −7,0 (−12,1 / −2,9) \\
Toplam ihracat (\$ mn) & 0 (−1 / 0) & −1 (−2 / 0) & −1 (−1 / 0) \\
Nihai mal ithalatı (\$ mn) & −1 (−2 / +1) & −1 (−3 / +1) & −2 (−4 / 0) \\
Üretim (\$ mn) & 0 (−1 / +2) & 0 (−1 / +1) & +2 (0 / +3) \\
İstihdam (kişi) & +2 (−3 / +7) & +1 (−3 / +5) & +6 (+1 / +12) \\
İstihdam (\%) & +0,1 (−0,1 / +0,2) & 0,0 (−0,1 / +0,2) & +0,2 (0,0 / +0,5) \\
Cari denge (\$ mn) & 0 (−1 / +1) & 0 (−1 / +1) & +1 (0 / +2) \\
TÜFE (bp) & 0,00 (0,00 / 0,00) & 0,00 (0,00 / 0,00) & +0,01 (+0,01 / +0,01) \\
\end{longtable}

\endgroup

\emph{Kaynak:} \etiketT{} f3e; toplamlar Faz 2B genel denge dağılımına kalibre. Yüzdeler 2024 düzeyine göre.

\hypertarget{nace_kartlar_16-18-c4-1623-diux11fer-bina-doux11framacux131lux131ux11fux131-ve-marangozluk-uxfcruxfcnlerinin-imalatux131}{%
\subsection{C.4 16.23 Diğer bina doğramacılığı ve marangozluk ürünlerinin imalatı}\label{nace_kartlar_16-18-c4-1623-diux11fer-bina-doux11framacux131lux131ux11fux131-ve-marangozluk-uxfcruxfcnlerinin-imalatux131}}

\Needspace{7\baselineskip}

\textbf{Tablo C.4 --- 16.23: veri sayfası}

\nopagebreak[4]

\begingroup\gbtablo\small

\begin{longtable}[]{@{}
  >{\raggedright\arraybackslash}p{(\columnwidth - 4\tabcolsep) * \real{0.1347}}
  >{\raggedright\arraybackslash}p{(\columnwidth - 4\tabcolsep) * \real{0.3505}}
  >{\raggedright\arraybackslash}p{(\columnwidth - 4\tabcolsep) * \real{0.5147}}@{}}
\toprule\noalign{}
\begin{minipage}[b]{\linewidth}\raggedright
\textbf{Alan}
\end{minipage} & \begin{minipage}[b]{\linewidth}\raggedright
\textbf{Gösterge}
\end{minipage} & \begin{minipage}[b]{\linewidth}\raggedright
\textbf{Değer}
\end{minipage} \\
\midrule\noalign{}
\endhead
\bottomrule\noalign{}
\endlastfoot
1 Kimlik & Rev. 2.1 karşılığı; CPA & 16.23, 16.25, 16.27 (1:3); CPA 16.23 \\
& Başlıca HS4 (ihracat, 2024) & 4418 \%72, 9406 \%28 \\
& BEC payları: ara / sermaye / tüketim & \%100 / \%0 / \%0 \\
2 Ölçek & Üretim; katma değer (GSYH payı) & \$963 mn; \$248 mn (\%0,02) \\
& Çalışan; girişim & 30.303; 11.871 \\
& Girişim ölçeği: 0 / 1--9 / 10--49 / 50--249 / 250+ çalışan & 144 / 12.505 / 240 / --- / --- \\
& 250+ girişimlerin istihdam payı; kişi başı KD (AB27 = 1) & \%1; 0,17 \\
& İstihdam değişimi, 2015--2024 & −\%22 \\
& İlk üç bölge (İBBS-2) & Aydın, Denizli, Muğla \%9; Kastamonu, Çankırı, Sinop \%8; İstanbul \%7 \\
& Çapa yatırımcılar & Ahşap kapı, pencere ve prefabrik ahşap yapı üreticileri; çoğu on kişiden az çalışanlı atölyeler \etiketH{} \\
3 Dış ticaret & İhracat 2024; 2020--2024 değişimi & \$227 mn; \%31 \\
& AB payı; ilk beş AB pazarı & \%22; Bulgaristan \%17; Romanya \%13; Almanya \%12; Fransa \%12; Hollanda \%9 \\
& İlk beş AB dışı pazar & Irak \%13; Libya \%8; Suudi Arabistan \%6; Kazakistan \%5; ABD \%5 \\
& İthalat 2024; AB ve Çin payı & \$49 mn; AB \%66, Çin \%4 \\
& İlk beş ithalat kaynağı & Avusturya \%34; Rusya \%22; Almanya \%13; İtalya \%6; Çin \%4 \\
& Dış ticaret dengesi & +\$178 mn \\
& İthal girdi içeriği (I-O); KD ÷ üretim & \%22; \%26 \\
& RCA: dünya 2020 → 2024; AB pazarı 2024 & 1,19 → 1,19; 1,38 \\
& Birim değer ÷ AB pazarı ortalaması & 1,04 \\
& AB pazarındaki pay (AB dışı tedarik içinde) 2020 → 2024; Türkiye'nin sırası & \%3,5 → \%5,5; 7. \\
4 AB tarafı & AB'ye ihracat 2025; rejim payları (GB / İTÜ / tarım) & €56,7 mn; \%100 / \%0 / \%0 \\
& S4'te AB tarifesi: 1/98 sürerse / askıda / tam MFN ortalaması & \%1,2 / \%1,2 / \%1,2 \\
& Tepe satırlar (S4 artışı ≥ \%10) & yok \\
& AB pazarında AB dışı rakipler (AB tercihi) & Çin \%23 (MFN/GSP); Birleşik Krallık \%12 (TCA); Endonezya \%9 (MFN/GSP); Ukrayna \%8 (STA/EPA); Bosna-Hersek \%7 (MFN/GSP) \\
& Tercihli rakiplerin payı; AME (puan; ham) & \%35; 27 (0,69) \\
& Menşe kuralı (PEM) & 4418: başlık değişikliği; hücreli ahşap panel, şindıl ve kiremit kullanılabilir \\
& Menşe riski S2 / S3; c\_RoO (\%); p\_uyum & D/D; 1,0 / 1,5; 0,95 / 0,92 \etiketH{} \\
& Değer zincirindeki rol; DİR bağımlılığı & nihai; düşük \\
& EUDR ve SKDM kapsamındaki AB'ye ihracat payı & \%67; \%0 \\
& PEM ortaklarına ihracat; kümülasyon maruziyeti; κ & \$28 mn; \$15 mn; 0,0 \etiketH{} \\
5 Türkiye tarafı & Uygulanan tarife 2026 (ithalat ağırlıklı): MFN; İGV dahil & \%0,3; \%2,0 \\
& S4'te AB menşeline tarife (c2) & \%1,0 \\
& Kritik girdiler: 2024 ithalatı, AB payı, ilk kaynaklar & --- \\
6 Kabiliyet & KS; sıra (ağırlık duyarlılığında P10--P90) & 43,4; 178 / 230 (129--201) \\
& KS bileşenleri (yüzdelik): RCA, PCI, yurt içi KD, Ar-Ge, tedarikçi, verimlilik, firma & 45, 41, 56, 25, 62, 2, 77 \\
& PCI (sınıf) & −0,16 \\
& Atlanabilir ürünler (yoğunluk üst çeyrek, PCI \textgreater{} sınıf) & 940610 (PCI +0,09, RCA 0,23) \\
7 Stratejik önem & SS; en yüksek boyut; ağırlıklı ortalama & 0,0; --- 0; 0,0 \etiketH{} \\
8 Makro duyarlılık & TÜFE ağırlığı (doğrudan); tüketimde ithal pay & \%0,00; \%38 \\
& ÜFE vekili (imalat üretimindeki pay) & \%0,14 \\
& EDE, 1 puan tarife: tüm menşeler / yalnız AB (bp; ρ = 1) & 0,003 / 0,002 \\
& CKE: tarife (\$ mn / bp); mali (\$ mn / mlr TL); puan & 5.019,9; 6,07; 50 \\
9 Senaryolar & Senaryo tablosu & Tablo C.4b (aşağıda) \\
\end{longtable}

\endgroup

\Needspace{7\baselineskip}

\textbf{Tablo C.4b --- 16.23 Doğrama ve ahşap yapı: senaryo etkileri (T+5'te yıllık; P50, parantez içinde P10 / P90)}

\nopagebreak[4]

\begingroup\gbtablo\small

\begin{longtable}[]{@{}lrrr@{}}
\toprule\noalign{}
\textbf{Değişken} & \textbf{S2} & \textbf{S3} & \textbf{S4} \\
\midrule\noalign{}
\endhead
\bottomrule\noalign{}
\endlastfoot
AB'ye ihracat (\$ mn) & −4 (−9 / +1) & −6 (−11 / −1) & −5 (−9 / −2) \\
AB'ye ihracat (\%) & −8,1 (−18,2 / +2,0) & −11,6 (−22,5 / −1,8) & −11,0 (−18,3 / −4,7) \\
Toplam ihracat (\$ mn) & −3 (−7 / +1) & −4 (−8 / −1) & −4 (−7 / −2) \\
Nihai mal ithalatı (\$ mn) & −1 (−4 / +2) & −2 (−5 / +1) & −2 (−3 / 0) \\
Üretim (\$ mn) & −2 (−3 / −1) & −2 (−4 / −2) & −3 (−5 / −2) \\
İstihdam (kişi) & −37 (−65 / −25) & −51 (−83 / −36) & −61 (−101 / −33) \\
İstihdam (\%) & −0,1 (−0,2 / −0,1) & −0,2 (−0,3 / −0,1) & −0,2 (−0,3 / −0,1) \\
Cari denge (\$ mn) & −1 (−2 / −1) & −2 (−3 / −1) & −2 (−4 / −1) \\
TÜFE (bp) & 0,00 (0,00 / 0,00) & 0,00 (0,00 / 0,00) & 0,00 (0,00 / 0,00) \\
\end{longtable}

\endgroup

\emph{Kaynak:} \etiketT{} f3e; toplamlar Faz 2B genel denge dağılımına kalibre. Yüzdeler 2024 düzeyine göre.

\hypertarget{nace_kartlar_16-18-c5-1624-ahux15fap-konteyner-imalatux131}{%
\subsection{C.5 16.24 Ahşap konteyner imalatı}\label{nace_kartlar_16-18-c5-1624-ahux15fap-konteyner-imalatux131}}

\Needspace{7\baselineskip}

\textbf{Tablo C.5 --- 16.24: veri sayfası}

\nopagebreak[4]

\begingroup\gbtablo\small

\begin{longtable}[]{@{}
  >{\raggedright\arraybackslash}p{(\columnwidth - 4\tabcolsep) * \real{0.1357}}
  >{\raggedright\arraybackslash}p{(\columnwidth - 4\tabcolsep) * \real{0.3597}}
  >{\raggedright\arraybackslash}p{(\columnwidth - 4\tabcolsep) * \real{0.5046}}@{}}
\toprule\noalign{}
\begin{minipage}[b]{\linewidth}\raggedright
\textbf{Alan}
\end{minipage} & \begin{minipage}[b]{\linewidth}\raggedright
\textbf{Gösterge}
\end{minipage} & \begin{minipage}[b]{\linewidth}\raggedright
\textbf{Değer}
\end{minipage} \\
\midrule\noalign{}
\endhead
\bottomrule\noalign{}
\endlastfoot
1 Kimlik & Rev. 2.1 karşılığı; CPA & 16.24, 16.27 (1:2); CPA 16.24 \\
& Başlıca HS4 (ihracat, 2024) & 4415 \%98, 4416 \%2 \\
& BEC payları: ara / sermaye / tüketim & \%100 / \%0 / \%0 \\
2 Ölçek & Üretim; katma değer (GSYH payı) & \$1.447 mn; \$266 mn (\%0,02) \\
& Çalışan; girişim & 14.721; 2.165 \\
& Girişim ölçeği: 0 / 1--9 / 10--49 / 50--249 / 250+ çalışan & 114 / 1.854 / 308 / 37 / 14 \\
& 250+ girişimlerin istihdam payı; kişi başı KD (AB27 = 1) & \%24; 0,31 \\
& İstihdam değişimi, 2015--2024 & \%30 \\
& İlk üç bölge (İBBS-2) & Kocaeli, Sakarya, Düzce, Bolu, Yalova \%20; İzmir \%14; Adana, Mersin \%9 \\
& Çapa yatırımcılar & Lojistik merkezlerine yakın palet üreticileri (Kocaeli, İstanbul, Bursa, Mersin) \etiketH{} \\
3 Dış ticaret & İhracat 2024; 2020--2024 değişimi & \$29 mn; \%71 \\
& AB payı; ilk beş AB pazarı & \%45; Almanya \%46; Bulgaristan \%14; İtalya \%6; Portekiz \%5; Fransa \%4 \\
& İlk beş AB dışı pazar & Ürdün \%21; Irak \%17; Mısır \%7; Fas \%6; Türkmenistan \%5 \\
& İthalat 2024; AB ve Çin payı & \$28 mn; AB \%62, Çin \%1 \\
& İlk beş ithalat kaynağı & Ukrayna \%31; Almanya \%13; İsveç \%11; Finlandiya \%9; Fransa \%6 \\
& Dış ticaret dengesi & +\$1 mn \\
& İthal girdi içeriği (I-O); KD ÷ üretim & \%22; \%18 \\
& RCA: dünya 2020 → 2024; AB pazarı 2024 & 0,36 → 0,39; 0,53 \\
& Birim değer ÷ AB pazarı ortalaması & 1,13 \\
& AB pazarındaki pay (AB dışı tedarik içinde) 2020 → 2024; Türkiye'nin sırası & \%1,7 → \%2,1; 9. \\
4 AB tarafı & AB'ye ihracat 2025; rejim payları (GB / İTÜ / tarım) & €12,2 mn; \%100 / \%0 / \%0 \\
& S4'te AB tarifesi: 1/98 sürerse / askıda / tam MFN ortalaması & \%3,3 / \%3,3 / \%3,3 \\
& Tepe satırlar (S4 artışı ≥ \%10) & yok \\
& AB pazarında AB dışı rakipler (AB tercihi) & ABD \%26 (MFN/GSP); Ukrayna \%19 (STA/EPA); Birleşik Krallık \%18 (TCA); Norveç \%8 (AEA/STA); İsviçre \%5 (AEA/STA) \\
& Tercihli rakiplerin payı; AME (puan; ham) & \%56; 70 (2,38) \\
& Menşe kuralı (PEM) & 4415--4416: başlık değişikliği \\
& Menşe riski S2 / S3; c\_RoO (\%); p\_uyum & D/D; 1,0 / 1,5; 0,95 / 0,92 \etiketH{} \\
& Değer zincirindeki rol; DİR bağımlılığı & ara; düşük \\
& EUDR ve SKDM kapsamındaki AB'ye ihracat payı & \%100; \%0 \\
& PEM ortaklarına ihracat; kümülasyon maruziyeti; κ & \$4 mn; \$2 mn; 0,0 \etiketH{} \\
5 Türkiye tarafı & Uygulanan tarife 2026 (ithalat ağırlıklı): MFN; İGV dahil & \%3,2; \%26,1 \\
& S4'te AB menşeline tarife (c2) & \%24,8 \\
& Kritik girdiler: 2024 ithalatı, AB payı, ilk kaynaklar & --- \\
6 Kabiliyet & KS; sıra (ağırlık duyarlılığında P10--P90) & 42,5; 185 / 230 (137--206) \\
& KS bileşenleri (yüzdelik): RCA, PCI, yurt içi KD, Ar-Ge, tedarikçi, verimlilik, firma & 29, 59, 43, 25, 62, 18, 67 \\
& PCI (sınıf) & 0,16 \\
& Atlanabilir ürünler (yoğunluk üst çeyrek, PCI \textgreater{} sınıf) & yok \\
7 Stratejik önem & SS; en yüksek boyut; ağırlıklı ortalama & 7,5; gıda güvenliği 10; 1,5 \etiketH{} \\
8 Makro duyarlılık & TÜFE ağırlığı (doğrudan); tüketimde ithal pay & \%0,00; \%95 \\
& ÜFE vekili (imalat üretimindeki pay) & \%0,22 \\
& EDE, 1 puan tarife: tüm menşeler / yalnız AB (bp; ρ = 1) & 0,001 / 0,001 \\
& CKE: tarife (\$ mn / bp); mali (\$ mn / mlr TL); puan & 6.138,1; 2,27; 48 \\
9 Senaryolar & Senaryo tablosu & Tablo C.5b (aşağıda) \\
\end{longtable}

\endgroup

\Needspace{7\baselineskip}

\textbf{Tablo C.5b --- 16.24 Palet ve ahşap ambalaj: senaryo etkileri (T+5'te yıllık; P50, parantez içinde P10 / P90)}

\nopagebreak[4]

\begingroup\gbtablo\small

\begin{longtable}[]{@{}lrrr@{}}
\toprule\noalign{}
\textbf{Değişken} & \textbf{S2} & \textbf{S3} & \textbf{S4} \\
\midrule\noalign{}
\endhead
\bottomrule\noalign{}
\endlastfoot
AB'ye ihracat (\$ mn) & −1 (−3 / 0) & −2 (−3 / 0) & −2 (−4 / −1) \\
AB'ye ihracat (\%) & −9,6 (−22,2 / +2,0) & −13,1 (−25,9 / −2,1) & −18,2 (−29,2 / −7,4) \\
Toplam ihracat (\$ mn) & −1 (−2 / 0) & −1 (−2 / 0) & −2 (−3 / −1) \\
Nihai mal ithalatı (\$ mn) & −2 (−6 / +1) & −3 (−6 / +1) & −5 (−8 / −1) \\
Üretim (\$ mn) & +1 (−1 / +4) & +1 (−1 / +4) & +3 (0 / +6) \\
İstihdam (kişi) & +7 (−7 / +29) & +9 (−7 / +30) & +20 (+1 / +40) \\
İstihdam (\%) & 0,0 (0,0 / +0,2) & +0,1 (0,0 / +0,2) & +0,1 (0,0 / +0,3) \\
Cari denge (\$ mn) & +1 (−1 / +3) & +1 (−1 / +3) & +2 (0 / +4) \\
TÜFE (bp) & 0,00 (0,00 / 0,00) & 0,00 (0,00 / 0,00) & +0,02 (+0,02 / +0,02) \\
\end{longtable}

\endgroup

\emph{Kaynak:} \etiketT{} f3e; toplamlar Faz 2B genel denge dağılımına kalibre. Yüzdeler 2024 düzeyine göre.

\hypertarget{nace_kartlar_16-18-c6-1629-aux11fauxe7tan-diux11fer-uxfcruxfcnlerin-imalatux131-mantar-saz-saman-ve-benzeri-uxf6rme-malzemelerinden-uxfcruxfcnlerin-imalatux131}{%
\subsection{C.6 16.29 Ağaçtan diğer ürünlerin imalatı; mantar, saz, saman ve benzeri örme malzemelerinden ürünlerin imalatı}\label{nace_kartlar_16-18-c6-1629-aux11fauxe7tan-diux11fer-uxfcruxfcnlerin-imalatux131-mantar-saz-saman-ve-benzeri-uxf6rme-malzemelerinden-uxfcruxfcnlerin-imalatux131}}

\Needspace{7\baselineskip}

\textbf{Tablo C.6 --- 16.29: veri sayfası}

\nopagebreak[4]

\begingroup\gbtablo\small

\begin{longtable}[]{@{}
  >{\raggedright\arraybackslash}p{(\columnwidth - 4\tabcolsep) * \real{0.1361}}
  >{\raggedright\arraybackslash}p{(\columnwidth - 4\tabcolsep) * \real{0.3642}}
  >{\raggedright\arraybackslash}p{(\columnwidth - 4\tabcolsep) * \real{0.4997}}@{}}
\toprule\noalign{}
\begin{minipage}[b]{\linewidth}\raggedright
\textbf{Alan}
\end{minipage} & \begin{minipage}[b]{\linewidth}\raggedright
\textbf{Gösterge}
\end{minipage} & \begin{minipage}[b]{\linewidth}\raggedright
\textbf{Değer}
\end{minipage} \\
\midrule\noalign{}
\endhead
\bottomrule\noalign{}
\endlastfoot
1 Kimlik & Rev. 2.1 karşılığı; CPA & 15.20, 16.26, 16.27, 16.28 (1:4); CPA 16.29 \\
& Başlıca HS4 (ihracat, 2024) & 4421 \%44, 4419 \%15, 4401 \%12, 4504 \%8 \\
& BEC payları: ara / sermaye / tüketim & \%52 / \%0 / \%48 \\
2 Ölçek & Üretim; katma değer (GSYH payı) & \$175 mn; \$46 mn (\%0,00) \\
& Çalışan; girişim & 6.039; 2.667 \\
& Girişim ölçeği: 0 / 1--9 / 10--49 / 50--249 / 250+ çalışan & 109 / 2.840 / 63 / 8 / 0 \\
& 250+ girişimlerin istihdam payı; kişi başı KD (AB27 = 1) & \%0; 0,20 \\
& İstihdam değişimi, 2015--2024 & \%85 \\
& İlk üç bölge (İBBS-2) & Kocaeli, Sakarya, Düzce, Bolu, Yalova \%17; İstanbul \%9; Kastamonu, Çankırı, Sinop \%8 \\
& Çapa yatırımcılar & --- \\
3 Dış ticaret & İhracat 2024; 2020--2024 değişimi & \$89 mn; \%142 \\
& AB payı; ilk beş AB pazarı & \%41; Almanya \%21; Fransa \%14; Bulgaristan \%11; Danimarka \%10; Slovenya \%8 \\
& İlk beş AB dışı pazar & ABD \%21; Rusya \%10; Birleşik Krallık \%7; Irak \%7; Suudi Arabistan \%6 \\
& İthalat 2024; AB ve Çin payı & \$116 mn; AB \%31, Çin \%51 \\
& İlk beş ithalat kaynağı & Çin \%51; Portekiz \%9; İspanya \%4; Romanya \%4; Rusya \%3 \\
& Dış ticaret dengesi & −\$28 mn \\
& İthal girdi içeriği (I-O); KD ÷ üretim & \%22; \%26 \\
& RCA: dünya 2020 → 2024; AB pazarı 2024 & 0,16 → 0,30; 0,22 \\
& Birim değer ÷ AB pazarı ortalaması & 1,83 \\
& AB pazarındaki pay (AB dışı tedarik içinde) 2020 → 2024; Türkiye'nin sırası & \%0,4 → \%0,9; 16. \\
4 AB tarafı & AB'ye ihracat 2025; rejim payları (GB / İTÜ / tarım) & €26,4 mn; \%100 / \%0 / \%0 \\
& S4'te AB tarifesi: 1/98 sürerse / askıda / tam MFN ortalaması & \%0,7 / \%0,7 / \%0,7 \\
& Tepe satırlar (S4 artışı ≥ \%10) & yok \\
& AB pazarında AB dışı rakipler (AB tercihi) & Çin \%51 (MFN/GSP); ABD \%10 (MFN/GSP); Vietnam \%8 (STA/EPA); Ukrayna \%4 (STA/EPA); Hindistan \%3 (MFN/GSP) \\
& Tercihli rakiplerin payı; AME (puan; ham) & \%24; 40 (1,09) \\
& Menşe kuralı (PEM) & 4414, 4417, 4419--4421 ve Fasıl 45--46: başlık değişikliği \\
& Menşe riski S2 / S3; c\_RoO (\%); p\_uyum & D/D; 1,0 / 1,5; 0,95 / 0,92 \etiketH{} \\
& Değer zincirindeki rol; DİR bağımlılığı & nihai; düşük \\
& EUDR ve SKDM kapsamındaki AB'ye ihracat payı & \%90; \%0 \\
& PEM ortaklarına ihracat; kümülasyon maruziyeti; κ & \$6 mn; \$3 mn; 0,0 \etiketH{} \\
5 Türkiye tarafı & Uygulanan tarife 2026 (ithalat ağırlıklı): MFN; İGV dahil & \%1,9; \%4,4 \\
& S4'te AB menşeline tarife (c2) & \%4,8 \\
& Kritik girdiler: 2024 ithalatı, AB payı, ilk kaynaklar & --- \\
6 Kabiliyet & KS; sıra (ağırlık duyarlılığında P10--P90) & 37,9; 213 / 230 (179--219) \\
& KS bileşenleri (yüzdelik): RCA, PCI, yurt içi KD, Ar-Ge, tedarikçi, verimlilik, firma & 39, 23, 57, 25, 62, 4, 68 \\
& PCI (sınıf) & −0,57 \\
& Atlanabilir ürünler (yoğunluk üst çeyrek, PCI \textgreater{} sınıf) & 442199 (PCI −0,44, RCA 0,46) \\
7 Stratejik önem & SS; en yüksek boyut; ağırlıklı ortalama & 0,0; --- 0; 0,0 \etiketH{} \\
8 Makro duyarlılık & TÜFE ağırlığı (doğrudan); tüketimde ithal pay & \%0,00; \%95 \\
& ÜFE vekili (imalat üretimindeki pay) & \%0,03 \\
& EDE, 1 puan tarife: tüm menşeler / yalnız AB (bp; ρ = 1) & 0,006 / 0,002 \\
& CKE: tarife (\$ mn / bp); mali (\$ mn / mlr TL); puan & 2.036,4; 42,55; 65 \\
9 Senaryolar & Senaryo tablosu & Tablo C.6b (aşağıda) \\
\end{longtable}

\endgroup

\Needspace{7\baselineskip}

\textbf{Tablo C.6b --- 16.29 Diğer ağaç ürünleri: senaryo etkileri (T+5'te yıllık; P50, parantez içinde P10 / P90)}

\nopagebreak[4]

\begingroup\gbtablo\small

\begin{longtable}[]{@{}lrrr@{}}
\toprule\noalign{}
\textbf{Değişken} & \textbf{S2} & \textbf{S3} & \textbf{S4} \\
\midrule\noalign{}
\endhead
\bottomrule\noalign{}
\endlastfoot
AB'ye ihracat (\$ mn) & −2 (−4 / +1) & −2 (−4 / 0) & −2 (−3 / −1) \\
AB'ye ihracat (\%) & −4,4 (−10,2 / +2,0) & −5,6 (−11,0 / −0,9) & −5,2 (−9,1 / −2,3) \\
Toplam ihracat (\$ mn) & −1 (−3 / 0) & −1 (−3 / 0) & −2 (−3 / −1) \\
Nihai mal ithalatı (\$ mn) & −1 (−2 / +1) & −1 (−3 / +1) & −2 (−3 / 0) \\
Üretim (\$ mn) & −1 (−1 / 0) & −1 (−1 / 0) & 0 (0 / +1) \\
İstihdam (kişi) & −13 (−22 / −4) & −13 (−24 / −1) & +4 (−8 / +18) \\
İstihdam (\%) & −0,2 (−0,4 / −0,1) & −0,2 (−0,4 / 0,0) & +0,1 (−0,1 / +0,3) \\
Cari denge (\$ mn) & 0 (−1 / 0) & 0 (−1 / 0) & 0 (0 / +1) \\
TÜFE (bp) & 0,00 (0,00 / 0,00) & 0,00 (0,00 / 0,00) & +0,01 (+0,01 / +0,01) \\
\end{longtable}

\endgroup

\emph{Kaynak:} \etiketT{} f3e; toplamlar Faz 2B genel denge dağılımına kalibre. Yüzdeler 2024 düzeyine göre.

\hypertarget{nace_kartlar_16-18-c7-1711-kaux11fux131t-hamuru-imalatux131}{%
\subsection{C.7 17.11 Kağıt hamuru imalatı}\label{nace_kartlar_16-18-c7-1711-kaux11fux131t-hamuru-imalatux131}}

\Needspace{7\baselineskip}

\textbf{Tablo C.7 --- 17.11: veri sayfası}

\nopagebreak[4]

\begingroup\gbtablo\small

\begin{longtable}[]{@{}
  >{\raggedright\arraybackslash}p{(\columnwidth - 4\tabcolsep) * \real{0.1362}}
  >{\raggedright\arraybackslash}p{(\columnwidth - 4\tabcolsep) * \real{0.3653}}
  >{\raggedright\arraybackslash}p{(\columnwidth - 4\tabcolsep) * \real{0.4984}}@{}}
\toprule\noalign{}
\begin{minipage}[b]{\linewidth}\raggedright
\textbf{Alan}
\end{minipage} & \begin{minipage}[b]{\linewidth}\raggedright
\textbf{Gösterge}
\end{minipage} & \begin{minipage}[b]{\linewidth}\raggedright
\textbf{Değer}
\end{minipage} \\
\midrule\noalign{}
\endhead
\bottomrule\noalign{}
\endlastfoot
1 Kimlik & Rev. 2.1 karşılığı; CPA & 17.11 (1:1); CPA 17.11 \\
& Başlıca HS4 (ihracat, 2024) & 4703 \%98, 4706 \%1, 4704 \%1, 4702 \%0 \\
& BEC payları: ara / sermaye / tüketim & \%100 / \%0 / \%0 \\
2 Ölçek & Üretim; katma değer (GSYH payı) & \$--- mn; \$--- mn (---) \\
& Çalışan; girişim & ---; --- \\
& Girişim ölçeği: 0 / 1--9 / 10--49 / 50--249 / 250+ çalışan & --- / --- / --- / --- / --- \\
& 250+ girişimlerin istihdam payı; kişi başı KD (AB27 = 1) & ---; --- \\
& İstihdam değişimi, 2015--2024 & --- \\
& İlk üç bölge (İBBS-2) & --- \\
& Çapa yatırımcılar & Yerli kâğıt hamuru üretimi yok denecek kadar az; ithalatçılar entegre hijyen ve kâğıt üreticileri \etiketH{} \\
3 Dış ticaret & İhracat 2024; 2020--2024 değişimi & \$117 mn; \%115 \\
& AB payı; ilk beş AB pazarı & \%0; İtalya \%40; Hollanda \%20; Almanya \%19; Finlandiya \%8; Romanya \%4 \\
& İlk beş AB dışı pazar & İran \%77; Rusya \%10; Moğolistan \%3; Afganistan \%3; Özbekistan \%3 \\
& İthalat 2024; AB ve Çin payı & \$1.131 mn; AB \%39, Çin \%0 \\
& İlk beş ithalat kaynağı & Brezilya \%25; ABD \%17; Finlandiya \%16; Uruguay \%11; Portekiz \%9 \\
& Dış ticaret dengesi & −\$1.014 mn \\
& İthal girdi içeriği (I-O); KD ÷ üretim & \%28; --- \\
& RCA: dünya 2020 → 2024; AB pazarı 2024 & 0,13 → 0,19; 0,00 \\
& Birim değer ÷ AB pazarı ortalaması & 2,44 \\
& AB pazarındaki pay (AB dışı tedarik içinde) 2020 → 2024; Türkiye'nin sırası & \%0,0 → \%0,0; 23. \\
4 AB tarafı & AB'ye ihracat 2025; rejim payları (GB / İTÜ / tarım) & €0,2 mn; \%100 / \%0 / \%0 \\
& S4'te AB tarifesi: 1/98 sürerse / askıda / tam MFN ortalaması & \%0,0 / \%0,0 / \%0,0 \\
& Tepe satırlar (S4 artışı ≥ \%10) & yok \\
& AB pazarında AB dışı rakipler (AB tercihi) & Brezilya \%49 (MFN/GSP); ABD \%19 (MFN/GSP); Uruguay \%18 (MFN/GSP); Şili \%6 (STA/EPA); Norveç \%4 (AEA/STA) \\
& Tercihli rakiplerin payı; AME (puan; ham) & \%14; 10 (0,01) \\
& Menşe kuralı (PEM) & Fasıl 47: başlık değişikliği \\
& Menşe riski S2 / S3; c\_RoO (\%); p\_uyum & D/D; 1,0 / 1,5; 0,95 / 0,92 \etiketH{} \\
& Değer zincirindeki rol; DİR bağımlılığı & ara; düşük \\
& EUDR ve SKDM kapsamındaki AB'ye ihracat payı & \%100; \%0 \\
& PEM ortaklarına ihracat; kümülasyon maruziyeti; κ & \$1 mn; \$0 mn; 0,0 \etiketH{} \\
5 Türkiye tarafı & Uygulanan tarife 2026 (ithalat ağırlıklı): MFN; İGV dahil & \%0,0; \%0,0 \\
& S4'te AB menşeline tarife (c2) & \%0,0 \\
& Kritik girdiler: 2024 ithalatı, AB payı, ilk kaynaklar & HS4 4401, 4707: \$354 mn (AB \%27; Birleşik Krallık \%25, Irak \%17, ABD \%16) \\
6 Kabiliyet & KS; sıra (ağırlık duyarlılığında P10--P90) & 47,9; 138 / 230 (89--183) \\
& KS bileşenleri (yüzdelik): RCA, PCI, yurt içi KD, Ar-Ge, tedarikçi, verimlilik, firma & 33, 64, 51, 25, 75, ---, --- \\
& PCI (sınıf) & 0,24 \\
& Atlanabilir ürünler (yoğunluk üst çeyrek, PCI \textgreater{} sınıf) & 470500 (PCI +0,34, RCA 0,00) \\
7 Stratejik önem & SS; en yüksek boyut; ağırlıklı ortalama & 30,7; kritik hammadde 40; 9,0 \etiketH{} \\
8 Makro duyarlılık & TÜFE ağırlığı (doğrudan); tüketimde ithal pay & \%0,00; \%5 \\
& ÜFE vekili (imalat üretimindeki pay) & --- \\
& EDE, 1 puan tarife: tüm menşeler / yalnız AB (bp; ρ = 1) & 0,092 / 0,036 \\
& CKE: tarife (\$ mn / bp); mali (\$ mn / mlr TL); puan & 0,0; ---; 3 \\
9 Senaryolar & Senaryo tablosu & Tablo C.7b (aşağıda) \\
\end{longtable}

\endgroup

\Needspace{7\baselineskip}

\textbf{Tablo C.7b --- 17.11 Kâğıt hamuru: senaryo etkileri (T+5'te yıllık; P50, parantez içinde P10 / P90)}

\nopagebreak[4]

\begingroup\gbtablo\small

\begin{longtable}[]{@{}lrrr@{}}
\toprule\noalign{}
\textbf{Değişken} & \textbf{S2} & \textbf{S3} & \textbf{S4} \\
\midrule\noalign{}
\endhead
\bottomrule\noalign{}
\endlastfoot
AB'ye ihracat (\$ mn) & 0 (0 / 0) & 0 (0 / 0) & 0 (0 / 0) \\
AB'ye ihracat (\%) & −4,8 (−12,0 / +2,0) & −8,2 (−17,3 / −1,3) & −7,6 (−14,2 / −3,3) \\
Toplam ihracat (\$ mn) & 0 (0 / 0) & 0 (0 / 0) & 0 (0 / 0) \\
Nihai mal ithalatı (\$ mn) & 0 (0 / 0) & 0 (0 / 0) & 0 (0 / 0) \\
Üretim (\$ mn) & 0 (0 / 0) & 0 (0 / 0) & 0 (0 / 0) \\
İstihdam (kişi) & 0 (0 / 0) & 0 (0 / 0) & 0 (0 / 0) \\
İstihdam (\%) & --- & --- & --- \\
Cari denge (\$ mn) & 0 (0 / 0) & 0 (0 / 0) & 0 (0 / 0) \\
TÜFE (bp) & +0,02 (+0,01 / +0,02) & +0,03 (+0,02 / +0,05) & +0,03 (+0,02 / +0,05) \\
\end{longtable}

\endgroup

\emph{Kaynak:} \etiketT{} f3e; toplamlar Faz 2B genel denge dağılımına kalibre. Yüzdeler 2024 düzeyine göre.

\hypertarget{nace_kartlar_16-18-c8-1712-kaux11fux131t-ve-mukavva-imalatux131}{%
\subsection{C.8 17.12 Kağıt ve mukavva imalatı}\label{nace_kartlar_16-18-c8-1712-kaux11fux131t-ve-mukavva-imalatux131}}

\Needspace{7\baselineskip}

\textbf{Tablo C.8 --- 17.12: veri sayfası}

\nopagebreak[4]

\begingroup\gbtablo\small

\begin{longtable}[]{@{}
  >{\raggedright\arraybackslash}p{(\columnwidth - 4\tabcolsep) * \real{0.1340}}
  >{\raggedright\arraybackslash}p{(\columnwidth - 4\tabcolsep) * \real{0.3430}}
  >{\raggedright\arraybackslash}p{(\columnwidth - 4\tabcolsep) * \real{0.5230}}@{}}
\toprule\noalign{}
\begin{minipage}[b]{\linewidth}\raggedright
\textbf{Alan}
\end{minipage} & \begin{minipage}[b]{\linewidth}\raggedright
\textbf{Gösterge}
\end{minipage} & \begin{minipage}[b]{\linewidth}\raggedright
\textbf{Değer}
\end{minipage} \\
\midrule\noalign{}
\endhead
\bottomrule\noalign{}
\endlastfoot
1 Kimlik & Rev. 2.1 karşılığı; CPA & 17.12 (1:1); CPA 17.12 \\
& Başlıca HS4 (ihracat, 2024) & 4803 \%45, 4811 \%25, 4805 \%12, 4810 \%6 \\
& BEC payları: ara / sermaye / tüketim & \%100 / \%0 / \%0 \\
2 Ölçek & Üretim; katma değer (GSYH payı) & \$3.106 mn; \$759 mn (\%0,06) \\
& Çalışan; girişim & 13.598; 441 \\
& Girişim ölçeği: 0 / 1--9 / 10--49 / 50--249 / 250+ çalışan & 43 / 323 / 63 / 27 / 13 \\
& 250+ girişimlerin istihdam payı; kişi başı KD (AB27 = 1) & \%62; 0,52 \\
& İstihdam değişimi, 2015--2024 & \%61 \\
& İlk üç bölge (İBBS-2) & İstanbul \%27; Hatay, Kahramanmaraş, Osmaniye \%10; İzmir \%9 \\
& Çapa yatırımcılar & Kartonsan, Modern Karton (Eren Holding), Mondi Tire Kutsan, Kipaş Kağıt ve Hayat Kağıt; Kocaeli, Tekirdağ, İzmir ve Kahramanmaraş \etiketH{} \\
3 Dış ticaret & İhracat 2024; 2020--2024 değişimi & \$1.372 mn; \%64 \\
& AB payı; ilk beş AB pazarı & \%25; Yunanistan \%27; Bulgaristan \%16; Romanya \%15; İtalya \%7; İspanya \%6 \\
& İlk beş AB dışı pazar & Birleşik Krallık \%20; Irak \%10; ABD \%8; İran \%8; Mısır \%6 \\
& İthalat 2024; AB ve Çin payı & \$2.526 mn; AB \%63, Çin \%16 \\
& İlk beş ithalat kaynağı & Çin \%16; Almanya \%15; İsveç \%10; Finlandiya \%9; Rusya \%6 \\
& Dış ticaret dengesi & −\$1.154 mn \\
& İthal girdi içeriği (I-O); KD ÷ üretim & \%28; \%24 \\
& RCA: dünya 2020 → 2024; AB pazarı 2024 & 0,79 → 1,02; 1,50 \\
& Birim değer ÷ AB pazarı ortalaması & 1,00 \\
& AB pazarındaki pay (AB dışı tedarik içinde) 2020 → 2024; Türkiye'nin sırası & \%2,9 → \%6,0; 6. \\
4 AB tarafı & AB'ye ihracat 2025; rejim payları (GB / İTÜ / tarım) & €321,2 mn; \%100 / \%0 / \%0 \\
& S4'te AB tarifesi: 1/98 sürerse / askıda / tam MFN ortalaması & \%0,0 / \%0,0 / \%0,0 \\
& Tepe satırlar (S4 artışı ≥ \%10) & yok \\
& AB pazarında AB dışı rakipler (AB tercihi) & Birleşik Krallık \%16 (TCA); ABD \%15 (MFN/GSP); Çin \%13 (MFN/GSP); İsviçre \%11 (AEA/STA); Norveç \%8 (AEA/STA) \\
& Tercihli rakiplerin payı; AME (puan; ham) & \%46; 23 (0,50) \\
& Menşe kuralı (PEM) & 4801--4811: başlık değişikliği; ithal kâğıt hamuru ve atık kâğıttan kâğıt menşe kazanır \\
& Menşe riski S2 / S3; c\_RoO (\%); p\_uyum & D/D; 1,0 / 1,5; 0,95 / 0,92 \etiketH{} \\
& Değer zincirindeki rol; DİR bağımlılığı & ara; orta \\
& EUDR ve SKDM kapsamındaki AB'ye ihracat payı & \%100; \%0 \\
& PEM ortaklarına ihracat; kümülasyon maruziyeti; κ & \$249 mn; \$120 mn; 0,0 \etiketH{} \\
5 Türkiye tarafı & Uygulanan tarife 2026 (ithalat ağırlıklı): MFN; İGV dahil & \%0,0; \%4,0 \\
& S4'te AB menşeline tarife (c2) & \%3,5 \\
& Kritik girdiler: 2024 ithalatı, AB payı, ilk kaynaklar & HS4 4703, 4707: \$1.356 mn (AB \%38; Brezilya \%21, ABD \%14, Finlandiya \%13) \\
6 Kabiliyet & KS; sıra (ağırlık duyarlılığında P10--P90) & 47,1; 145 / 230 (111--176) \\
& KS bileşenleri (yüzdelik): RCA, PCI, yurt içi KD, Ar-Ge, tedarikçi, verimlilik, firma & 55, 45, 39, 25, 75, 64, 34 \\
& PCI (sınıf) & −0,10 \\
& Atlanabilir ürünler (yoğunluk üst çeyrek, PCI \textgreater{} sınıf) & 480890 (PCI +0,08, RCA 0,58); 481099 (PCI +0,03, RCA 0,58) \\
7 Stratejik önem & SS; en yüksek boyut; ağırlıklı ortalama & 23,2; kritik hammadde 30; 7,5 \etiketH{} \\
8 Makro duyarlılık & TÜFE ağırlığı (doğrudan); tüketimde ithal pay & \%0,00; \%3 \\
& ÜFE vekili (imalat üretimindeki pay) & \%0,47 \\
& EDE, 1 puan tarife: tüm menşeler / yalnız AB (bp; ρ = 1) & 0,206 / 0,130 \\
& CKE: tarife (\$ mn / bp); mali (\$ mn / mlr TL); puan & 3.386,6; 59,02; 74 \\
9 Senaryolar & Senaryo tablosu & Tablo 5.2 \\
\end{longtable}

\endgroup

\hypertarget{nace_kartlar_16-18-c9-1721-oluklu-kaux11fux131t-ve-mukavva-ile-kaux11fux131t-ve-mukavvadan-muhafazalarux131n-imalatux131}{%
\subsection{C.9 17.21 Oluklu kağıt ve mukavva ile kağıt ve mukavvadan muhafazaların imalatı}\label{nace_kartlar_16-18-c9-1721-oluklu-kaux11fux131t-ve-mukavva-ile-kaux11fux131t-ve-mukavvadan-muhafazalarux131n-imalatux131}}

\Needspace{7\baselineskip}

\textbf{Tablo C.9 --- 17.21: veri sayfası}

\nopagebreak[4]

\begingroup\gbtablo\small

\begin{longtable}[]{@{}
  >{\raggedright\arraybackslash}p{(\columnwidth - 4\tabcolsep) * \real{0.1356}}
  >{\raggedright\arraybackslash}p{(\columnwidth - 4\tabcolsep) * \real{0.3593}}
  >{\raggedright\arraybackslash}p{(\columnwidth - 4\tabcolsep) * \real{0.5050}}@{}}
\toprule\noalign{}
\begin{minipage}[b]{\linewidth}\raggedright
\textbf{Alan}
\end{minipage} & \begin{minipage}[b]{\linewidth}\raggedright
\textbf{Gösterge}
\end{minipage} & \begin{minipage}[b]{\linewidth}\raggedright
\textbf{Değer}
\end{minipage} \\
\midrule\noalign{}
\endhead
\bottomrule\noalign{}
\endlastfoot
1 Kimlik & Rev. 2.1 karşılığı; CPA & 17.21 (1:1); CPA 17.21 \\
& Başlıca HS4 (ihracat, 2024) & 4819 \%99, 4808 \%1 \\
& BEC payları: ara / sermaye / tüketim & \%100 / \%0 / \%0 \\
2 Ölçek & Üretim; katma değer (GSYH payı) & \$6.015 mn; \$1.435 mn (\%0,11) \\
& Çalışan; girişim & 40.967; 2.426 \\
& Girişim ölçeği: 0 / 1--9 / 10--49 / 50--249 / 250+ çalışan & 130 / 1.853 / 408 / 139 / 22 \\
& 250+ girişimlerin istihdam payı; kişi başı KD (AB27 = 1) & \%29; 0,44 \\
& İstihdam değişimi, 2015--2024 & \%39 \\
& İlk üç bölge (İBBS-2) & İstanbul \%30; Gaziantep, Adıyaman, Kilis \%9; İzmir \%7 \\
& Çapa yatırımcılar & Mondi Tire Kutsan, Dentaş ve entegre karton gruplarının oluklu ambalaj tesisleri; Bursa, İzmir, Kocaeli ve Gaziantep kümeleri \etiketH{} \\
3 Dış ticaret & İhracat 2024; 2020--2024 değişimi & \$834 mn; \%44 \\
& AB payı; ilk beş AB pazarı & \%39; Almanya \%19; Fransa \%18; Belçika \%8; Hollanda \%7; Romanya \%6 \\
& İlk beş AB dışı pazar & Birleşik Krallık \%21; ABD \%17; Irak \%12; Gürcistan \%5; Azerbaycan \%5 \\
& İthalat 2024; AB ve Çin payı & \$90 mn; AB \%76, Çin \%15 \\
& İlk beş ithalat kaynağı & Almanya \%29; Çin \%15; Polonya \%10; İtalya \%9; İspanya \%9 \\
& Dış ticaret dengesi & +\$744 mn \\
& İthal girdi içeriği (I-O); KD ÷ üretim & \%28; \%24 \\
& RCA: dünya 2020 → 2024; AB pazarı 2024 & 2,12 → 2,24; 3,31 \\
& Birim değer ÷ AB pazarı ortalaması & 0,81 \\
& AB pazarındaki pay (AB dışı tedarik içinde) 2020 → 2024; Türkiye'nin sırası & \%10,5 → \%13,2; 3. \\
4 AB tarafı & AB'ye ihracat 2025; rejim payları (GB / İTÜ / tarım) & €318,9 mn; \%100 / \%0 / \%0 \\
& S4'te AB tarifesi: 1/98 sürerse / askıda / tam MFN ortalaması & \%0,0 / \%0,0 / \%0,0 \\
& Tepe satırlar (S4 artışı ≥ \%10) & yok \\
& AB pazarında AB dışı rakipler (AB tercihi) & Çin \%45 (MFN/GSP); Birleşik Krallık \%13 (TCA); İsviçre \%6 (AEA/STA); Sırbistan \%5 (MFN/GSP); ABD \%3 (MFN/GSP) \\
& Tercihli rakiplerin payı; AME (puan; ham) & \%28; 31 (0,79) \\
& Menşe kuralı (PEM) & 4819: başlık değişikliği; ithal kartondan kutu menşe kazanır \\
& Menşe riski S2 / S3; c\_RoO (\%); p\_uyum & D/D; 1,0 / 1,5; 0,95 / 0,92 \etiketH{} \\
& Değer zincirindeki rol; DİR bağımlılığı & ara; düşük \\
& EUDR ve SKDM kapsamındaki AB'ye ihracat payı & \%100; \%0 \\
& PEM ortaklarına ihracat; kümülasyon maruziyeti; κ & \$94 mn; \$38 mn; 0,0 \etiketH{} \\
5 Türkiye tarafı & Uygulanan tarife 2026 (ithalat ağırlıklı): MFN; İGV dahil & \%0,0; \%24,9 \\
& S4'te AB menşeline tarife (c2) & \%24,9 \\
& Kritik girdiler: 2024 ithalatı, AB payı, ilk kaynaklar & HS4 4804, 4805: \$701 mn (AB \%66; Finlandiya \%12, Almanya \%12, Rusya \%10) \\
6 Kabiliyet & KS; sıra (ağırlık duyarlılığında P10--P90) & 47,7; 141 / 230 (86--172) \\
& KS bileşenleri (yüzdelik): RCA, PCI, yurt içi KD, Ar-Ge, tedarikçi, verimlilik, firma & 69, 28, 39, 25, 75, 44, 69 \\
& PCI (sınıf) & −0,39 \\
& Atlanabilir ürünler (yoğunluk üst çeyrek, PCI \textgreater{} sınıf) & yok \\
7 Stratejik önem & SS; en yüksek boyut; ağırlıklı ortalama & 22,8; gıda güvenliği 30; 6,0 \etiketH{} \\
8 Makro duyarlılık & TÜFE ağırlığı (doğrudan); tüketimde ithal pay & \%0,00; \%1 \\
& ÜFE vekili (imalat üretimindeki pay) & \%0,90 \\
& EDE, 1 puan tarife: tüm menşeler / yalnız AB (bp; ρ = 1) & 0,007 / 0,006 \\
& CKE: tarife (\$ mn / bp); mali (\$ mn / mlr TL); puan & 3.321,8; 1,71; 37 \\
9 Senaryolar & Senaryo tablosu & Tablo 5.3 \\
\end{longtable}

\endgroup

\hypertarget{nace_kartlar_16-18-c10-1722-kaux11fux131ttan-yapux131lan-ev-eux15fyasux131-sux131hhi-malzemeler-ve-tuvalet-malzemeleri-imalatux131}{%
\subsection{C.10 17.22 Kağıttan yapılan ev eşyası, sıhhi malzemeler ve tuvalet malzemeleri imalatı}\label{nace_kartlar_16-18-c10-1722-kaux11fux131ttan-yapux131lan-ev-eux15fyasux131-sux131hhi-malzemeler-ve-tuvalet-malzemeleri-imalatux131}}

\Needspace{7\baselineskip}

\textbf{Tablo C.10 --- 17.22: veri sayfası}

\nopagebreak[4]

\begingroup\gbtablo\small

\begin{longtable}[]{@{}
  >{\raggedright\arraybackslash}p{(\columnwidth - 4\tabcolsep) * \real{0.1354}}
  >{\raggedright\arraybackslash}p{(\columnwidth - 4\tabcolsep) * \real{0.3571}}
  >{\raggedright\arraybackslash}p{(\columnwidth - 4\tabcolsep) * \real{0.5075}}@{}}
\toprule\noalign{}
\begin{minipage}[b]{\linewidth}\raggedright
\textbf{Alan}
\end{minipage} & \begin{minipage}[b]{\linewidth}\raggedright
\textbf{Gösterge}
\end{minipage} & \begin{minipage}[b]{\linewidth}\raggedright
\textbf{Değer}
\end{minipage} \\
\midrule\noalign{}
\endhead
\bottomrule\noalign{}
\endlastfoot
1 Kimlik & Rev. 2.1 karşılığı; CPA & 17.22 (1:1); CPA 17.22 \\
& Başlıca HS4 (ihracat, 2024) & 9619 \%64, 4818 \%18, 5601 \%10, 4823 \%8 \\
& BEC payları: ara / sermaye / tüketim & \%14 / \%0 / \%86 \\
2 Ölçek & Üretim; katma değer (GSYH payı) & \$5.940 mn; \$1.281 mn (\%0,09) \\
& Çalışan; girişim & 24.641; 1.302 \\
& Girişim ölçeği: 0 / 1--9 / 10--49 / 50--249 / 250+ çalışan & 113 / 1.167 / 123 / 41 / 20 \\
& 250+ girişimlerin istihdam payı; kişi başı KD (AB27 = 1) & \%57; 0,46 \\
& İstihdam değişimi, 2015--2024 & \%63 \\
& İlk üç bölge (İBBS-2) & İstanbul \%25; Gaziantep, Adıyaman, Kilis \%20; Kocaeli, Sakarya, Düzce, Bolu, Yalova \%19 \\
& Çapa yatırımcılar & Hayat Kimya (Molfix, Papia) ve Eczacıbaşı (Selpak) başta olmak üzere entegre hijyen üreticileri; Kocaeli, Sakarya ve İstanbul \etiketH{} \\
3 Dış ticaret & İhracat 2024; 2020--2024 değişimi & \$1.284 mn; \%23 \\
& AB payı; ilk beş AB pazarı & \%18; Bulgaristan \%17; Almanya \%15; Fransa \%13; Romanya \%8; İspanya \%7 \\
& İlk beş AB dışı pazar & Irak \%27; Gürcistan \%6; İran \%5; Azerbaycan \%5; Libya \%5 \\
& İthalat 2024; AB ve Çin payı & \$186 mn; AB \%51, Çin \%10 \\
& İlk beş ithalat kaynağı & Almanya \%12; Çin \%10; Macaristan \%9; Çekya \%9; Ukrayna \%8 \\
& Dış ticaret dengesi & +\$1.098 mn \\
& İthal girdi içeriği (I-O); KD ÷ üretim & \%28; \%22 \\
& RCA: dünya 2020 → 2024; AB pazarı 2024 & 3,07 → 2,77; 2,84 \\
& Birim değer ÷ AB pazarı ortalaması & 0,85 \\
& AB pazarındaki pay (AB dışı tedarik içinde) 2020 → 2024; Türkiye'nin sırası & \%8,4 → \%11,3; 3. \\
4 AB tarafı & AB'ye ihracat 2025; rejim payları (GB / İTÜ / tarım) & €228,6 mn; \%100 / \%0 / \%0 \\
& S4'te AB tarifesi: 1/98 sürerse / askıda / tam MFN ortalaması & \%0,8 / \%0,8 / \%0,8 \\
& Tepe satırlar (S4 artışı ≥ \%10) & yok \\
& AB pazarında AB dışı rakipler (AB tercihi) & Çin \%41 (MFN/GSP); Birleşik Krallık \%18 (TCA); Ukrayna \%4 (STA/EPA); İsviçre \%4 (AEA/STA); ABD \%3 (MFN/GSP) \\
& Tercihli rakiplerin payı; AME (puan; ham) & \%33; 24 (0,51) \\
& Menşe kuralı (PEM) & 4818 ve 9619: başlık değişikliği; bazı kâğıt eşyada Fasıl 47 malzemesinden imalat \etiketD{} \\
& Menşe riski S2 / S3; c\_RoO (\%); p\_uyum & D/D; 1,0 / 1,5; 0,95 / 0,92 \etiketH{} \\
& Değer zincirindeki rol; DİR bağımlılığı & nihai; orta \\
& EUDR ve SKDM kapsamındaki AB'ye ihracat payı & \%56; \%0 \\
& PEM ortaklarına ihracat; kümülasyon maruziyeti; κ & \$179 mn; \$74 mn; 0,0 \etiketH{} \\
5 Türkiye tarafı & Uygulanan tarife 2026 (ithalat ağırlıklı): MFN; İGV dahil & \%3,9; \%7,7 \\
& S4'te AB menşeline tarife (c2) & \%6,7 \\
& Kritik girdiler: 2024 ithalatı, AB payı, ilk kaynaklar & HS4 4703, 3906: \$1.575 mn (AB \%41; Brezilya \%18, ABD \%13, Finlandiya \%11) \\
6 Kabiliyet & KS; sıra (ağırlık duyarlılığında P10--P90) & 44,0; 174 / 230 (132--196) \\
& KS bileşenleri (yüzdelik): RCA, PCI, yurt içi KD, Ar-Ge, tedarikçi, verimlilik, firma & 65, 27, 34, 25, 75, 48, 43 \\
& PCI (sınıf) & −0,45 \\
& Atlanabilir ürünler (yoğunluk üst çeyrek, PCI \textgreater{} sınıf) & yok \\
7 Stratejik önem & SS; en yüksek boyut; ağırlıklı ortalama & 29,8; sağlık 40; 6,0 \etiketH{} \\
8 Makro duyarlılık & TÜFE ağırlığı (doğrudan); tüketimde ithal pay & \%0,41; \%1 \\
& ÜFE vekili (imalat üretimindeki pay) & \%0,89 \\
& EDE, 1 puan tarife: tüm menşeler / yalnız AB (bp; ρ = 1) & 0,018 / 0,009 \\
& CKE: tarife (\$ mn / bp); mali (\$ mn / mlr TL); puan & 2.220,0; 3,05; 37 \\
9 Senaryolar & Senaryo tablosu & Tablo 5.4 \\
\end{longtable}

\endgroup

\hypertarget{nace_kartlar_16-18-c11-1723-kaux11fux131t-kux131rtasiye-uxfcruxfcnleri-imalatux131}{%
\subsection{C.11 17.23 Kağıt kırtasiye ürünleri imalatı}\label{nace_kartlar_16-18-c11-1723-kaux11fux131t-kux131rtasiye-uxfcruxfcnleri-imalatux131}}

\Needspace{7\baselineskip}

\textbf{Tablo C.11 --- 17.23: veri sayfası}

\nopagebreak[4]

\begingroup\gbtablo\small

\begin{longtable}[]{@{}
  >{\raggedright\arraybackslash}p{(\columnwidth - 4\tabcolsep) * \real{0.1363}}
  >{\raggedright\arraybackslash}p{(\columnwidth - 4\tabcolsep) * \real{0.3663}}
  >{\raggedright\arraybackslash}p{(\columnwidth - 4\tabcolsep) * \real{0.4973}}@{}}
\toprule\noalign{}
\begin{minipage}[b]{\linewidth}\raggedright
\textbf{Alan}
\end{minipage} & \begin{minipage}[b]{\linewidth}\raggedright
\textbf{Gösterge}
\end{minipage} & \begin{minipage}[b]{\linewidth}\raggedright
\textbf{Değer}
\end{minipage} \\
\midrule\noalign{}
\endhead
\bottomrule\noalign{}
\endlastfoot
1 Kimlik & Rev. 2.1 karşılığı; CPA & 17.23 (1:1); CPA 17.23 \\
& Başlıca HS4 (ihracat, 2024) & 4820 \%63, 4817 \%29, 4823 \%4, 4816 \%4 \\
& BEC payları: ara / sermaye / tüketim & \%71 / \%0 / \%29 \\
2 Ölçek & Üretim; katma değer (GSYH payı) & \$603 mn; \$163 mn (\%0,01) \\
& Çalışan; girişim & 4.538; 391 \\
& Girişim ölçeği: 0 / 1--9 / 10--49 / 50--249 / 250+ çalışan & 15 / 349 / 37 / 17 / 2 \\
& 250+ girişimlerin istihdam payı; kişi başı KD (AB27 = 1) & \%18; 0,53 \\
& İstihdam değişimi, 2015--2024 & −\%5 \\
& İlk üç bölge (İBBS-2) & İstanbul \%37; Ankara \%23; Kocaeli, Sakarya, Düzce, Bolu, Yalova \%8 \\
& Çapa yatırımcılar & --- \\
3 Dış ticaret & İhracat 2024; 2020--2024 değişimi & \$69 mn; \%23 \\
& AB payı; ilk beş AB pazarı & \%57; Almanya \%22; İtalya \%22; Romanya \%12; Fransa \%9; Yunanistan \%8 \\
& İlk beş AB dışı pazar & ABD \%15; Birleşik Krallık \%14; Azerbaycan \%7; Irak \%5; Sırbistan \%4 \\
& İthalat 2024; AB ve Çin payı & \$23 mn; AB \%32, Çin \%45 \\
& İlk beş ithalat kaynağı & Çin \%45; Endonezya \%18; Almanya \%9; İtalya \%7; İspanya \%7 \\
& Dış ticaret dengesi & +\$46 mn \\
& İthal girdi içeriği (I-O); KD ÷ üretim & \%28; \%27 \\
& RCA: dünya 2020 → 2024; AB pazarı 2024 & 0,97 → 0,86; 1,19 \\
& Birim değer ÷ AB pazarı ortalaması & 1,19 \\
& AB pazarındaki pay (AB dışı tedarik içinde) 2020 → 2024; Türkiye'nin sırası & \%4,4 → \%4,7; 4. \\
4 AB tarafı & AB'ye ihracat 2025; rejim payları (GB / İTÜ / tarım) & €30,5 mn; \%100 / \%0 / \%0 \\
& S4'te AB tarifesi: 1/98 sürerse / askıda / tam MFN ortalaması & \%0,0 / \%0,0 / \%0,0 \\
& Tepe satırlar (S4 artışı ≥ \%10) & yok \\
& AB pazarında AB dışı rakipler (AB tercihi) & Çin \%65 (MFN/GSP); Birleşik Krallık \%8 (TCA); ABD \%6 (MFN/GSP); Hindistan \%3 (MFN/GSP); Tunus \%3 (STA/EPA) \\
& Tercihli rakiplerin payı; AME (puan; ham) & \%17; 42 (1,13) \\
& Menşe kuralı (PEM) & 4817 ve 4820: başlık değişikliği \\
& Menşe riski S2 / S3; c\_RoO (\%); p\_uyum & D/D; 1,0 / 1,5; 0,95 / 0,92 \etiketH{} \\
& Değer zincirindeki rol; DİR bağımlılığı & nihai; düşük \\
& EUDR ve SKDM kapsamındaki AB'ye ihracat payı & \%100; \%0 \\
& PEM ortaklarına ihracat; kümülasyon maruziyeti; κ & \$6 mn; \$3 mn; 0,0 \etiketH{} \\
5 Türkiye tarafı & Uygulanan tarife 2026 (ithalat ağırlıklı): MFN; İGV dahil & \%0,0; \%15,6 \\
& S4'te AB menşeline tarife (c2) & \%13,8 \\
& Kritik girdiler: 2024 ithalatı, AB payı, ilk kaynaklar & --- \\
6 Kabiliyet & KS; sıra (ağırlık duyarlılığında P10--P90) & 39,8; 205 / 230 (167--214) \\
& KS bileşenleri (yüzdelik): RCA, PCI, yurt içi KD, Ar-Ge, tedarikçi, verimlilik, firma & 32, 18, 45, 25, 75, 66, 49 \\
& PCI (sınıf) & −0,67 \\
& Atlanabilir ürünler (yoğunluk üst çeyrek, PCI \textgreater{} sınıf) & 482090 (PCI −0,55, RCA 0,39) \\
7 Stratejik önem & SS; en yüksek boyut; ağırlıklı ortalama & 0,0; --- 0; 0,0 \etiketH{} \\
8 Makro duyarlılık & TÜFE ağırlığı (doğrudan); tüketimde ithal pay & \%0,00; \%1 \\
& ÜFE vekili (imalat üretimindeki pay) & \%0,09 \\
& EDE, 1 puan tarife: tüm menşeler / yalnız AB (bp; ρ = 1) & 0,002 / 0,001 \\
& CKE: tarife (\$ mn / bp); mali (\$ mn / mlr TL); puan & 2.507,8; 3,23; 38 \\
9 Senaryolar & Senaryo tablosu & Tablo C.11b (aşağıda) \\
\end{longtable}

\endgroup

\Needspace{7\baselineskip}

\textbf{Tablo C.11b --- 17.23 Kâğıt kırtasiye: senaryo etkileri (T+5'te yıllık; P50, parantez içinde P10 / P90)}

\nopagebreak[4]

\begingroup\gbtablo\small

\begin{longtable}[]{@{}lrrr@{}}
\toprule\noalign{}
\textbf{Değişken} & \textbf{S2} & \textbf{S3} & \textbf{S4} \\
\midrule\noalign{}
\endhead
\bottomrule\noalign{}
\endlastfoot
AB'ye ihracat (\$ mn) & −1 (−2 / +1) & −2 (−3 / 0) & −1 (−3 / −1) \\
AB'ye ihracat (\%) & −2,3 (−5,7 / +2,0) & −4,0 (−8,4 / −0,7) & −3,8 (−6,9 / −1,5) \\
Toplam ihracat (\$ mn) & −1 (−2 / +1) & −1 (−2 / 0) & −1 (−2 / 0) \\
Nihai mal ithalatı (\$ mn) & 0 (−1 / +1) & −1 (−2 / 0) & −1 (−3 / 0) \\
Üretim (\$ mn) & 0 (0 / 0) & −1 (−1 / 0) & 0 (0 / +1) \\
İstihdam (kişi) & −1 (−3 / +2) & −3 (−5 / 0) & 0 (−2 / +4) \\
İstihdam (\%) & 0,0 (−0,1 / 0,0) & −0,1 (−0,1 / 0,0) & 0,0 (0,0 / +0,1) \\
Cari denge (\$ mn) & 0 (0 / 0) & 0 (−1 / 0) & 0 (0 / +1) \\
TÜFE (bp) & 0,00 (0,00 / 0,00) & 0,00 (0,00 / 0,00) & +0,01 (+0,01 / +0,01) \\
\end{longtable}

\endgroup

\emph{Kaynak:} \etiketT{} f3e; toplamlar Faz 2B genel denge dağılımına kalibre. Yüzdeler 2024 düzeyine göre.

\hypertarget{nace_kartlar_16-18-c12-1724-duvar-kaux11fux131dux131-imalatux131}{%
\subsection{C.12 17.24 Duvar kağıdı imalatı}\label{nace_kartlar_16-18-c12-1724-duvar-kaux11fux131dux131-imalatux131}}

\Needspace{7\baselineskip}

\textbf{Tablo C.12 --- 17.24: veri sayfası}

\nopagebreak[4]

\begingroup\gbtablo\small

\begin{longtable}[]{@{}
  >{\raggedright\arraybackslash}p{(\columnwidth - 4\tabcolsep) * \real{0.1356}}
  >{\raggedright\arraybackslash}p{(\columnwidth - 4\tabcolsep) * \real{0.3592}}
  >{\raggedright\arraybackslash}p{(\columnwidth - 4\tabcolsep) * \real{0.5052}}@{}}
\toprule\noalign{}
\begin{minipage}[b]{\linewidth}\raggedright
\textbf{Alan}
\end{minipage} & \begin{minipage}[b]{\linewidth}\raggedright
\textbf{Gösterge}
\end{minipage} & \begin{minipage}[b]{\linewidth}\raggedright
\textbf{Değer}
\end{minipage} \\
\midrule\noalign{}
\endhead
\bottomrule\noalign{}
\endlastfoot
1 Kimlik & Rev. 2.1 karşılığı; CPA & 17.24 (1:1); CPA 17.24 \\
& Başlıca HS4 (ihracat, 2024) & 4814 \%96, 5905 \%4 \\
& BEC payları: ara / sermaye / tüketim & \%100 / \%0 / \%0 \\
2 Ölçek & Üretim; katma değer (GSYH payı) & \$62 mn; \$27 mn (\%0,00) \\
& Çalışan; girişim & 495; 37 \\
& Girişim ölçeği: 0 / 1--9 / 10--49 / 50--249 / 250+ çalışan & 2 / 29 / 7 / 2 / 0 \\
& 250+ girişimlerin istihdam payı; kişi başı KD (AB27 = 1) & \%0; 0,63 \\
& İstihdam değişimi, 2015--2024 & \%162 \\
& İlk üç bölge (İBBS-2) & Kocaeli, Sakarya, Düzce, Bolu, Yalova \%35; Adana, Mersin \%24; Tekirdağ, Edirne, Kırklareli \%17 \\
& Çapa yatırımcılar & --- \\
3 Dış ticaret & İhracat 2024; 2020--2024 değişimi & \$24 mn; \%34 \\
& AB payı; ilk beş AB pazarı & \%12; Polonya \%28; Almanya \%14; Romanya \%10; Bulgaristan \%8; Hollanda \%7 \\
& İlk beş AB dışı pazar & ABD \%20; Azerbaycan \%17; Irak \%14; Ukrayna \%5; Gürcistan \%4 \\
& İthalat 2024; AB ve Çin payı & \$8 mn; AB \%63, Çin \%22 \\
& İlk beş ithalat kaynağı & Çin \%22; Almanya \%20; İtalya \%15; Hollanda \%11; Belçika \%9 \\
& Dış ticaret dengesi & +\$16 mn \\
& İthal girdi içeriği (I-O); KD ÷ üretim & \%28; \%43 \\
& RCA: dünya 2020 → 2024; AB pazarı 2024 & 1,10 → 1,39; 0,77 \\
& Birim değer ÷ AB pazarı ortalaması & 0,78 \\
& AB pazarındaki pay (AB dışı tedarik içinde) 2020 → 2024; Türkiye'nin sırası & \%0,8 → \%3,1; 6. \\
4 AB tarafı & AB'ye ihracat 2025; rejim payları (GB / İTÜ / tarım) & €2,5 mn; \%100 / \%0 / \%0 \\
& S4'te AB tarifesi: 1/98 sürerse / askıda / tam MFN ortalaması & \%0,4 / \%0,4 / \%0,4 \\
& Tepe satırlar (S4 artışı ≥ \%10) & yok \\
& AB pazarında AB dışı rakipler (AB tercihi) & Birleşik Krallık \%39 (TCA); Çin \%24 (MFN/GSP); ABD \%10 (MFN/GSP); Ukrayna \%8 (STA/EPA); Güney Kore \%5 (STA/EPA) \\
& Tercihli rakiplerin payı; AME (puan; ham) & \%58; 18 (0,28) \\
& Menşe kuralı (PEM) & 4814: başlık değişikliği \\
& Menşe riski S2 / S3; c\_RoO (\%); p\_uyum & D/D; 1,0 / 1,5; 0,95 / 0,92 \etiketH{} \\
& Değer zincirindeki rol; DİR bağımlılığı & nihai; düşük \\
& EUDR ve SKDM kapsamındaki AB'ye ihracat payı & \%94; \%0 \\
& PEM ortaklarına ihracat; kümülasyon maruziyeti; κ & \$3 mn; \$1 mn; 0,0 \etiketH{} \\
5 Türkiye tarafı & Uygulanan tarife 2026 (ithalat ağırlıklı): MFN; İGV dahil & \%1,3; \%1,3 \\
& S4'te AB menşeline tarife (c2) & \%1,8 \\
& Kritik girdiler: 2024 ithalatı, AB payı, ilk kaynaklar & --- \\
6 Kabiliyet & KS; sıra (ağırlık duyarlılığında P10--P90) & 61,1; 21 / 230 (12--65) \\
& KS bileşenleri (yüzdelik): RCA, PCI, yurt içi KD, Ar-Ge, tedarikçi, verimlilik, firma & 60, 69, 73, 25, 75, 80, 51 \\
& PCI (sınıf) & 0,35 \\
& Atlanabilir ürünler (yoğunluk üst çeyrek, PCI \textgreater{} sınıf) & yok \\
7 Stratejik önem & SS; en yüksek boyut; ağırlıklı ortalama & 0,0; --- 0; 0,0 \etiketH{} \\
8 Makro duyarlılık & TÜFE ağırlığı (doğrudan); tüketimde ithal pay & \%0,00; \%1 \\
& ÜFE vekili (imalat üretimindeki pay) & \%0,01 \\
& EDE, 1 puan tarife: tüm menşeler / yalnız AB (bp; ρ = 1) & 0,001 / 0,000 \\
& CKE: tarife (\$ mn / bp); mali (\$ mn / mlr TL); puan & 3.526,9; 20,26; 57 \\
9 Senaryolar & Senaryo tablosu & Tablo C.12b (aşağıda) \\
\end{longtable}

\endgroup

\Needspace{7\baselineskip}

\textbf{Tablo C.12b --- 17.24 Duvar kâğıdı: senaryo etkileri (T+5'te yıllık; P50, parantez içinde P10 / P90)}

\nopagebreak[4]

\begingroup\gbtablo\small

\begin{longtable}[]{@{}lrrr@{}}
\toprule\noalign{}
\textbf{Değişken} & \textbf{S2} & \textbf{S3} & \textbf{S4} \\
\midrule\noalign{}
\endhead
\bottomrule\noalign{}
\endlastfoot
AB'ye ihracat (\$ mn) & 0 (0 / 0) & 0 (0 / 0) & 0 (0 / 0) \\
AB'ye ihracat (\%) & −5,9 (−13,6 / +2,0) & −8,0 (−15,9 / −1,3) & −7,4 (−12,9 / −3,1) \\
Toplam ihracat (\$ mn) & 0 (0 / 0) & 0 (0 / 0) & 0 (0 / 0) \\
Nihai mal ithalatı (\$ mn) & 0 (−1 / 0) & 0 (−1 / 0) & 0 (−1 / 0) \\
Üretim (\$ mn) & 0 (0 / +1) & 0 (0 / +1) & 0 (0 / 0) \\
İstihdam (kişi) & +1 (−2 / +3) & +1 (−1 / +3) & +1 (0 / +2) \\
İstihdam (\%) & +0,1 (−0,3 / +0,6) & +0,2 (−0,3 / +0,7) & +0,2 (0,0 / +0,5) \\
Cari denge (\$ mn) & 0 (0 / 0) & 0 (0 / 0) & 0 (0 / 0) \\
TÜFE (bp) & 0,00 (0,00 / 0,00) & 0,00 (0,00 / 0,00) & 0,00 (0,00 / 0,00) \\
\end{longtable}

\endgroup

\emph{Kaynak:} \etiketT{} f3e; toplamlar Faz 2B genel denge dağılımına kalibre. Yüzdeler 2024 düzeyine göre.

\hypertarget{nace_kartlar_16-18-c13-1729-kaux11fux131t-ve-mukavvadan-diux11fer-uxfcruxfcnlerin-imalatux131}{%
\subsection{C.13 17.29 Kağıt ve mukavvadan diğer ürünlerin imalatı}\label{nace_kartlar_16-18-c13-1729-kaux11fux131t-ve-mukavvadan-diux11fer-uxfcruxfcnlerin-imalatux131}}

\Needspace{7\baselineskip}

\textbf{Tablo C.13 --- 17.29: veri sayfası}

\nopagebreak[4]

\begingroup\gbtablo\small

\begin{longtable}[]{@{}
  >{\raggedright\arraybackslash}p{(\columnwidth - 4\tabcolsep) * \real{0.1360}}
  >{\raggedright\arraybackslash}p{(\columnwidth - 4\tabcolsep) * \real{0.3628}}
  >{\raggedright\arraybackslash}p{(\columnwidth - 4\tabcolsep) * \real{0.5012}}@{}}
\toprule\noalign{}
\begin{minipage}[b]{\linewidth}\raggedright
\textbf{Alan}
\end{minipage} & \begin{minipage}[b]{\linewidth}\raggedright
\textbf{Gösterge}
\end{minipage} & \begin{minipage}[b]{\linewidth}\raggedright
\textbf{Değer}
\end{minipage} \\
\midrule\noalign{}
\endhead
\bottomrule\noalign{}
\endlastfoot
1 Kimlik & Rev. 2.1 karşılığı; CPA & 17.25 (1:1); CPA 17.29 \\
& Başlıca HS4 (ihracat, 2024) & 4821 \%45, 4823 \%40, 4822 \%9, 4813 \%5 \\
& BEC payları: ara / sermaye / tüketim & \%98 / \%0 / \%2 \\
2 Ölçek & Üretim; katma değer (GSYH payı) & \$1.644 mn; \$520 mn (\%0,04) \\
& Çalışan; girişim & 14.587; 925 \\
& Girişim ölçeği: 0 / 1--9 / 10--49 / 50--249 / 250+ çalışan & 66 / 696 / 166 / 61 / 5 \\
& 250+ girişimlerin istihdam payı; kişi başı KD (AB27 = 1) & \%16; 0,48 \\
& İstihdam değişimi, 2015--2024 & \%44 \\
& İlk üç bölge (İBBS-2) & İstanbul \%37; Bursa, Eskişehir, Bilecik \%12; İzmir \%11 \\
& Çapa yatırımcılar & Etiket, filtre ve gıda ambalajı kâğıdı üreticileri (İstanbul, Kocaeli, İzmir) \etiketH{} \\
3 Dış ticaret & İhracat 2024; 2020--2024 değişimi & \$187 mn; \%59 \\
& AB payı; ilk beş AB pazarı & \%35; İtalya \%16; Romanya \%16; Almanya \%15; İspanya \%9; Hollanda \%7 \\
& İlk beş AB dışı pazar & Mısır \%15; ABD \%9; Irak \%7; Birleşik Krallık \%6; Gürcistan \%5 \\
& İthalat 2024; AB ve Çin payı & \$147 mn; AB \%77, Çin \%12 \\
& İlk beş ithalat kaynağı & Almanya \%22; İtalya \%20; Çin \%12; Avusturya \%11; Çekya \%6 \\
& Dış ticaret dengesi & +\$40 mn \\
& İthal girdi içeriği (I-O); KD ÷ üretim & \%28; \%32 \\
& RCA: dünya 2020 → 2024; AB pazarı 2024 & 0,84 → 0,95; 1,45 \\
& Birim değer ÷ AB pazarı ortalaması & 0,84 \\
& AB pazarındaki pay (AB dışı tedarik içinde) 2020 → 2024; Türkiye'nin sırası & \%4,4 → \%5,8; 5. \\
4 AB tarafı & AB'ye ihracat 2025; rejim payları (GB / İTÜ / tarım) & €65,4 mn; \%100 / \%0 / \%0 \\
& S4'te AB tarifesi: 1/98 sürerse / askıda / tam MFN ortalaması & \%0,0 / \%0,0 / \%0,0 \\
& Tepe satırlar (S4 artışı ≥ \%10) & yok \\
& AB pazarında AB dışı rakipler (AB tercihi) & Çin \%40 (MFN/GSP); Birleşik Krallık \%17 (TCA); ABD \%9 (MFN/GSP); Hindistan \%8 (MFN/GSP); İsviçre \%4 (AEA/STA) \\
& Tercihli rakiplerin payı; AME (puan; ham) & \%32; 28 (0,70) \\
& Menşe kuralı (PEM) & 4821--4823: başlık değişikliği \\
& Menşe riski S2 / S3; c\_RoO (\%); p\_uyum & D/D; 1,0 / 1,5; 0,95 / 0,92 \etiketH{} \\
& Değer zincirindeki rol; DİR bağımlılığı & ara; düşük \\
& EUDR ve SKDM kapsamındaki AB'ye ihracat payı & \%100; \%0 \\
& PEM ortaklarına ihracat; kümülasyon maruziyeti; κ & \$42 mn; \$17 mn; 0,0 \etiketH{} \\
5 Türkiye tarafı & Uygulanan tarife 2026 (ithalat ağırlıklı): MFN; İGV dahil & \%0,0; \%16,7 \\
& S4'te AB menşeline tarife (c2) & \%16,0 \\
& Kritik girdiler: 2024 ithalatı, AB payı, ilk kaynaklar & --- \\
6 Kabiliyet & KS; sıra (ağırlık duyarlılığında P10--P90) & 51,7; 94 / 230 (68--131) \\
& KS bileşenleri (yüzdelik): RCA, PCI, yurt içi KD, Ar-Ge, tedarikçi, verimlilik, firma & 48, 52, 57, 25, 75, 55, 66 \\
& PCI (sınıf) & 0,05 \\
& Atlanabilir ürünler (yoğunluk üst çeyrek, PCI \textgreater{} sınıf) & yok \\
7 Stratejik önem & SS; en yüksek boyut; ağırlıklı ortalama & 7,9; sağlık 10; 3,0 \etiketH{} \\
8 Makro duyarlılık & TÜFE ağırlığı (doğrudan); tüketimde ithal pay & \%0,00; \%2 \\
& ÜFE vekili (imalat üretimindeki pay) & \%0,25 \\
& EDE, 1 puan tarife: tüm menşeler / yalnız AB (bp; ρ = 1) & 0,012 / 0,009 \\
& CKE: tarife (\$ mn / bp); mali (\$ mn / mlr TL); puan & 3.333,3; 9,36; 47 \\
9 Senaryolar & Senaryo tablosu & Tablo C.13b (aşağıda) \\
\end{longtable}

\endgroup

\Needspace{7\baselineskip}

\textbf{Tablo C.13b --- 17.29 Diğer kâğıt ürünleri: senaryo etkileri (T+5'te yıllık; P50, parantez içinde P10 / P90)}

\nopagebreak[4]

\begingroup\gbtablo\small

\begin{longtable}[]{@{}lrrr@{}}
\toprule\noalign{}
\textbf{Değişken} & \textbf{S2} & \textbf{S3} & \textbf{S4} \\
\midrule\noalign{}
\endhead
\bottomrule\noalign{}
\endlastfoot
AB'ye ihracat (\$ mn) & −2 (−5 / +1) & −3 (−7 / −1) & −3 (−6 / −1) \\
AB'ye ihracat (\%) & −3,1 (−7,7 / +2,0) & −5,3 (−11,3 / −0,9) & −4,9 (−9,2 / −2,1) \\
Toplam ihracat (\$ mn) & −1 (−4 / +1) & −2 (−5 / 0) & −3 (−5 / −1) \\
Nihai mal ithalatı (\$ mn) & −9 (−31 / +8) & −13 (−33 / +5) & −26 (−47 / −6) \\
Üretim (\$ mn) & +8 (−7 / +27) & +10 (−5 / +28) & +23 (+4 / +43) \\
İstihdam (kişi) & +47 (−43 / +169) & +64 (−32 / +178) & +145 (+26 / +271) \\
İstihdam (\%) & +0,3 (−0,3 / +1,2) & +0,4 (−0,2 / +1,2) & +1,0 (+0,2 / +1,9) \\
Cari denge (\$ mn) & +5 (−5 / +20) & +7 (−4 / +20) & +17 (+3 / +31) \\
TÜFE (bp) & +0,02 (+0,01 / +0,02) & +0,03 (+0,03 / +0,04) & +0,13 (+0,11 / +0,14) \\
\end{longtable}

\endgroup

\emph{Kaynak:} \etiketT{} f3e; toplamlar Faz 2B genel denge dağılımına kalibre. Yüzdeler 2024 düzeyine göre.

\hypertarget{nace_kartlar_16-18-c14-1811-gazetelerin-basux131mux131}{%
\subsection{C.14 18.11 Gazetelerin basımı}\label{nace_kartlar_16-18-c14-1811-gazetelerin-basux131mux131}}

\Needspace{7\baselineskip}

\textbf{Tablo C.14 --- 18.11: veri sayfası}

\nopagebreak[4]

\begingroup\gbtablo\small

\begin{longtable}[]{@{}
  >{\raggedright\arraybackslash}p{(\columnwidth - 4\tabcolsep) * \real{0.1427}}
  >{\raggedright\arraybackslash}p{(\columnwidth - 4\tabcolsep) * \real{0.4288}}
  >{\raggedright\arraybackslash}p{(\columnwidth - 4\tabcolsep) * \real{0.4285}}@{}}
\toprule\noalign{}
\begin{minipage}[b]{\linewidth}\raggedright
\textbf{Alan}
\end{minipage} & \begin{minipage}[b]{\linewidth}\raggedright
\textbf{Gösterge}
\end{minipage} & \begin{minipage}[b]{\linewidth}\raggedright
\textbf{Değer}
\end{minipage} \\
\midrule\noalign{}
\endhead
\bottomrule\noalign{}
\endlastfoot
1 Kimlik & Rev. 2.1 karşılığı; CPA & 18.11 (1:1); CPA 18.11 \\
& Başlıca HS4 (ihracat, 2024) & --- \\
& BEC payları: ara / sermaye / tüketim & --- \\
2 Ölçek & Üretim; katma değer (GSYH payı) & \$79 mn; \$25 mn (\%0,00) \\
& Çalışan; girişim & 1.965; 407 \\
& Girişim ölçeği: 0 / 1--9 / 10--49 / 50--249 / 250+ çalışan & 19 / 370 / 39 / 4 / 0 \\
& 250+ girişimlerin istihdam payı; kişi başı KD (AB27 = 1) & \%0; 0,22 \\
& İstihdam değişimi, 2015--2024 & −\%47 \\
& İlk üç bölge (İBBS-2) & İstanbul \%29; Adana, Mersin \%10; Trabzon, Ordu, Giresun, Rize, Artvin, Gümüşhane \%5 \\
& Çapa yatırımcılar & --- \\
3 Dış ticaret & İhracat ve ithalat 2024 & Mal ticareti yok (hizmet sınıfı ya da ürünleri başka sınıflara eşleniyor) \\
& İthal girdi içeriği (I-O); KD ÷ üretim & \%26; \%32 \\
4 AB tarafı & AB'ye ihracat ve S4 tarifesi & AB'ye mal ihracatı yok \\
& Menşe kuralı (PEM) & Hizmet (gazete basımı); mal ticareti yok \\
& Menşe riski S2 / S3; c\_RoO (\%); p\_uyum & D/D; 1,0 / 1,5; 0,95 / 0,92 \etiketH{} \\
& Değer zincirindeki rol; DİR bağımlılığı & ara; düşük \\
& EUDR ve SKDM kapsamındaki AB'ye ihracat payı & ---; --- \\
& PEM ortaklarına ihracat; kümülasyon maruziyeti; κ & \$0 mn; \$0 mn; 0,0 \etiketH{} \\
5 Türkiye tarafı & Uygulanan tarife 2026 (ithalat ağırlıklı): MFN; İGV dahil & ---; --- \\
& S4'te AB menşeline tarife (c2) & \%0,0 \\
& Kritik girdiler: 2024 ithalatı, AB payı, ilk kaynaklar & --- \\
6 Kabiliyet & KS; sıra (ağırlık duyarlılığında P10--P90) & 41,7; 198 / 230 (126--216) \\
& KS bileşenleri (yüzdelik): RCA, PCI, yurt içi KD, Ar-Ge, tedarikçi, verimlilik, firma & ---, ---, 58, 25, 60, 5, 62 \\
& PCI (sınıf) & --- \\
& Atlanabilir ürünler (yoğunluk üst çeyrek, PCI \textgreater{} sınıf) & yok \\
7 Stratejik önem & SS; en yüksek boyut; ağırlıklı ortalama & 0,0; --- 0; 0,0 \etiketH{} \\
8 Makro duyarlılık & TÜFE ağırlığı (doğrudan); tüketimde ithal pay & \%0,09; \%0 \\
& ÜFE vekili (imalat üretimindeki pay) & \%0,01 \\
& EDE, 1 puan tarife: tüm menşeler / yalnız AB (bp; ρ = 1) & 0,000 / 0,000 \\
& CKE: tarife (\$ mn / bp); mali (\$ mn / mlr TL); puan & ---; 0,00; 5 \\
9 Senaryolar & Senaryo tablosu & Tablo C.14b (aşağıda) \\
\end{longtable}

\endgroup

\Needspace{7\baselineskip}

\textbf{Tablo C.14b --- 18.11 Gazete basımı: senaryo etkileri (T+5'te yıllık; P50, parantez içinde P10 / P90)}

\nopagebreak[4]

\begingroup\gbtablo\small

\begin{longtable}[]{@{}lrrr@{}}
\toprule\noalign{}
\textbf{Değişken} & \textbf{S2} & \textbf{S3} & \textbf{S4} \\
\midrule\noalign{}
\endhead
\bottomrule\noalign{}
\endlastfoot
AB'ye ihracat (\$ mn) & 0 (0 / 0) & 0 (0 / 0) & 0 (0 / 0) \\
AB'ye ihracat (\%) & --- & --- & --- \\
Toplam ihracat (\$ mn) & 0 (0 / 0) & 0 (0 / 0) & 0 (0 / 0) \\
Nihai mal ithalatı (\$ mn) & 0 (0 / 0) & 0 (0 / 0) & 0 (0 / 0) \\
Üretim (\$ mn) & 0 (0 / 0) & 0 (0 / 0) & 0 (0 / 0) \\
İstihdam (kişi) & 0 (0 / 0) & 0 (0 / 0) & 0 (0 / 0) \\
İstihdam (\%) & 0,0 (0,0 / 0,0) & 0,0 (0,0 / 0,0) & 0,0 (0,0 / 0,0) \\
Cari denge (\$ mn) & 0 (0 / 0) & 0 (0 / 0) & 0 (0 / 0) \\
TÜFE (bp) & 0,00 (0,00 / 0,00) & 0,00 (0,00 / 0,00) & 0,00 (0,00 / 0,00) \\
\end{longtable}

\endgroup

\emph{Kaynak:} \etiketT{} f3e; toplamlar Faz 2B genel denge dağılımına kalibre. Yüzdeler 2024 düzeyine göre.

\hypertarget{nace_kartlar_16-18-c15-1812-diux11fer-matbaacux131lux131k}{%
\subsection{C.15 18.12 Diğer matbaacılık}\label{nace_kartlar_16-18-c15-1812-diux11fer-matbaacux131lux131k}}

\Needspace{7\baselineskip}

\textbf{Tablo C.15 --- 18.12: veri sayfası}

\nopagebreak[4]

\begingroup\gbtablo\small

\begin{longtable}[]{@{}
  >{\raggedright\arraybackslash}p{(\columnwidth - 4\tabcolsep) * \real{0.1331}}
  >{\raggedright\arraybackslash}p{(\columnwidth - 4\tabcolsep) * \real{0.3338}}
  >{\raggedright\arraybackslash}p{(\columnwidth - 4\tabcolsep) * \real{0.5332}}@{}}
\toprule\noalign{}
\begin{minipage}[b]{\linewidth}\raggedright
\textbf{Alan}
\end{minipage} & \begin{minipage}[b]{\linewidth}\raggedright
\textbf{Gösterge}
\end{minipage} & \begin{minipage}[b]{\linewidth}\raggedright
\textbf{Değer}
\end{minipage} \\
\midrule\noalign{}
\endhead
\bottomrule\noalign{}
\endlastfoot
1 Kimlik & Rev. 2.1 karşılığı; CPA & 13.30, 18.12 (1:2); CPA 18.12 \\
& Başlıca HS4 (ihracat, 2024) & --- \\
& BEC payları: ara / sermaye / tüketim & --- \\
2 Ölçek & Üretim; katma değer (GSYH payı) & \$3.944 mn; \$1.069 mn (\%0,08) \\
& Çalışan; girişim & 47.966; 10.768 \\
& Girişim ölçeği: 0 / 1--9 / 10--49 / 50--249 / 250+ çalışan & 258 / 10.399 / 590 / 87 / 14 \\
& 250+ girişimlerin istihdam payı; kişi başı KD (AB27 = 1) & \%7; 0,42 \\
& İstihdam değişimi, 2015--2024 & \%1 \\
& İlk üç bölge (İBBS-2) & İstanbul \%42; Ankara \%9; İzmir \%9 \\
& Çapa yatırımcılar & Ofset, kitap ve ambalaj matbaaları (İstanbul'da Bağcılar ve Topkapı matbaacılar siteleri, Ankara, İzmir); güvenlikli basımda TCMB Banknot Matbaası \etiketH{} \\
3 Dış ticaret & İhracat ve ithalat 2024 & Mal ticareti yok (hizmet sınıfı ya da ürünleri başka sınıflara eşleniyor) \\
& İthal girdi içeriği (I-O); KD ÷ üretim & \%26; \%27 \\
4 AB tarafı & AB'ye ihracat ve S4 tarifesi & AB'ye mal ihracatı yok \\
& Menşe kuralı (PEM) & Hizmet (basım); basılı ürün ihracatı (Fasıl 49) yayıncılık sınıflarında (58.xx), ambalaj baskısı 17.21 ve 17.29'da \\
& Menşe riski S2 / S3; c\_RoO (\%); p\_uyum & D/D; 1,0 / 1,5; 0,95 / 0,92 \etiketH{} \\
& Değer zincirindeki rol; DİR bağımlılığı & ara; düşük \\
& EUDR ve SKDM kapsamındaki AB'ye ihracat payı & ---; --- \\
& PEM ortaklarına ihracat; kümülasyon maruziyeti; κ & \$0 mn; \$0 mn; 0,0 \etiketH{} \\
5 Türkiye tarafı & Uygulanan tarife 2026 (ithalat ağırlıklı): MFN; İGV dahil & ---; --- \\
& S4'te AB menşeline tarife (c2) & \%0,0 \\
& Kritik girdiler: 2024 ithalatı, AB payı, ilk kaynaklar & HS4 4802, 4810, 3215: \$1.373 mn (AB \%57; Çin \%20, İsveç \%13, Almanya \%11) \\
6 Kabiliyet & KS; sıra (ağırlık duyarlılığında P10--P90) & 48,4; 131 / 230 (70--179) \\
& KS bileşenleri (yüzdelik): RCA, PCI, yurt içi KD, Ar-Ge, tedarikçi, verimlilik, firma & ---, ---, 47, 25, 60, 40, 83 \\
& PCI (sınıf) & --- \\
& Atlanabilir ürünler (yoğunluk üst çeyrek, PCI \textgreater{} sınıf) & yok \\
7 Stratejik önem & SS; en yüksek boyut; ağırlıklı ortalama & 15,2; savunma 20; 4,0 \etiketH{} \\
8 Makro duyarlılık & TÜFE ağırlığı (doğrudan); tüketimde ithal pay & \%0,00; \%0 \\
& ÜFE vekili (imalat üretimindeki pay) & \%0,59 \\
& EDE, 1 puan tarife: tüm menşeler / yalnız AB (bp; ρ = 1) & 0,000 / 0,000 \\
& CKE: tarife (\$ mn / bp); mali (\$ mn / mlr TL); puan & ---; 0,00; 5 \\
9 Senaryolar & Senaryo tablosu & Tablo C.15b (aşağıda) \\
\end{longtable}

\endgroup

\Needspace{7\baselineskip}

\textbf{Tablo C.15b --- 18.12 Matbaacılık: senaryo etkileri (T+5'te yıllık; P50, parantez içinde P10 / P90)}

\nopagebreak[4]

\begingroup\gbtablo\small

\begin{longtable}[]{@{}lrrr@{}}
\toprule\noalign{}
\textbf{Değişken} & \textbf{S2} & \textbf{S3} & \textbf{S4} \\
\midrule\noalign{}
\endhead
\bottomrule\noalign{}
\endlastfoot
AB'ye ihracat (\$ mn) & 0 (0 / 0) & 0 (0 / 0) & 0 (0 / 0) \\
AB'ye ihracat (\%) & --- & --- & --- \\
Toplam ihracat (\$ mn) & 0 (0 / 0) & 0 (0 / 0) & 0 (0 / 0) \\
Nihai mal ithalatı (\$ mn) & 0 (0 / 0) & 0 (0 / 0) & 0 (0 / 0) \\
Üretim (\$ mn) & 0 (0 / 0) & 0 (0 / 0) & 0 (0 / 0) \\
İstihdam (kişi) & 0 (0 / 0) & 0 (0 / 0) & 0 (0 / 0) \\
İstihdam (\%) & 0,0 (0,0 / 0,0) & 0,0 (0,0 / 0,0) & 0,0 (0,0 / 0,0) \\
Cari denge (\$ mn) & 0 (0 / 0) & 0 (0 / 0) & 0 (0 / 0) \\
TÜFE (bp) & 0,00 (0,00 / 0,00) & 0,00 (0,00 / 0,00) & 0,00 (0,00 / 0,00) \\
\end{longtable}

\endgroup

\emph{Kaynak:} \etiketT{} f3e; toplamlar Faz 2B genel denge dağılımına kalibre. Yüzdeler 2024 düzeyine göre.

\hypertarget{nace_kartlar_16-18-c16-1813-basux131m-ve-yayux131m-uxf6ncesi-hizmetler}{%
\subsection{C.16 18.13 Basım ve yayım öncesi hizmetler}\label{nace_kartlar_16-18-c16-1813-basux131m-ve-yayux131m-uxf6ncesi-hizmetler}}

\Needspace{7\baselineskip}

\textbf{Tablo C.16 --- 18.13: veri sayfası}

\nopagebreak[4]

\begingroup\gbtablo\small

\begin{longtable}[]{@{}
  >{\raggedright\arraybackslash}p{(\columnwidth - 4\tabcolsep) * \real{0.1361}}
  >{\raggedright\arraybackslash}p{(\columnwidth - 4\tabcolsep) * \real{0.3635}}
  >{\raggedright\arraybackslash}p{(\columnwidth - 4\tabcolsep) * \real{0.5004}}@{}}
\toprule\noalign{}
\begin{minipage}[b]{\linewidth}\raggedright
\textbf{Alan}
\end{minipage} & \begin{minipage}[b]{\linewidth}\raggedright
\textbf{Gösterge}
\end{minipage} & \begin{minipage}[b]{\linewidth}\raggedright
\textbf{Değer}
\end{minipage} \\
\midrule\noalign{}
\endhead
\bottomrule\noalign{}
\endlastfoot
1 Kimlik & Rev. 2.1 karşılığı; CPA & 18.13 (1:1); CPA 18.13 \\
& Başlıca HS4 (ihracat, 2024) & 8442 \%100 \\
& BEC payları: ara / sermaye / tüketim & \%100 / \%0 / \%0 \\
2 Ölçek & Üretim; katma değer (GSYH payı) & \$179 mn; \$60 mn (\%0,00) \\
& Çalışan; girişim & 3.170; 1.118 \\
& Girişim ölçeği: 0 / 1--9 / 10--49 / 50--249 / 250+ çalışan & 41 / 1.159 / 29 / 7 / 0 \\
& 250+ girişimlerin istihdam payı; kişi başı KD (AB27 = 1) & \%0; 0,47 \\
& İstihdam değişimi, 2015--2024 & \%22 \\
& İlk üç bölge (İBBS-2) & İstanbul \%48; İzmir \%12; Gaziantep, Adıyaman, Kilis \%5 \\
& Çapa yatırımcılar & --- \\
3 Dış ticaret & İhracat 2024; 2020--2024 değişimi & \$16 mn; \%26 \\
& AB payı; ilk beş AB pazarı & \%63; Almanya \%43; Hollanda \%17; Polonya \%7; Slovakya \%7; Avusturya \%7 \\
& İlk beş AB dışı pazar & Rusya \%22; Kazakistan \%9; Birleşik Krallık \%9; Endonezya \%8; İran \%5 \\
& İthalat 2024; AB ve Çin payı & \$17 mn; AB \%52, Çin \%29 \\
& İlk beş ithalat kaynağı & Almanya \%36; Çin \%29; İtalya \%11; İsviçre \%6; Birleşik Krallık \%4 \\
& Dış ticaret dengesi & −\$2 mn \\
& İthal girdi içeriği (I-O); KD ÷ üretim & \%26; \%34 \\
& RCA: dünya 2020 → 2024; AB pazarı 2024 & 1,76 → 1,90; 3,21 \\
& Birim değer ÷ AB pazarı ortalaması & 0,43 \\
& AB pazarındaki pay (AB dışı tedarik içinde) 2020 → 2024; Türkiye'nin sırası & \%9,5 → \%12,8; 4. \\
4 AB tarafı & AB'ye ihracat 2025; rejim payları (GB / İTÜ / tarım) & €8,9 mn; \%100 / \%0 / \%0 \\
& S4'te AB tarifesi: 1/98 sürerse / askıda / tam MFN ortalaması & \%0,0 / \%0,0 / \%0,0 \\
& Tepe satırlar (S4 artışı ≥ \%10) & yok \\
& AB pazarında AB dışı rakipler (AB tercihi) & İsviçre \%40 (AEA/STA); Birleşik Krallık \%14 (TCA); Çin \%13 (MFN/GSP); ABD \%8 (MFN/GSP); Japonya \%5 (STA/EPA) \\
& Tercihli rakiplerin payı; AME (puan; ham) & \%61; 47 (1,25) \\
& Menşe kuralı (PEM) & 8442 50: başlık değişikliği (baskı kalıbı ve klişe) \\
& Menşe riski S2 / S3; c\_RoO (\%); p\_uyum & D/D; 1,0 / 1,5; 0,95 / 0,92 \etiketH{} \\
& Değer zincirindeki rol; DİR bağımlılığı & ara; düşük \\
& EUDR ve SKDM kapsamındaki AB'ye ihracat payı & \%0; \%0 \\
& PEM ortaklarına ihracat; kümülasyon maruziyeti; κ & \$1 mn; \$0 mn; 0,0 \etiketH{} \\
5 Türkiye tarafı & Uygulanan tarife 2026 (ithalat ağırlıklı): MFN; İGV dahil & \%1,7; \%1,7 \\
& S4'te AB menşeline tarife (c2) & \%1,7 \\
& Kritik girdiler: 2024 ithalatı, AB payı, ilk kaynaklar & --- \\
6 Kabiliyet & KS; sıra (ağırlık duyarlılığında P10--P90) & 57,5; 44 / 230 (29--83) \\
& KS bileşenleri (yüzdelik): RCA, PCI, yurt içi KD, Ar-Ge, tedarikçi, verimlilik, firma & 66, 72, 63, 25, 60, 51, 58 \\
& PCI (sınıf) & 0,40 \\
& Atlanabilir ürünler (yoğunluk üst çeyrek, PCI \textgreater{} sınıf) & yok \\
7 Stratejik önem & SS; en yüksek boyut; ağırlıklı ortalama & 0,0; --- 0; 0,0 \etiketH{} \\
8 Makro duyarlılık & TÜFE ağırlığı (doğrudan); tüketimde ithal pay & \%0,00; \%0 \\
& ÜFE vekili (imalat üretimindeki pay) & \%0,03 \\
& EDE, 1 puan tarife: tüm menşeler / yalnız AB (bp; ρ = 1) & 0,001 / 0,001 \\
& CKE: tarife (\$ mn / bp); mali (\$ mn / mlr TL); puan & 1.902,3; 4,97; 37 \\
9 Senaryolar & Senaryo tablosu & Tablo C.16b (aşağıda) \\
\end{longtable}

\endgroup

\Needspace{7\baselineskip}

\textbf{Tablo C.16b --- 18.13 Baskı öncesi hizmetler: senaryo etkileri (T+5'te yıllık; P50, parantez içinde P10 / P90)}

\nopagebreak[4]

\begingroup\gbtablo\small

\begin{longtable}[]{@{}lrrr@{}}
\toprule\noalign{}
\textbf{Değişken} & \textbf{S2} & \textbf{S3} & \textbf{S4} \\
\midrule\noalign{}
\endhead
\bottomrule\noalign{}
\endlastfoot
AB'ye ihracat (\$ mn) & 0 (0 / 0) & 0 (−1 / 0) & 0 (0 / 0) \\
AB'ye ihracat (\%) & −1,5 (−3,7 / +2,0) & −2,7 (−5,6 / −0,5) & −2,5 (−4,7 / −1,1) \\
Toplam ihracat (\$ mn) & 0 (0 / 0) & 0 (0 / 0) & 0 (0 / 0) \\
Nihai mal ithalatı (\$ mn) & 0 (−1 / +1) & 0 (−1 / 0) & 0 (−1 / 0) \\
Üretim (\$ mn) & 0 (0 / 0) & 0 (0 / 0) & 0 (0 / 0) \\
İstihdam (kişi) & +1 (−6 / +5) & +1 (−5 / +5) & +1 (0 / +3) \\
İstihdam (\%) & 0,0 (−0,2 / +0,2) & 0,0 (−0,2 / +0,2) & 0,0 (0,0 / +0,1) \\
Cari denge (\$ mn) & 0 (0 / 0) & 0 (0 / 0) & 0 (0 / 0) \\
TÜFE (bp) & 0,00 (0,00 / 0,00) & 0,00 (0,00 / 0,00) & 0,00 (0,00 / 0,00) \\
\end{longtable}

\endgroup

\emph{Kaynak:} \etiketT{} f3e; toplamlar Faz 2B genel denge dağılımına kalibre. Yüzdeler 2024 düzeyine göre.

\hypertarget{nace_kartlar_16-18-c17-1814-ciltuxe7ilik-ve-ilgili-hizmetler}{%
\subsection{C.17 18.14 Ciltçilik ve ilgili hizmetler}\label{nace_kartlar_16-18-c17-1814-ciltuxe7ilik-ve-ilgili-hizmetler}}

\Needspace{7\baselineskip}

\textbf{Tablo C.17 --- 18.14: veri sayfası}

\nopagebreak[4]

\begingroup\gbtablo\small

\begin{longtable}[]{@{}
  >{\raggedright\arraybackslash}p{(\columnwidth - 4\tabcolsep) * \real{0.1458}}
  >{\raggedright\arraybackslash}p{(\columnwidth - 4\tabcolsep) * \real{0.4602}}
  >{\raggedright\arraybackslash}p{(\columnwidth - 4\tabcolsep) * \real{0.3939}}@{}}
\toprule\noalign{}
\begin{minipage}[b]{\linewidth}\raggedright
\textbf{Alan}
\end{minipage} & \begin{minipage}[b]{\linewidth}\raggedright
\textbf{Gösterge}
\end{minipage} & \begin{minipage}[b]{\linewidth}\raggedright
\textbf{Değer}
\end{minipage} \\
\midrule\noalign{}
\endhead
\bottomrule\noalign{}
\endlastfoot
1 Kimlik & Rev. 2.1 karşılığı; CPA & 18.14 (1:1); CPA 18.14 \\
& Başlıca HS4 (ihracat, 2024) & --- \\
& BEC payları: ara / sermaye / tüketim & --- \\
2 Ölçek & Üretim; katma değer (GSYH payı) & \$60 mn; \$25 mn (\%0,00) \\
& Çalışan; girişim & 1.422; 322 \\
& Girişim ölçeği: 0 / 1--9 / 10--49 / 50--249 / 250+ çalışan & 11 / 313 / 22 / 1 / 0 \\
& 250+ girişimlerin istihdam payı; kişi başı KD (AB27 = 1) & \%0; 0,46 \\
& İstihdam değişimi, 2015--2024 & −\%19 \\
& İlk üç bölge (İBBS-2) & İstanbul \%29; Ankara \%16; Antalya, Isparta, Burdur \%6 \\
& Çapa yatırımcılar & --- \\
3 Dış ticaret & İhracat ve ithalat 2024 & Mal ticareti yok (hizmet sınıfı ya da ürünleri başka sınıflara eşleniyor) \\
& İthal girdi içeriği (I-O); KD ÷ üretim & \%26; \%42 \\
4 AB tarafı & AB'ye ihracat ve S4 tarifesi & AB'ye mal ihracatı yok \\
& Menşe kuralı (PEM) & Hizmet (ciltleme); mal ticareti yok \\
& Menşe riski S2 / S3; c\_RoO (\%); p\_uyum & D/D; 1,0 / 1,5; 0,95 / 0,92 \etiketH{} \\
& Değer zincirindeki rol; DİR bağımlılığı & ara; düşük \\
& EUDR ve SKDM kapsamındaki AB'ye ihracat payı & ---; --- \\
& PEM ortaklarına ihracat; kümülasyon maruziyeti; κ & \$0 mn; \$0 mn; 0,0 \etiketH{} \\
5 Türkiye tarafı & Uygulanan tarife 2026 (ithalat ağırlıklı): MFN; İGV dahil & ---; --- \\
& S4'te AB menşeline tarife (c2) & \%0,0 \\
& Kritik girdiler: 2024 ithalatı, AB payı, ilk kaynaklar & --- \\
6 Kabiliyet & KS; sıra (ağırlık duyarlılığında P10--P90) & 51,7; 92 / 230 (49--153) \\
& KS bileşenleri (yüzdelik): RCA, PCI, yurt içi KD, Ar-Ge, tedarikçi, verimlilik, firma & ---, ---, 74, 25, 60, 48, 55 \\
& PCI (sınıf) & --- \\
& Atlanabilir ürünler (yoğunluk üst çeyrek, PCI \textgreater{} sınıf) & yok \\
7 Stratejik önem & SS; en yüksek boyut; ağırlıklı ortalama & 0,0; --- 0; 0,0 \etiketH{} \\
8 Makro duyarlılık & TÜFE ağırlığı (doğrudan); tüketimde ithal pay & \%0,00; \%0 \\
& ÜFE vekili (imalat üretimindeki pay) & \%0,01 \\
& EDE, 1 puan tarife: tüm menşeler / yalnız AB (bp; ρ = 1) & 0,000 / 0,000 \\
& CKE: tarife (\$ mn / bp); mali (\$ mn / mlr TL); puan & ---; 0,00; 5 \\
9 Senaryolar & Senaryo tablosu & Tablo C.17b (aşağıda) \\
\end{longtable}

\endgroup

\Needspace{7\baselineskip}

\textbf{Tablo C.17b --- 18.14 Ciltçilik: senaryo etkileri (T+5'te yıllık; P50, parantez içinde P10 / P90)}

\nopagebreak[4]

\begingroup\gbtablo\small

\begin{longtable}[]{@{}lrrr@{}}
\toprule\noalign{}
\textbf{Değişken} & \textbf{S2} & \textbf{S3} & \textbf{S4} \\
\midrule\noalign{}
\endhead
\bottomrule\noalign{}
\endlastfoot
AB'ye ihracat (\$ mn) & 0 (0 / 0) & 0 (0 / 0) & 0 (0 / 0) \\
AB'ye ihracat (\%) & --- & --- & --- \\
Toplam ihracat (\$ mn) & 0 (0 / 0) & 0 (0 / 0) & 0 (0 / 0) \\
Nihai mal ithalatı (\$ mn) & 0 (0 / 0) & 0 (0 / 0) & 0 (0 / 0) \\
Üretim (\$ mn) & 0 (0 / 0) & 0 (0 / 0) & 0 (0 / 0) \\
İstihdam (kişi) & 0 (0 / 0) & 0 (0 / 0) & 0 (0 / 0) \\
İstihdam (\%) & 0,0 (0,0 / 0,0) & 0,0 (0,0 / 0,0) & 0,0 (0,0 / 0,0) \\
Cari denge (\$ mn) & 0 (0 / 0) & 0 (0 / 0) & 0 (0 / 0) \\
TÜFE (bp) & 0,00 (0,00 / 0,00) & 0,00 (0,00 / 0,00) & 0,00 (0,00 / 0,00) \\
\end{longtable}

\endgroup

\emph{Kaynak:} \etiketT{} f3e; toplamlar Faz 2B genel denge dağılımına kalibre. Yüzdeler 2024 düzeyine göre.

\hypertarget{nace_kartlar_16-18-c18-1820-kayux131tlux131-medyanux131n-uxe7oux11faltux131lmasux131}{%
\subsection{C.18 18.20 Kayıtlı medyanın çoğaltılması}\label{nace_kartlar_16-18-c18-1820-kayux131tlux131-medyanux131n-uxe7oux11faltux131lmasux131}}

\Needspace{7\baselineskip}

\textbf{Tablo C.18 --- 18.20: veri sayfası}

\nopagebreak[4]

\begingroup\gbtablo\small

\begin{longtable}[]{@{}
  >{\raggedright\arraybackslash}p{(\columnwidth - 4\tabcolsep) * \real{0.1458}}
  >{\raggedright\arraybackslash}p{(\columnwidth - 4\tabcolsep) * \real{0.4602}}
  >{\raggedright\arraybackslash}p{(\columnwidth - 4\tabcolsep) * \real{0.3939}}@{}}
\toprule\noalign{}
\begin{minipage}[b]{\linewidth}\raggedright
\textbf{Alan}
\end{minipage} & \begin{minipage}[b]{\linewidth}\raggedright
\textbf{Gösterge}
\end{minipage} & \begin{minipage}[b]{\linewidth}\raggedright
\textbf{Değer}
\end{minipage} \\
\midrule\noalign{}
\endhead
\bottomrule\noalign{}
\endlastfoot
1 Kimlik & Rev. 2.1 karşılığı; CPA & 18.20 (1:1); CPA 18.20 \\
& Başlıca HS4 (ihracat, 2024) & --- \\
& BEC payları: ara / sermaye / tüketim & --- \\
2 Ölçek & Üretim; katma değer (GSYH payı) & \$12 mn; \$5 mn (\%0,00) \\
& Çalışan; girişim & 199; 51 \\
& Girişim ölçeği: 0 / 1--9 / 10--49 / 50--249 / 250+ çalışan & --- / 65 / --- / --- / 0 \\
& 250+ girişimlerin istihdam payı; kişi başı KD (AB27 = 1) & \%0; --- \\
& İstihdam değişimi, 2015--2024 & \%169 \\
& İlk üç bölge (İBBS-2) & İstanbul \%47; Ankara \%11; Kayseri, Sivas, Yozgat \%4 \\
& Çapa yatırımcılar & --- \\
3 Dış ticaret & İhracat ve ithalat 2024 & Mal ticareti yok (hizmet sınıfı ya da ürünleri başka sınıflara eşleniyor) \\
& İthal girdi içeriği (I-O); KD ÷ üretim & \%26; \%42 \\
4 AB tarafı & AB'ye ihracat ve S4 tarifesi & AB'ye mal ihracatı yok \\
& Menşe kuralı (PEM) & 8523: başlık değişikliği \\
& Menşe riski S2 / S3; c\_RoO (\%); p\_uyum & D/D; 1,0 / 1,5; 0,95 / 0,92 \etiketH{} \\
& Değer zincirindeki rol; DİR bağımlılığı & nihai; düşük \\
& EUDR ve SKDM kapsamındaki AB'ye ihracat payı & ---; --- \\
& PEM ortaklarına ihracat; kümülasyon maruziyeti; κ & \$0 mn; \$0 mn; 0,0 \etiketH{} \\
5 Türkiye tarafı & Uygulanan tarife 2026 (ithalat ağırlıklı): MFN; İGV dahil & ---; --- \\
& S4'te AB menşeline tarife (c2) & \%0,0 \\
& Kritik girdiler: 2024 ithalatı, AB payı, ilk kaynaklar & --- \\
6 Kabiliyet & KS; sıra (ağırlık duyarlılığında P10--P90) & 50,9; 107 / 230 (46--171) \\
& KS bileşenleri (yüzdelik): RCA, PCI, yurt içi KD, Ar-Ge, tedarikçi, verimlilik, firma & ---, ---, 73, 25, 60, ---, 47 \\
& PCI (sınıf) & --- \\
& Atlanabilir ürünler (yoğunluk üst çeyrek, PCI \textgreater{} sınıf) & yok \\
7 Stratejik önem & SS; en yüksek boyut; ağırlıklı ortalama & 0,0; --- 0; 0,0 \etiketH{} \\
8 Makro duyarlılık & TÜFE ağırlığı (doğrudan); tüketimde ithal pay & \%0,00; \%0 \\
& ÜFE vekili (imalat üretimindeki pay) & \%0,00 \\
& EDE, 1 puan tarife: tüm menşeler / yalnız AB (bp; ρ = 1) & 0,000 / 0,000 \\
& CKE: tarife (\$ mn / bp); mali (\$ mn / mlr TL); puan & ---; 0,00; 5 \\
9 Senaryolar & Senaryo tablosu & Tablo C.18b (aşağıda) \\
\end{longtable}

\endgroup

\Needspace{7\baselineskip}

\textbf{Tablo C.18b --- 18.20 Kayıtlı ortam çoğaltma: senaryo etkileri (T+5'te yıllık; P50, parantez içinde P10 / P90)}

\nopagebreak[4]

\begingroup\gbtablo\small

\begin{longtable}[]{@{}lrrr@{}}
\toprule\noalign{}
\textbf{Değişken} & \textbf{S2} & \textbf{S3} & \textbf{S4} \\
\midrule\noalign{}
\endhead
\bottomrule\noalign{}
\endlastfoot
AB'ye ihracat (\$ mn) & 0 (0 / 0) & 0 (0 / 0) & 0 (0 / 0) \\
AB'ye ihracat (\%) & --- & --- & --- \\
Toplam ihracat (\$ mn) & 0 (0 / 0) & 0 (0 / 0) & 0 (0 / 0) \\
Nihai mal ithalatı (\$ mn) & 0 (0 / 0) & 0 (0 / 0) & 0 (0 / 0) \\
Üretim (\$ mn) & 0 (0 / 0) & 0 (0 / 0) & 0 (0 / 0) \\
İstihdam (kişi) & 0 (0 / 0) & 0 (0 / 0) & 0 (0 / 0) \\
İstihdam (\%) & 0,0 (0,0 / 0,0) & 0,0 (0,0 / 0,0) & 0,0 (0,0 / 0,0) \\
Cari denge (\$ mn) & 0 (0 / 0) & 0 (0 / 0) & 0 (0 / 0) \\
TÜFE (bp) & 0,00 (0,00 / 0,00) & 0,00 (0,00 / 0,00) & 0,00 (0,00 / 0,00) \\
\end{longtable}

\endgroup

\emph{Kaynak:} \etiketT{} f3e; toplamlar Faz 2B genel denge dağılımına kalibre. Yüzdeler 2024 düzeyine göre.

\hypertarget{nace_kartlar_16-18-ek-d--uxf6nceki-paruxe7alara-ve-fazlara-duxfczeltmeler}{%
\section{Ek D --- Önceki parçalara ve fazlara düzeltmeler}\label{nace_kartlar_16-18-ek-d--uxf6nceki-paruxe7alara-ve-fazlara-duxfczeltmeler}}

\textbf{D-7. Kâğıt hamurunda iç ikame ve imalat toplamları.} Kâğıt hamurunda (17.11) iç ikame tavanının sıfırlanması, sınıfın S2'deki 129 ve S4'teki 98 kişilik anlamsız istihdam artışını kaldırdı. Bu artış önceki parçaların imalat geneli toplamlarına giriyordu. İmalatın S4 istihdam etkisi medyanda −46,2 bin kişi oldu (ikinci parçada −46,1 bin).

\textbf{D-8. İlk iki parçanın S2 ve S3 sonuçları.} S2'de AB'ye ihracattaki toplam kayıp genel denge sonucuna bağlı ve sınıflara kısmi denge kayıp paylarıyla dağıtılıyor. Bu parçadaki iki değişiklik bu dağılımı kaydırdı. Birincisi ve büyüğü, S2'de menşe maliyetinin yalnız tercihin kurtardığı tarife kadar sayılması (Ek B, V-49). İkincisi, ağaç, kâğıt ve basımın menşe parametrelerinin varsayılan değerlerden (c\_RoO \%1,5, p\_uyum 0,90) sınıfa özgü değerlere (\%1, 0,95) çekilmesi. İkisi de tarifesi düşük sınıfların S2 kaybını küçültüp tarifeli sınıflarınkini büyüttü. İkinci parçanın S2 istihdam etkisi −15.453 kişi oldu (önceki sürümde −14.419; P10 −32.928), ilk parçanınki −318 (önceki sürümde −286). İmalatın S2 istihdam etkisi −24,9 bin kişi (önceki sürümde −24,3 bin). S3'ün genel denge toplamı, S2 ile S4'ün kısmi denge toplamları arasındaki konumuyla belirlendiği için S3 sonuçları da değişti. S4 sonuçları yalnız CE kanalının ahşap sınıflarına eklenmesiyle çok az değişti. İkinci parçanın temel bulgusu değişmiyor, hatta güçleniyor: giyimde iyi bir anlaşmanın değeri büyük ve S2'de bile kayıp kalıcı. İlk iki parçanın belgeleri bu sayılarla yeniden üretildi; metinlerindeki bütün sayılar veriden geldiği için ayrıca düzeltme gerekmedi.

\textbf{D-9. Faz 2C'nin son tarihleri.} Faz 2C'nin listelediği iki son tarih bu belgenin tarihine çok yaklaştı. DS629'da (Çin--Türkiye) panel raporunun kabulü ya da temyizi için son gün 24 Kasım 2026, yani bu belgenin tarihinden iki gün sonra; bir temyiz başvurusuna ilişkin kayıt bulunamadı \etiketD{}. 2027--2029 dönemi için GATT XXVIII:5 çekincesinin son günü 31 Aralık 2026; Türkiye'nin bir bildirimine ilişkin kayıt da bulunamadı \etiketD{}. İki hamlenin (H-D1, H-D2) önceliği Faz 2C'de belirlendi ve değişmedi; durumları \kod{DURUM\_\allowbreak{}OZETI.\allowbreak{}md}'de açık konular arasında izleniyor.

\hypertarget{nace_kartlar_16-18-ek-e--varsayux131mlar-ve-veri-boux15fluklarux131}{%
\section{Ek E --- Varsayımlar ve veri boşlukları}\label{nace_kartlar_16-18-ek-e--varsayux131mlar-ve-veri-boux15fluklarux131}}

Bu parçada varsayım kaydına altı yeni madde eklendi, bir madde güncellendi, beş veri boşluğu açıldı ve biri güncellendi. VB-32, yani DTÖ bağlı oranları, bu parçada da bağlayıcı olmadı: c2 uygulanan MFN oranlarını değiştirmiyor ve önerilen tek tarife değişikliği (E-17.12-01) bir indirim.

\Needspace{7\baselineskip}

\textbf{Tablo E.1 --- Varsayım kaydına eklenen maddeler}

\nopagebreak[4]

\begingroup\gbtablo\small

\begin{longtable}[]{@{}
  >{\raggedright\arraybackslash}p{(\columnwidth - 2\tabcolsep) * \real{0.1365}}
  >{\raggedright\arraybackslash}p{(\columnwidth - 2\tabcolsep) * \real{0.8635}}@{}}
\toprule\noalign{}
\begin{minipage}[b]{\linewidth}\raggedright
\textbf{Kimlik}
\end{minipage} & \begin{minipage}[b]{\linewidth}\raggedright
\textbf{Varsayım}
\end{minipage} \\
\midrule\noalign{}
\endhead
\bottomrule\noalign{}
\endlastfoot
V-35 (güncellendi) & SS boyut puanları ve ağırlıkları \etiketH{} (veri/f3\_ss\_boyutlar.csv; 10--18 gözden geçirildi, diğerleri ilk tahmin) \\
V-44 & iç ikame tavanı (veri/f3\_sinif\_varsayim.csv, ikame\_tavani): yerli üretimi AB ithalatının yerini alamayan sınıflarda iç ikame payı sınırlanır; kâğıt hamurunda (17.11) 0 \etiketH{} \\
V-45 & CE ve onaylanmış kuruluş bağımlılığı (ce\_payi): ahşap levhada 0,15, ahşap döşemede 0,5, doğramada 0,3 (AB Yapı Ürünleri Tüzüğü; AB'ye ihracatta üçüncü taraf deneyi gereken pay) \etiketH{} \\
V-46 & ağaç, kâğıt ve basımda menşe: Fasıl 44, 47--49 ve 4818'de başlık değişikliği; S2'de risk düşük (c\_RoO \%1, p\_uyum 0,95), S3'te \%1,5 ve 0,92 \etiketH{}; kurallar \etiketV{} \\
V-47 & EUDR payı KN koduna göre (f3d); tamamen geri dönüştürülmüş ürün ve yalnız taşıma amaçlı ambalaj istisnası payı hesaba katılmadı (üst sınır, VB-41) \\
V-48 & E-17.12-01'in (ambalaj kâğıdında c2 düzeltmesi) etkisi: TÜFE −0,2 bp, mali maliyet AB dışı menşelerin İGV kaybı (≈ \$16 mn ≈ 0,7 mlr TL/yıl); AB menşeli ürünün c2 geliri bugün olmadığı için maliyet sayılmadı \etiketT{} \\
V-49 & S2 ve S3'te menşe maliyeti yalnız tercihin kurtardığı tarife kadar: τ\_etkin = (1 − p\_uyum)·τ + q·p\_uyum·min(c\_RoO, τ); ithalatta min(c\_AB, τ\_TR,c2). MFN'si sıfır üründe ihracatçı tercih istemez. Sınıf ortalaması düzeyinde yaklaşık \etiketT{} \\
\end{longtable}

\endgroup

\Needspace{7\baselineskip}

\textbf{Tablo E.2 --- Açık veri boşlukları}

\nopagebreak[4]

\begingroup\gbtablo\small

\begin{longtable}[]{@{}
  >{\raggedright\arraybackslash}p{(\columnwidth - 2\tabcolsep) * \real{0.1331}}
  >{\raggedright\arraybackslash}p{(\columnwidth - 2\tabcolsep) * \real{0.8669}}@{}}
\toprule\noalign{}
\begin{minipage}[b]{\linewidth}\raggedright
\textbf{Kimlik}
\end{minipage} & \begin{minipage}[b]{\linewidth}\raggedright
\textbf{Veri boşluğu}
\end{minipage} \\
\midrule\noalign{}
\endhead
\bottomrule\noalign{}
\endlastfoot
VB-40 & levha, kontrplak ve kâğıtta ithal odun ve kâğıt hamurunun firma düzeyinde menşe dağılımı (Rusya ve Belarus payı; EUDR ve anti-damping riski) \\
VB-41 & tamamen geri dönüştürülmüş lifle üretilen kâğıt ve ambalajın ihracattaki payı (EUDR kapsamı dışında; sınıfların EUDR payı bu yüzden üst sınır) \\
VB-42 & basılı ürün ve baskı hizmeti ihracatının sınıf eşlemesi (Fasıl 49 yayıncılık sınıflarına, 58.xx, gidiyor; 18.11, 18.12, 18.14 ve 18.20'nin dış ticareti ölçülemiyor) \\
VB-43 & c2 rejiminde yurt içi kabiliyet bayrağı HS6 düzeyinde hesaplanıyor; sıvı ambalaj kartonu (4811 59) ve kraftliner (4804 11) gibi kalitelerde yanlış sınıflama, 8--12 haneli üretim verisi gerekir \\
VB-44 & kâğıt hamuru (17.11) YSHİ verisi gizli; sınıfın çalışanı ve üretimi bilinmiyor \\
VB-32 (sürüyor) & DTÖ bağlı oranları, satır düzeyi. 13--18'de bağlayıcı değil: c2 uygulanan MFN oranlarını değiştirmiyor, yalnız AB tercihini kaldırıyor; 16--18'deki tek tarife önerisi (E-17.12-01) bir indirim \\
\end{longtable}

\endgroup

\begin{gbilerleme}

İLERLEME: Faz 3 parça 3 / 16 tamamlandı (72 / 230 C sınıfı kartı) --- SONRAKİ ADIM: Faz 3 parça 4 (19--21: kok ve rafine petrol, kimyasallar, eczacılık); DEVAM bekleniyor.

\end{gbilerleme}
